\documentclass[a4paper,11pt]{article}
% %%%%%%%%%%%%%%%%%%%%%%%%%%%%%%%%%%%%%%%%%%%%%%%%%%%%%%%%%%%%%%%%%%%%%
\usepackage[a4paper, total={6.5in, 9.8in}]{geometry}
% %%%%%%%%%%%%%%%%%%%%%%%%%%%%%%%%%%%%%%%%%%%%%%%%%%%%%%%%%%%%%%%%%%%%%
\usepackage{amssymb}
\usepackage{amsfonts}
\usepackage{amsmath}
\usepackage{makeidx}
\usepackage{fontspec}
\usepackage{xunicode}
\usepackage{xltxtra}
\usepackage{listings}
\usepackage{xcolor}
\usepackage{graphicx}
\usepackage{morefloats}
\usepackage[yyyymmdd]{datetime}
\usepackage{appendix}
\usepackage{tikz}
\usepackage{enumitem}
\usepackage[most]{tcolorbox}
\usepackage{hyperref}
\hypersetup{
  pdftitle={Matrix Module Documentation (matrix\textunderscore mod\textunderscore matrix)},
  pdfauthor={Konstantinos N. Anagnostopoulos},
  colorlinks=true,
  linkcolor=blue!70!black,
  filecolor=magenta,
  citecolor=blue!70!black,
  anchorcolor=blue!70!black,
  urlcolor=blue!70!black,
  bookmarksopen=true,
  bookmarksdepth=3
}
\urlstyle{same}
\setlength{\marginparwidth}{2cm}
\emergencystretch=1.5em

\setmainfont[Mapping=tex-text]{GFS Didot}

% Load user's LaTeX macros and listings configuration
%\newcommand{\part}{\partial}
%\renewcommand{\Huge}{\Large} % προβλημα με ελληνικες γραμματοσειρες
%\renewcommand{\huge}{\Large} % προβλημα με ελληνικες γραμματοσειρες
%%\renewcommand\contentsname{ΠΕΡΙΕΧΟΜΕΝΑ}
%%\renewcommand\listfigurename{ΚΑΤΑΛΟΓΟΣ ΤΩΝ ΕΙΚΟΝΩΝ}
%%\renewcommand\listtablename{ΚΑΤΑΛΟΓΟΣ ΠΙΝΑΚΩΝ}
%%%\newcommand\refname{ΑΝΑΦΟΡΕΣ}
%%\renewcommand\refname{ΑΝΑΦΟΡΕΣ}
%%\renewcommand\indexname{Index}
%%\renewcommand\figurename{Σχήμα}
%%\renewcommand\tablename{Πίνακας}
%%\renewcommand\partname{Μέρος}
%%\renewcommand\chaptername{ΚΕΦΑΛΑΙΟ}
%%\renewcommand\appendixname{ΠΑΡΑΡΤΗΜΑ}
%%%\newcommand\abstractname{ΠΕΡΙΛΗΨΗ}
%%\renewcommand\abstractname{ΠΕΡΙΛΗΨΗ}
%%\renewcommand\bibname{ΒΙΒΛΙΟΓΡΑΦΙΑ}
%%\def\today{\number\day \space\ifcase\month\or
%%  Ιανουαρίου \or Φεβρουαρίου\or Μαρτίου\or Απριλίου\or Μαΐου\or Ιουνίου\or
%%  Ιουλίου\or Αυγούστου\or Σεπτεμβρίου\or Οκτωβρίου\or Νοεμβρίου\or
%%  Δεκεμβρίου\fi 
%%   \space\number\year}

% %%%%%%%%%%%%%%%%%%%%%%%%%%%%%%%%%%%%%%%%%%%%%%%%%%%%%%%%%%
% %%%%%%%%%%%%%%%%%%%%%%%%%%%%%%%%%%%%%%%%%%%%%%%%%%%%%%%%%%
%                     math definitions
% %%%%%%%%%%%%%%%%%%%%%%%%%%%%%%%%%%%%%%%%%%%%%%%%%%%%%%%%%%
\newcommand  {\e}       [1]{{\rm e}^{#1}}
\newcommand  {\vev}     [1]{\ensuremath{\langle #1 \rangle}}
\newcommand  {\veveq}   [1]{\ensuremath{{\langle #1 \rangle}_{\mbox{\scriptsize eq}}}}
\newcommand  {\norm}    [1]{\ensuremath{\left\lVert #1 \right\rVert}}  % double vertical line norm like ||x||
\newcommand  {\rf}      [1]{(\ref{#1})}
\newcommand  {\eq}      [1]{Eq.(\protect\ref{#1})}
\newcommand  {\eqs}     [2]{Eqs.(\protect\ref{#1})--(\protect\ref{#2})}
\newcommand  {\fig}     [1]{Figure \protect\ref{#1}}
\newcommand  {\figs}    [2]{Figures \protect\ref{#1} -- \protect\ref{#2}}
\newcommand  {\tbl}     [1]{Table  \protect\ref{#1}}
\newcommand  {\code}    [1]{\colorbox{listinggray}{{\tt #1}}}                           % \definecolor{listinggray}{gray}{0.9}  -- defined below
\newcommand  {\lcode}   [1]{\colorbox{listinggray}{\lstinline[language=bash       ]{#1}}}   
\newcommand  {\lcodeb}  [1]{\colorbox{listinggray}{\lstinline[language=bash       ]{#1}}}   
\newcommand  {\lcodef}  [1]{\colorbox{listinggray}{\lstinline[language=Fortran    ]{#1}}}    
\newcommand  {\lcodec}  [1]{\colorbox{listinggray}{\lstinline[language=C          ]{#1}}}   
\newcommand  {\lcodea}  [1]{\colorbox{listinggray}{\lstinline[language=Awk        ]{#1}}}   
\newcommand  {\lcodeg}  [1]{\colorbox{listinggray}{\lstinline[language=Gnuplot    ]{#1}}}   
\newcommand  {\lcodem}  [1]{\colorbox{listinggray}{\lstinline[language=Mathematica]{#1}}}   
\newcommand  {\cbox}    [1]{\colorbox{listinggray}{{\tt #1}}}                           
\renewcommand{\implies}    {\ensuremath{\Rightarrow}}
\newcommand  {\nn}         {\nonumber\\}
\newcommand  {\ket}     [1]{\ensuremath{\,| #1 \rangle}}
\newcommand  {\bra}     [1]{\ensuremath{\langle #1 |\,}}
\newcommand  {\braket}  [2]{\ensuremath{\langle #1 | #2 \rangle}}
\newcommand  {\calO}       {\ensuremath{\mathcal{O}}}
\newcommand  {\calA}       {\ensuremath{\mathcal{A}}}
\newcommand  {\calM}       {\ensuremath{\mathcal{M}}}
\newcommand  {\calD}       {\ensuremath{\mathcal{D}}}
\newcommand  {\calN}       {\ensuremath{\mathcal{N}}}
\newcommand  {\calQ}       {\ensuremath{\mathcal{Q}}}
\newcommand  {\smbox}   [1]{\ensuremath{\mbox{\eng{\scriptsize{#1}}}}}
\newcommand  {\embox}   [1]{\ensuremath{\mbox{\eng{{#1}}}}}
\newcommand  {\padd}       {\ensuremath{P_{\smbox{add}}}}
\newcommand  {\abs}     [1]{\ensuremath{\left| #1 \right|}}
\newcommand  {\pd}      [2]{\ensuremath{\frac{\partial   #1}{\partial #2}}}
\newcommand  {\pdd}     [2]{\ensuremath{\frac{\partial^2 #1}{\partial {#2}^2}}}
\newcommand  {\tr}         {\ensuremath{{\rm tr\,}}}
\newcommand  {\Tr}         {\ensuremath{{\rm Tr\,}}}
\newcommand  {\Reals}      {\ensuremath{\mathbb{R}}}
\newcommand  {\Complexes}  {\ensuremath{\mathbb{C}}}
\newcommand  {\Naturals}   {\ensuremath{\mathbb{N}}}
\newcommand  {\Rationals}  {\ensuremath{\mathbb{Q}}}
\newcommand  {\Integers}   {\ensuremath{\mathbb{Z}}}
\newcommand  {\UnitMatrix} {\ensuremath{\mathbb{I}}}
\newcommand  {\el}         {\ensuremath{\in}}        % Mathematica notation
\renewcommand{\Re}         {\ensuremath{\mbox{Re}}}
\renewcommand{\Im}         {\ensuremath{\mbox{Im}}}
%\newcommand  {\Seff}       {\ensuremath{S_{\text{\scriptsize{eff}}}}}
%\newcommand  {\tSeff}      {\ensuremath{\tilde{S}_{\text{\scriptsize{eff}}}}}
\newcommand  {\Seff}       {\ensuremath{S_{\text{eff}}}}
\newcommand  {\tSeff}      {\ensuremath{\tilde{S}_{\text{eff}}}}
\newcommand  {\dtau}       {\ensuremath{d\tau}}
\newcommand  {\taut}       {\ensuremath{\tau_{\mbox{\scriptsize aut}}}}
\newcommand  {\tint}       {\ensuremath{\tau_{\mbox{\scriptsize int}}}}
\newcommand  {\MyPhi}      {\ensuremath{P}}                             % notation of pdf of stochastic variables
\newcommand  {\MyPhieq}    {\ensuremath{\MyPhi_{\mbox{\scriptsize eq}}}}
\newcommand  {\mscript} [2]{\ensuremath{{#1}_{\mbox{\scriptsize #2}}}}
\newcommand  {\mscr}    [2]{\ensuremath{{#1}_{\mbox{\scriptsize #2}}}}
\newcommand  {\Pf}         {\ensuremath{\text{Pf}\,}}
\newcommand  {\PfM}         {\ensuremath{\text{Pf}\,\mathcal{M}}}
\newcommand  {\detM}        {\ensuremath{\det\mathcal{M}}}
\newcommand  {\detD}        {\ensuremath{\det\mathcal{D}}}
\newcommand  {\mf}          {\ensuremath{m_{\text{f}}}}
\newcommand  {\ncg}         {\ensuremath{N_{\text{cg}}}}
\newcommand  {\NS}          {\ensuremath{\mathscr{D}}}  % symbol for spinor dimensionality 2^{D/2-1}. Choices: N_S,N_s, \mathscr{D}, \mathfrak{D}, d_s
% %%%%%%%%%%%%%%%%%%%%%%%%%%%%%%%%%%%%%%%%%%%%%%%%%%%%%%%%%%
% tofo notes: \usepackage[colorinlistoftodos,shadow]{todonotes}, \listoftodos, \todo[optios]{text}, \missingfigure[figwidth=0.5\textwidth]{text}, \todo[inline,size=\Large,color=red!40]{#1} (option disable: \usepackage[disable,...]{todonotes})
\usepackage[colorinlistoftodos,shadow]{todonotes}
\newcommand {\XXX}     [1]{\todo[size=\tiny,color=green!40]{#1}}
\newcommand {\XXXS}    [1]{\todo[size=\small,color=green!40]{#1}}
\newcommand {\XXXs}    [1]{\todo[size=\tiny ,color=green!40]{#1}}
\newcommand {\XXXi}    [1]{\todo[inline,color=green!40]{#1}}
\newcommand {\XXXX}    [1]{\todo[inline,color=green!40]{#1}}
\newcommand {\XXXtodo} [1]{\todo{#1}}
\newcommand {\XXXred}  [1]{\todo[size=\small,color=red!40]{#1}}
\newcommand {\XXXXXX}  [1]{\todo[inline,size=\Large,color=red!40]{#1}}
\newcommand {\XXXref}  [1]{\todo[size=\small,linecolor=green!70!white,backgroundcolor=blue!20!white,bordercolor=red]{Add Ref.: #1}}
\newcommand {\XXXfig}  [1]{\todo[inline,color=green!40]{#1}} %insert in captions
\newcommand {\XXXtab}  [1]{\todo[inline,color=green!40]{#1}}
\newcommand {\XXXaddfig}[1]{\missingfigure[figwidth=0.5\textwidth]{#1}}
% %%%%%%%%%%%%%%%%%%%%%%%%%%%%%%%%%%%%%%%%%%%%%%%%%%%%%%%%%%
% index alias:
\newcommand {\ii}      [1]{\index{#1}}
% %%%%%%%%%%%%%%%%%%%%%%%%%%%%%%%%%%%%%%%%%%%%%%%%%%%%%%%%%%
% Other packages:
\usepackage{booktabs} % For tables in chapter 06_Parallel
%\usepackage{afterpage} % For big floats placement. Use \afterpage{ \begin{figure}[!htb] ...  \end{figure} }
\usepackage{float}    % For big floats. use  \begin{figure}[H] to actually place it here!
\usepackage{mathrsfs} % \mathscr{D}
% %%%%%%%%%%%%%%%%%%%%%%%%%%%%%%%%%%%%%%%%%%%%%%%%%%%%%%%%%%
%                     hyphenation
% %%%%%%%%%%%%%%%%%%%%%%%%%%%%%%%%%%%%%%%%%%%%%%%%%%%%%%%%%%
%\hyphenation{every-where}
\hyphenation{Lan-ge-vin}
% %%%%%%%%%%%%%%%%%%%%%%%%%%%%%%%%%%%%%%%%%%%%%%%%%%%%%%%%%%
%                     include subsubsections
% %%%%%%%%%%%%%%%%%%%%%%%%%%%%%%%%%%%%%%%%%%%%%%%%%%%%%%%%%%
\setcounter{secnumdepth}{3}
\setcounter{tocdepth}{3}


% %%%%%%%%%%%%%%%%%%%%%%%%%%%%%%%%%%%%%%%%%%%%%%%%%%%%%%%%%
%   Code listing: 
% %%%%%%%%%%%%%%%%%%%%%%%%%%%%%%%%%%%%%%%%%%%%%%%%%%%%%%%%%
%\usepackage{listings}
% /usr/share/doc/texlive-doc/latex/listings/listings.pdf
% Languages: C, Fortran, [77]Fortran,  Gnuplot, Awk, csh, TeX,
% LaTeX,make, Perl, Lisp,HTML....
%\usepackage{color}
%%\lstset{language=Gnuplot,alsolanguage=csh,basicstyle=\scriptsize,deletekeywords={gawk,emacs},morekeywords={print,printf}}
%%\begin{lstlisting}[caption={Hello there},label=lbl]
%%\begin{lstlisting}[title={My title}]
%% ..................
%%\end{lstlisting}
\definecolor{listinggray}{gray}{0.9}
\definecolor{lbcolor}{rgb}{0.9,0.9,0.9}
\lstset{
        backgroundcolor=\color{lbcolor},
        tabsize=4,
        rulecolor=,
        language=Fortran,
        basicstyle=\footnotesize,%\scriptsize,
        upquote=true,
        aboveskip={1.5\baselineskip},
        columns=fixed,
        showstringspaces=false,
        extendedchars=true,
        inputencoding=utf8x,
        breaklines=true,
        prebreak = \raisebox{0ex}[0ex][0ex]{\ensuremath{\hookleftarrow}},
        frame=single,
        showtabs=false,
        showspaces=false,
        showstringspaces=false,
        identifierstyle=\ttfamily,
        keywordstyle=\color[rgb]{0,0,1},
        commentstyle=\color[rgb]{0.133,0.545,0.133},
        stringstyle=\color[rgb]{0.627,0.126,0.941},
}
% from /usr/share/texlive/texmf-dist/tex/latex/listings/lstlang1.sty
\lstdefinelanguage[95]{Fortran}[77]{Fortran}%
  {deletekeywords=SAVE,%
   morekeywords={subroutine,ACTION,ADVANCE,ALLOCATE,ALLOCATABLE,ASSIGNMENT,CASE,%
      CONTAINS,CYCLE,DEALLOCATE,DEFAULT,DELIM,EXIT,INCLUDE,IN,NONE,IN,%
      OUT,INTENT,INTERFACE,IOLENGTH,KIND,LEN,MODULE,NAME,NAMELIST,NMT,%
      NULLIFY,ONLY,OPERATOR,OPTIONAL,OUT,PAD,POINTER,POSITION,PRIVATE,%
      PUBLIC,READWRITE,RECURSIVE,RESULT,SELECT,SEQUENCE,SIZE,STAT,%
      TARGET,USE,WHERE,WHILE,BLOCKDATA,DOUBLEPRECISION,%
      ENDBLOCKDATA,ENDFILE,ENDFUNCTION,ENDINTERFACE,%
      ENDMODULE,ENDPROGRAM,ENDSELECT,ENDSUBROUTINE,ENDTYPE,ENDWHERE,%
      INOUT,SELECTCASE,FORALL,RESHAPE,SUM,PRODUCT,ELSEWHERE,%
      sin,cos,tan,atan2,exp,cexp,dexp,cdexp,log,abs,conjg,dconjg,%
      dble,dimag,cdabs,imag,tiny,huge,epsilon,range,trim,atan,pure,%
      ifdef,ifndef,define,endif,elif,%
      MATMUL,TRANSPOSE,UBOUND,LBOUND,SAVE,%
      MPI_LOGICAL,MPI_INTEGER,MPI_INTEGER8,MPI_REAL,MPI_REAL8,MPI_COMPLEX,MPI_COMPEX16,MPI_CHARACTER,MPI_STATUS_IGNORE,%
      MPI_SUM,MPI_PROD,MPI_MAX,MPI_MIN,%
      concurrent,CONCURRENT,random_number,%
      co_broadcast,co_sum,co_max,com_min,co_reduce,sync,all,this_image,image_index,num_images,error,lcobound,ucobound,codimension,%
      mpi_init,mpi_comm_size,mpi_comm_rank,mpi_finalize,mpi_gather,mpi_allgather,mpi_send,mpi_recv,%
      mpi_reduce,mpi_allreduce,MPI_COMM_WORLD,mpi_bcast,mpi_wtime,mpi_get_count,mpi_scatter,mpi_barrier,mpi_sendrecv,MPI_TAG,MPI_SOURCE,%
      mpi_comm_group,mpi_group_incl,mpi_comm_create,mpi_comm_split,mpi_comm_null,mpi_comm_free,mpi_sendrecv_replace,%
      blacs_get,blacs_gridinit,blacs_gridinfo,blacs_exit,blacs_pinfo,blacs_barrier,blacs_exit,blacs_pnum,blacs_gridmap,%
      pdelget,pdelset,pdtran,dgemm,pdgemm,numroc,descinit,zheev,pzheev,dsyev,pdsyev,%
      mpiu_init,mpiu_end,mpiu_sync,mpiu_abort,mpiu_world_set,mpiu_bcast,mpiu_scatter,mpiu_gather,mpiu_sndrcv,mpiu_sum,mpiu_prod,%
      indxl2g,indxg2l,indxg2p},%
   deletecomment=[f],% no fixed comment line: 1998 Magne Rudshaug
   morecomment=[l]!%
  }%
%% csh definition (c) 1998 Kai Below
%%
\lstdefinelanguage{csh}
  {morekeywords={alias,awk,cat,echo,else,end,endif,endsw,exec,exit,%
      foreach,glob,goto,history,if,logout,nice,nohup,onintr,repeat,sed,%
      set,setenv,shift,source,switch,then,time,while,umask,unalias,%
      unset,wait,while,@,env,argv,child,home,ignoreeof,noclobber,%
      noglob,nomatch,path,prompt,shell,status,verbose,print,printf,%
      sqrt,BEGIN,END,gfortran,ls,rm,less,chmod,perl,cd,pwd,mkdir,rmdir,%
      sed,cp,mv,which,where,whatis,cmd,man,info,locate,head,tail,sort,grep,%
      gawk,emacs,du,df,top,ps,sudo,apt-get},%
   morecomment=[l]\#,%
   morestring=[d]"%
  }[keywords,comments,strings]%

% %%%%%%%%%%%%%%%%%%%%%%%%%%%%%%%%%%%%%%%%%%%%%%%%%%%%%%%%%
%   Making emacs refcard table in 01computer.tex
% %%%%%%%%%%%%%%%%%%%%%%%%%%%%%%%%%%%%%%%%%%%%%%%%%%%%%%%%%
% We assume 3 columns in the table. Change multicolumn commands and &
% appropriately if changed
\usepackage{longtable}
\newcommand{\emtentry}   [1]{      \multicolumn{3}{l}{\parbox[t]{10cm}{\scriptsize{\eng{#1}}}}\\}
\newcommand{\emtsection} [1]{\hline\multicolumn{3}{l}{\eng{\bfseries{{#1}}}}\\\hline}
\newcommand{\emtkey}     [2]{\eng{\scriptsize{#1}} & \engt{\scriptsize{#2}} & \\}
\newcommand{\kbd}        [1]{{\tt\scriptsize {#1}}\null}
\newcommand{\emtthreecol}[3]{\eng{\scriptsize{#1}}&\engt{\scriptsize{#2}}&\engt{\scriptsize{#3}}\\}
%\newcommand{\emtkey}[2]{\eng{#1} & \multicolumn{2}{l}{\engt{#2}}\\}
% refer to problems: askhseis:
\newcommand{\pref}[1]{\thechapter.\ref{#1}}

\newenvironment{latintext}{}{}
%\newcommand{\latintext}{}
\newcommand{\greektext}{}

\usepackage{rotating}
% rotate tables:
%%\begin{sidewaystable}
%%    \centering
%%    \caption{Your caption here}
%%   \begin{tabular}{ll}
%%    First First & First Second\\
%%    Second First & Second Second
%%    \end{tabular}
%%\end{sidewaystable}


% Load custom documentation box and man-page styling macros
% ===================================================================
% Custom Macros and Styling for matrix_mod_matrix Extended Documentation
% ===================================================================

% Color Palette
\definecolor{modPrimary}{RGB}{22, 64, 122}      % Deep Navy Blue
\definecolor{modSecondary}{RGB}{31, 111, 139}   % Teal / Slate Blue
\definecolor{modAccent}{RGB}{168, 42, 42}       % Crimson / Burgundy for Procedure Names
\definecolor{modArg}{RGB}{0, 102, 51}           % Forest Green for Dummy Arguments
\definecolor{modType}{RGB}{102, 51, 153}        % Royal Purple for Types / Intents
\definecolor{modBoxBg}{RGB}{244, 247, 251}      % Very Light Ice Blue background
\definecolor{modBoxFrame}{RGB}{44, 88, 150}     % Medium Navy frame
\definecolor{modWarnBg}{RGB}{254, 249, 238}     % Warm Amber background
\definecolor{modWarnFrame}{RGB}{204, 122, 0}    % Amber frame
\definecolor{modCodeBg}{RGB}{242, 244, 247}     % Clean light slate gray for listings
\definecolor{modSynBg}{RGB}{236, 241, 248}      % Soft blue-gray for Synopsis blocks

% Enhance default listings appearance while preserving mymacros.tex settings
\lstset{
    backgroundcolor=\color{modCodeBg},
    rulecolor=\color{modBoxFrame!50},
    frame=single,
    basicstyle=\footnotesize\ttfamily,
    keywordstyle=\color[rgb]{0,0,0.75}\bfseries,
    morekeywords={class,extends,generic,pass,type,procedure,aimag,cmplx,dot_product,maxval,minval,sqrt},
    commentstyle=\color[rgb]{0.13,0.50,0.13}\itshape,
    stringstyle=\color[rgb]{0.58,0.10,0.75},
    upquote=true,
    breaklines=true,
    columns=fixed,
    showstringspaces=false
}

% Typography Macros for Procedures, Arguments, Types, and Files
\newcommand{\procname}[1]{\textcolor{modAccent}{\texttt{\textbf{#1}}}}
\newcommand{\argname}[1]{\textcolor{modArg}{\texttt{\textbf{#1}}}}
\newcommand{\typename}[1]{\textcolor{modType}{\texttt{#1}}}
\newcommand{\intentname}[1]{\textcolor{modSecondary}{\texttt{#1}}}
\newcommand{\modname}[1]{\textcolor{modPrimary}{\texttt{\textbf{#1}}}}
\newcommand{\filename}[1]{\textcolor{modSecondary}{\texttt{#1}}}

% Man-Page Section Label Macro
\newcommand{\manheading}[1]{%
  \par\nopagebreak[4]\vspace{0.35em}\nopagebreak[4]\noindent%
  {\sffamily\bfseries\footnotesize\textcolor{modPrimary}{\MakeUppercase{#1}}}\par\nopagebreak[4]\vspace{0.15em}\nopagebreak[4]\noindent%
}

% Colored Informational & Warning Boxes
\newtcolorbox{notebox}[1][Note]{
  enhanced,
  breakable,
  colback=modBoxBg,
  colframe=modBoxFrame,
  coltitle=white,
  fonttitle=\sffamily\bfseries\small,
  title={#1},
  arc=1.5mm,
  boxrule=0.6pt,
  left=8pt, right=8pt, top=6pt, bottom=6pt
}

\newtcolorbox{warnbox}[1][Important]{
  enhanced,
  breakable,
  colback=modWarnBg,
  colframe=modWarnFrame,
  coltitle=white,
  fonttitle=\sffamily\bfseries\small,
  title={#1},
  arc=1.5mm,
  boxrule=0.6pt,
  left=8pt, right=8pt, top=6pt, bottom=6pt
}

% Visual Subgroup Banner inside Operators / Procedures sections (not added to TOC)
\newcommand{\mansubgroup}[1]{%
  \par\vspace{1.0em}\noindent%
  \colorbox{modPrimary!10}{%
    \parbox{\dimexpr\linewidth-2\fboxsep\relax}{%
      \sffamily\bfseries\normalsize\textcolor{modPrimary}{#1}%
    }%
  }\par\nopagebreak[4]\vspace{0.4em}%
}

% Man-Page Procedure Entry Header Box & TOC Registration
% Usage: \manentry{kind}{proc_display_name}{toc_name}{brief_summary}{label_id}
\newcommand{\manentry}[5]{%
  \par\vspace{1.2em}%
  \phantomsection%
  \label{#5}%
  \addcontentsline{toc}{subsubsection}{\texorpdfstring{\texttt{\textbf{#3}} \textnormal{\footnotesize--- #4}}{#3}}%
  \begin{tcolorbox}[
    enhanced,
    colback=modBoxBg,
    colframe=modPrimary,
    boxrule=0.8pt,
    leftrule=4.5pt,
    arc=1.5mm,
    left=8pt, right=8pt, top=6pt, bottom=6pt,
    before skip=0.5em,
    after={\par\nopagebreak[4]\vspace{0.4em}\nopagebreak[4]}
  ]
  \noindent{\sffamily\bfseries\scriptsize\textcolor{modSecondary}{\MakeUppercase{#1}}}\hfill{\sffamily\scriptsize\textcolor{gray}{matrix\_mod\_matrix(3f)}}\\[2pt]
  {\large\procname{#2}} \quad {\normalsize --- \textbf{#4}}
  \end{tcolorbox}%
  \noindent\textbf{\sffamily\footnotesize\textcolor{modPrimary}{NAME}}\par\nopagebreak[4]\vspace{0.1em}\nopagebreak[4]\noindent%
  \hspace{1.2em}\procname{#2} --- #4\par\nopagebreak[4]%
}

% Argument Item Macro inside Dummy Arguments list
% Usage: \argitem{arg_name}{type_and_intent_spec}{description}
\newcommand{\argitem}[3]{%
  \item[\argname{#1}] (\typename{#2}): #3%
}


\setcounter{secnumdepth}{3}
\setcounter{tocdepth}{3}

\title{%
  \vspace{-1.2cm}%
  \begin{tcolorbox}[
    enhanced,
    colback=modPrimary,
    colframe=modPrimary!80!black,
    coltext=white,
    arc=2mm,
    boxrule=1pt,
    top=12pt, bottom=12pt, left=12pt, right=12pt
  ]
  \centering
  {\LARGE\bfseries Matrix Module Documentation}\\[6pt]
  {\large Extended Reference Manual for the Fortran \texttt{matrix\_mod\_matrix} Module}
  \end{tcolorbox}%
  \vspace{-0.4cm}%
}
\author{\textbf{Konstantinos N. Anagnostopoulos}\\
  {\small Physics Department, National Technical University of Athens}\\
  {\small \texttt{konstant@mail.ntua.gr} $\cdot$ \url{https://www.physics.ntua.gr/konstant}}}
\date{\small \today}

\begin{document}

\maketitle

\begin{abstract}
\noindent
This document provides the complete reference manual for the public derived types, constructors, type-bound (component) procedures, overloaded operators, and generic procedures of the Fortran \modname{matrix\_mod\_matrix} module in the \texttt{Matrix} library (\filename{Software/Tools/Matrix/}). Built as an object-oriented layer on top of \modname{matrix\_mod\_array} (\modname{array\_mod}), \texttt{BLAS}, \texttt{LAPACK}, and \texttt{PFAPACK}, \modname{matrix\_mod\_matrix} introduces four high-level derived types---\typename{Matrix} and \typename{DMatrix} for double-precision complex and real matrices, and \typename{Vector} and \typename{DVector} for double-precision complex and real vectors---with built-in dimension, index-bound, symmetry-tag, and label metadata tracking. It supports natural algebraic notation via overloaded operators (\procname{=}, \procname{+}, \procname{-}, \procname{*}, \procname{/}, \procname{**}), zero-allocation \texttt{BLAS} multiplication subroutines (\procname{mmmult}, \procname{mvmult}, \procname{vmmult}), matrix inversion, determinants, Pfaffians, spectral decompositions, element-wise mathematical functions, random matrix generation, and formatted terminal/file/Mathematica I/O. Every public entity is documented in Unix man-page style with its name, synopsis, brief summary, detailed behavior, dummy arguments (or data components), and usage examples.
\end{abstract}

\vspace{0.3cm}
\begin{tcolorbox}[
  enhanced,
  breakable,
  colback=modBoxBg!60,
  colframe=modPrimary!60,
  arc=1.5mm,
  boxrule=0.6pt,
  left=10pt, right=10pt, top=6pt, bottom=8pt
]
\tableofcontents
\end{tcolorbox}
\newpage

% Section 1: Usage, Compiling and Linking, Download, Module-Wide Utilities
% ===================================================================
% Section 1: Usage, Compiling and Linking, Download, Utilities
% ===================================================================
\section{Usage}
\label{sec:usage}

The \modname{matrix\_mod\_matrix} module (part of the \texttt{Matrix} library) provides an object-oriented, high-level modern Fortran interface for linear algebra calculations using encapsulated matrix and vector derived types. Built on top of \modname{matrix\_mod\_array} (whose standalone array API is exported via \modname{array\_mod}), \modname{matrix\_mod\_matrix} wraps performance-critical \texttt{BLAS}, \texttt{LAPACK}, and \texttt{PFAPACK} routines behind intuitive type constructors, type-bound methods, and overloaded mathematical operators.

Unlike \modname{array\_mod}, which operates directly on raw Fortran arrays, \modname{matrix\_mod\_matrix} manages four public derived types that store both numerical values and descriptive metadata:
\begin{itemize}
  \item \typename{type(Matrix)} and \typename{type(DMatrix)}: Double-precision complex (\typename{complex(8)}) and real (\typename{real(8)}) matrices carrying dimensions $(m,n)$, arbitrary lower/upper index bounds (\texttt{is:ie}, \texttt{js:je}), a symmetry flag (\texttt{mtype}: \texttt{'ZG'}/\texttt{'ZH'} or \texttt{'DG'}/\texttt{'DS'}), and an optional label (\texttt{name}).
  \item \typename{type(Vector)} and \typename{type(DVector)}: Double-precision complex (\typename{complex(8)}) and real (\typename{real(8)}) vectors carrying length $n$, arbitrary index bounds (\texttt{is:ie}), and an optional label (\texttt{name}).
  \item \textbf{Operator Overloading \& Zero-Allocation Subroutines}: Full support for assignment (\procname{=}), addition (\procname{+}), subtraction (\procname{-}), scalar/matrix/vector multiplication (\procname{*}), scalar division (\procname{/}), and element-wise exponentiation (\procname{**}), alongside low-overhead \texttt{BLAS} subroutines (\procname{mmmult}, \procname{mvmult}, \procname{vmmult}).
  \item \textbf{Symmetry-Aware Linear Algebra \& I/O}: Constructors, uniform/Gaussian random generators, symmetry projections and tests, traces, norms, inverses, determinants (\procname{determinant}, \procname{lndet}), Pfaffians (\procname{Pfaffian}, \procname{lnPfaffian}), eigensolvers (\procname{eigenvalues}, \procname{eigenvectors}), element-wise math intrinsics, and formatted terminal/file/Mathematica I/O.
\end{itemize}

\begin{notebox}[Ergonomics vs.\ Raw Array Performance]
The derived types \typename{Matrix}, \typename{DMatrix}, \typename{Vector}, and \typename{DVector} are designed primarily for expressive, readable scientific code and educational use. Operator overloading simplifies matrix algebra at the cost of temporary allocations for small arrays ($N \lesssim 200$--$500$). For maximum throughput in compute-intensive inner loops, use the subroutine interfaces (\procname{mmmult}, \procname{mvmult}, \procname{vmmult}) or the raw array module \modname{array\_mod} (see Section~\ref{sec:performance} for detailed benchmarks).
\end{notebox}

A minimal motivating program demonstrating Gaussian random matrix/vector construction, overloaded matrix--vector multiplication, eigenvalue computation, and formatted output is shown below:
\lstset{language=[95]Fortran}
\begin{lstlisting}
program hello_world
  use :: matrix_mod_matrix             ! The name of the module
  implicit none
  type(Matrix) :: m
  type(Vector) :: v1, v2
  integer      :: n

  call matrix_random_init              ! seed PRNG from /dev/urandom
  f_mout = 6                           ! unit for the module messages

  n  = 4
  m  = Matrix('g', n)                  ! Gaussian random n x n matrix
  v1 = Vector('g', n)                  ! Gaussian random n-vector

  v2 = m * v1                          ! matrix-vector multiplication
  v2 = eigenvalues(m)                  ! compute eigenvalues of m
  call v2 % print(f_mout)              ! print the eigenvalues
end program hello_world
\end{lstlisting}

% -------------------------------------------------------------------
\subsection{Compiling and Linking}
\label{subsec:compiling}

To use \modname{matrix\_mod\_matrix}, you must first compile the \texttt{Matrix} library archive (\filename{libMatrix.a}) and its Fortran module description files (\filename{*.mod}), and then link your program against \texttt{libMatrix}, \texttt{LAPACK}, and \texttt{BLAS}.

If your program file \filename{hello.f90} is placed inside \filename{Software/Tools/Matrix/}, build the library and compile using \texttt{gfortran} with:
\lstset{language=bash}
\begin{lstlisting}
make gnu
gfortran -I. hello.f90 -L. -lMatrix -llapack -lblas -o hello
\end{lstlisting}
Here \lcodeb{make gnu} builds the bundled \texttt{PFAPACK} library first, compiles all modules (including \filename{matrix\_mod\_common.f90}, \filename{matrix\_mod\_array.f90}, and \filename{matrix\_mod\_matrix.f90}), archives everything into \filename{libMatrix.a}, and copies the artifacts into \filename{lib/gnu/}. Once \texttt{PFAPACK} has already been compiled, running \lcodeb{make lib} suffices to recompile module objects and rebuild \filename{libMatrix.a}.

When compiling a program outside the library tree, pass the paths to the compiled \filename{.mod} and \filename{libMatrix.a} directory (e.g.\ \filename{/path/to/Matrix/}):
\lstset{language=bash}
\begin{lstlisting}
gfortran -I/path/to/Matrix hello.f90 -L/path/to/Matrix -lMatrix \
         -llapack -lblas -o hello
\end{lstlisting}

\begin{notebox}[Makefile Targets and Compiler Options]
The included \filename{Makefile} provides tuned configurations for GNU, Intel, AMD, and NVIDIA toolchains:
\begin{itemize}
  \item \textbf{GNU Fortran (\texttt{gfortran})}: Run \lcodeb{make gnu}. Default variables:
\lstset{language=bash}
\begin{lstlisting}
FC     = gfortran
FFLAGS = -O3 -mtune=native -funroll-loops -fopenmp -pthread
LIBS   = -llapack -lblas
\end{lstlisting}
  If \texttt{libMatrix.a} was compiled with OpenMP (\texttt{-fopenmp}), pass \texttt{-fopenmp} when linking executables and set \lcodeb{export OMP_NUM_THREADS=8} at runtime as needed.
  \item \textbf{Intel Fortran (\texttt{ifx} / \texttt{ifort} with MKL)}: Run \lcodeb{make intel} (\texttt{-O3 -xHost -qopenmp -qmkl=parallel}), saving artifacts to \filename{lib/intel/}.
  \item \textbf{HPC Cluster Targets}: Targets are included for AMD AOCC/\texttt{flang} with AOCL (\lcodeb{make aris}), GNU with Intel MKL (\lcodeb{make arismkl}), and NVIDIA HPC SDK \texttt{nvfortran} (\lcodeb{make nvaris}, \lcodeb{make icarus}, \lcodeb{make icarusnogpu}).
  \item \textbf{Debug Builds}: Run \lcodeb{make debug} (or \lcodeb{make d}) to compile with \texttt{FFLAGS=-g}.
\end{itemize}
\end{notebox}

% -------------------------------------------------------------------
\subsection{Download}
\label{subsec:download}

The complete source code archive of the \texttt{Matrix} library can be downloaded from:
\begin{center}
  \url{https://physics.ntua.gr/konstant/PUB/Matrix.tgz}
\end{center}
The official GitHub repository is maintained at:
\begin{center}
  \url{https://github.com/KNAnagnostop/FortranMatrix}
\end{center}

% -------------------------------------------------------------------
\subsection{Module-Wide Error Handling and \texorpdfstring{\texttt{mtype}}{mtype} Conventions}
\label{subsec:error_mtype}

Module \modname{matrix\_mod\_matrix} distinguishes two kinds of failure.

\textbf{Shape errors abort.} A wrong shape is a programming error: it depends on the code, not on the data. The program stops at once with a non-zero exit status and a message naming the procedure where it happened. Ordinary procedures call \texttt{matrix\_error}, which prints a banner; \texttt{pure} procedures use \texttt{error stop '<procedure>: ...'}, which is allowed in \texttt{pure} code since Fortran~2018. Shape errors are:
\begin{itemize}
  \item \procname{+} and \procname{-} with non-conformable \typename{Matrix}, \typename{DMatrix}, or array operands;
  \item \procname{*} (matrix--matrix and matrix--vector), \procname{mmmult}, \procname{mvmult}, and \procname{vmmult} when the inner dimensions differ;
  \item \typename{Vector} and \typename{DVector} \procname{+}, \procname{-}, \procname{dot\_product}, and array or vector assignment with a size mismatch;
  \item \procname{eigenvalues}, \procname{eigenvectors}, \procname{hermitian\_set}, \procname{symmetric\_set}, and \procname{antisymmetric\_set} on a non-square matrix;
  \item scalar assignment to an unconstructed \typename{Matrix}, \typename{DMatrix}, \typename{Vector}, or \typename{DVector}.
\end{itemize}

\textbf{Data-dependent numerical failures return \procname{NaN}.} These can occur in a correct program, for example for one configuration of a Monte Carlo simulation, so the program continues and can test the result with \procname{isNaN}: a singular matrix in \procname{inverse}, and a \texttt{LAPACK} convergence failure in \procname{eigenvalues} and \procname{eigenvectors}.

{\sloppy The functions \procname{trace}, \procname{trace2}, \procname{trace2c}, \procname{traceless}, \procname{isHermitian}, \procname{isSymmetric}, \procname{isAntiSymmetric}, \procname{inverse}, \procname{determinant}, \procname{lndet}, \procname{Pfaffian}, and \procname{lnPfaffian} pass their argument directly to \modname{matrix\_mod\_array}, and therefore still return \procname{NaN} for a non-square matrix, following the \modname{array\_mod} convention.\par}

\textbf{\texttt{mtype} semantics.} \texttt{'ZH'} (\typename{Matrix}) and \texttt{'DS'} (\typename{DMatrix}) select the Hermitian or symmetric \texttt{BLAS} and \texttt{LAPACK} kernels (\texttt{zhemm}, \texttt{zhemv}, \texttt{zheev} and \texttt{dsymm}, \texttt{dsymv}, \texttt{dsyev}), which read only the upper triangle of the stored array. Operations that may destroy hermiticity or symmetry reset \texttt{mtype} to \texttt{'ZG'} or \texttt{'DG'}; call \procname{hermitian\_set} or \procname{symmetric\_set} to re-tag a matrix that is known to be Hermitian or symmetric.

% -------------------------------------------------------------------
\subsection{Module-Wide Public Variable and Utilities}
\label{subsec:mod_utilities}

Before introducing the four derived types in detail, this subsection documents the module-wide public variable \procname{f\_mout} and the two standalone utilities \procname{matrix\_random\_init} and \procname{NaN}, which are shared across all classes of \modname{matrix\_mod\_matrix}.

% ===================================================================
% f_mout
% ===================================================================
\manentry{Public Module Variable}{f\_mout}{f\textunderscore mout}{Default Fortran logical I/O unit used for module output, saving, and messages}{mod:f_mout}

\manheading{Synopsis}
\lstset{language=[95]Fortran}
\begin{lstlisting}
integer :: f_mout = stdout   ! output_unit from iso_fortran_env (typically 6)
\end{lstlisting}

\manheading{Brief}
Specifies the default Fortran I/O unit number for output produced by \procname{\% print}, \procname{\% save}, and \procname{matrix\_random\_init}.

\manheading{Description}
Declared in \modname{matrix\_mod\_common} and re-exported publicly by \modname{matrix\_mod\_matrix}, \procname{f\_mout} is initialized to \texttt{stdout} (\texttt{output\_unit} from \texttt{iso\_fortran\_env}, usually unit \texttt{6}). Whenever the optional \argname{unit} dummy argument is omitted in a call to \procname{MAT \% print}, \procname{MAT \% save}, \procname{vec \% print}, or \procname{vec \% save}, output is directed to \procname{f\_mout}. Users may reassign \procname{f\_mout} to an open file unit globally rather than passing \argname{unit=...} to every call. Note that the companion input unit \texttt{f\_minput} (initialized to \texttt{stdin}) is \emph{not} exported publicly by \modname{matrix\_mod\_matrix}, so use \texttt{call obj \% read(unit=un)} to read from an explicit file unit.

\manheading{Specification}
\begin{description}[leftmargin=2.2cm, style=nextline]
  \argitem{f\_mout}{integer, public module variable, saved}{Default output unit number (initial value: \texttt{stdout} / \texttt{iso\_fortran\_env:output\_unit}).}
\end{description}

\manheading{Example}
\lstset{language=[95]Fortran}
\begin{lstlisting}
use :: matrix_mod_matrix
type(Matrix) :: m

m = Matrix('gaussian', 4, name='G')
call m % print()        ! writes to default stdout (f_mout = 6)

open(newunit=f_mout, file='out.dat', status='replace')
call m % print()        ! now writes to out.dat
close(f_mout)
f_mout = 6              ! restore stdout
\end{lstlisting}

% ===================================================================
% matrix_random_init
% ===================================================================
\manentry{Subroutine}{matrix\_random\_init}{matrix\textunderscore random\textunderscore init}{Seeds Fortran's intrinsic PRNG from \texttt{/dev/urandom} with a portable fallback}{mod:matrix_random_init}

\manheading{Synopsis}
\lstset{language=[95]Fortran}
\begin{lstlisting}
call matrix_random_init()
\end{lstlisting}

\manheading{Brief}
Initializes Fortran's pseudo-random number generator (\texttt{random\_seed}) using high-entropy bytes from \filename{/dev/urandom}, or falls back gracefully to \lcodef{call random_seed()}.

\manheading{Description}
\procname{matrix\_random\_init} queries the processor seed array size with \lcodef{call random_seed(size=nseeds)}, opens \filename{/dev/urandom} as an unformatted stream file, and reads \texttt{nseeds} integers. To ensure valid positive non-degenerate seeds across compilers, it clears the sign bit with \lcodef{iand(seeds, huge(seeds))} and forces the least-significant bit of \texttt{seeds(1)} to $1$ with \lcodef{ior(seeds(1), 1)} before calling \lcodef{call random_seed(put=seeds)}. An informational status message is written to \procname{f\_mout}.

If \filename{/dev/urandom} is unavailable or unreadable (for example, in restricted containers or on Windows), \procname{matrix\_random\_init} prints a warning to \procname{f\_mout} and falls back cleanly to \lcodef{call random_seed()} without aborting.

\manheading{Dummy Arguments}
\begin{description}[leftmargin=2.2cm, style=nextline]
  \argitem{(none)}{}{This subroutine takes no arguments.}
\end{description}

\manheading{Example}
\lstset{language=[95]Fortran}
\begin{lstlisting}
use :: matrix_mod_matrix
type(Matrix) :: m

call matrix_random_init
m = Matrix('gaussian', 8, sigma=1.0_8)
\end{lstlisting}

% ===================================================================
% NaN
% ===================================================================
\manentry{Pure Generic Function}{NaN}{NaN}{Returns an IEEE quiet Not-a-Number scalar for real or complex assignments}{mod:NaN}

\manheading{Synopsis}
\lstset{language=[95]Fortran}
\begin{lstlisting}
z = NaN()    ! returns complex(8) with both Re and Im set to IEEE quiet NaN
z = NaN(z0)  ! returns complex(8) with both Re and Im set to IEEE quiet NaN
r = NaN(r0)  ! returns real(8) IEEE quiet NaN
\end{lstlisting}

\manheading{Brief}
Produces a double-precision real or complex IEEE quiet \texttt{NaN} (\texttt{IEEE\_QUIET\_NAN}) via Fortran's intrinsic \texttt{ieee\_arithmetic} module.

\manheading{Description}
\procname{NaN} provides three pure overloads exported from \modname{matrix\_mod\_common}:
\begin{itemize}
  \item \lcodef{NaN()} (no arguments) and \lcodef{NaN(z)} (\typename{complex(8)} argument) return a \typename{complex(8)} value \texttt{CMPLX(qnan, qnan, kind=dp)}, setting \emph{both} real and imaginary parts to IEEE quiet \texttt{NaN}.
  \item \lcodef{NaN(r)} (\typename{real(8)} argument) returns a \typename{real(8)} IEEE quiet \texttt{NaN}. Note that assigning \cbox{m \% v(i,j) = NaN(r)} to a complex entry sets its real part to \texttt{NaN} and its imaginary part to \texttt{0.0\_8} (and still causes \procname{isNaN} to return \texttt{.true.}).
\end{itemize}
When assigning \texttt{m \% v = NaN()} or \texttt{v \% v = NaN()} to an entire array component, Fortran broadcasts the scalar \texttt{NaN} across all elements. Always use \procname{isNaN} to test for \texttt{NaN}, never equality (\texttt{==} / \texttt{/=}).

\manheading{Dummy Arguments}
\begin{description}[leftmargin=2.2cm, style=nextline]
  \argitem{z / r}{complex(8) or real(8), scalar, intent(in), optional}{Type-selecting mold argument (when omitted, returns \typename{complex(8)}).}
  \argitem{result}{complex(8) or real(8), scalar, return value}{IEEE quiet \texttt{NaN} value matching the selected precision and type.}
\end{description}

\manheading{Example}
\lstset{language=[95]Fortran}
\begin{lstlisting}
use :: matrix_mod_matrix
type(Matrix)  :: m1
type(DMatrix) :: d1
real(8)       :: r = 0.0_8

m1 = Matrix(4)
d1 = DMatrix(4)
m1 % v      = NaN()       ! set both Re and Im of all entries to NaN
m1 % v(1,1) = NaN(r)      ! set Re part of entry (1,1) to NaN
d1 % v      = NaN(r)      ! set all real entries of d1 to NaN
if (isNaN(m1)) print *, 'm1 contains NaN'
\end{lstlisting}


% Section 2: Class Matrix (Components, Constructors, Component Procedures, Operators, Procedures)
% ===================================================================
% Section 2: Class Matrix (Components, Constructors, Component Procedures, Operators)
% ===================================================================
\section{Class \texttt{Matrix}}
\label{sec:class_matrix}

The \typename{Matrix} derived type provides a structured interface for double-precision complex (\typename{complex(8)}) matrix manipulation. It extends a private base type \texttt{MatrixClass} that holds common matrix dimensions, custom index bounds, symmetry classification, label metadata, and shared I/O and random-filling methods.

% -------------------------------------------------------------------
\subsection{Components}
\label{subsec:mat_components}

% ===================================================================
% type(Matrix)
% ===================================================================
\manentry{Derived Type}{type(Matrix)}{Matrix (type)}{Double-precision complex matrix class extending the private \texttt{MatrixClass}}{mat:type}

\manheading{Synopsis}
\lstset{language=[95]Fortran}
\begin{lstlisting}
type(Matrix) :: mat
! Access elements and metadata:
! mat % m, mat % n, mat % is, mat % ie, mat % js, mat % je
! mat % mtype, mat % name, mat % v(is:ie, js:je)
\end{lstlisting}

\manheading{Brief}
Encapsulates an allocatable $m \times n$ \typename{complex(8)} array \texttt{v(is:ie, js:je)} together with dimension, index-bound, symmetry-type, and name metadata.

\manheading{Description}
An object of type \typename{Matrix} stores an $m \times n$ complex matrix whose row indices span \texttt{is:ie} (\texttt{ie = is + m - 1}) and whose column indices span \texttt{js:je} (\texttt{je = js + n - 1}). The default constructor values initialize \texttt{m = n = is = ie = js = je = 0} and \texttt{mtype = 'GG'}, indicating an uninitialized matrix. Once constructed, \texttt{mtype} is either \texttt{'ZG'} (general double-precision complex matrix) or \texttt{'ZH'} (Hermitian complex matrix, $M^\dagger = M$). Whenever \texttt{mtype = 'ZH'}, routines such as \procname{*}, \procname{mmmult}, \procname{mvmult}, \procname{trace2}, \procname{trace2c}, \procname{eigenvalues}, and \procname{eigenvectors} automatically dispatch to specialized Hermitian \texttt{BLAS}/\texttt{LAPACK} kernels (which read only the upper triangle) or $O(N^2)$ formulas.

The parent type \texttt{MatrixClass} is private to \modname{matrix\_mod\_matrix} and cannot be instantiated directly; all user matrices are declared as \lcodef{type(Matrix)} or \lcodef{class(Matrix)}.

\manheading{Data Components}
\begin{description}[leftmargin=2.2cm, style=nextline]
  \argitem{m, n}{integer, default 0}{Number of rows (\texttt{m}) and columns (\texttt{n}) of the matrix; \texttt{m == 0} indicates uninitialized state.}
  \argitem{is, ie}{integer, default 0}{Lower and upper row index bounds of \texttt{v}, satisfying \texttt{ie = is + m - 1}.}
  \argitem{js, je}{integer, default 0}{Lower and upper column index bounds of \texttt{v}, satisfying \texttt{je = js + n - 1}.}
  \argitem{mtype}{character(4), default \texttt{'GG'}}{\texttt{LAPACK}-style symmetry identifier (\texttt{mtype\_len = 4}): \texttt{'ZG'} for general complex matrices, \texttt{'ZH'} for Hermitian matrices, or \texttt{'GG'} when uninitialized.}
  \argitem{name}{character(20), default \texttt{''}}{User-defined matrix label (\texttt{name\_len = 20}) printed by \procname{\% print}.}
  \argitem{v(:,:)}{complex(8), allocatable}{Allocated rank-2 complex array storing matrix elements with bounds \texttt{v(is:ie, js:je)}.}
\end{description}

\manheading{Example}
\lstset{language=[95]Fortran}
\begin{lstlisting}
use :: matrix_mod_matrix
type(Matrix) :: mat

mat = Matrix(4, 3, is=0, js=5, name='A')
print *, mat % m, mat % n               ! 4, 3
print *, mat % is, mat % ie             ! 0, 3
print *, mat % js, mat % je             ! 5, 7
print *, lbound(mat % v), ubound(mat%v) ! 0, 5 and 3, 7
\end{lstlisting}

% -------------------------------------------------------------------
\subsection{Constructors}
\label{subsec:mat_constructors}

% ===================================================================
% Matrix(...) constructor
% ===================================================================
\manentry{Generic Constructor}{Matrix(...)}{Matrix (constructor)}{Constructs and allocates a \texttt{type(Matrix)} object (zero, constant, random, identity, or from 2D array)}{mat:constructor}

\manheading{Synopsis}
\lstset{language=[95]Fortran}
\begin{lstlisting}
mat = Matrix(m, [n], [is], [js], [mtype], [name])
mat = Matrix(C,      [is], [js], [mtype], [name])
mat = Matrix(r, m, [n], [is], [js], [mtype], [name])
mat = Matrix(c, m, [n], [is], [js], [mtype], [name])
mat = Matrix(rtype, m, [n], [is], [js], [mtype], [name], [sigma])
\end{lstlisting}

\manheading{Brief}
Allocates a \typename{type(Matrix)} instance, sets its metadata and index bounds, and initializes \texttt{mat \% v} to zero, a constant scalar, a 2D array \argname{C}, an identity matrix, or uniform/Gaussian random numbers.

\manheading{Description}
The structure constructor interface \procname{Matrix} overloads five construction forms, all sharing common metadata defaults set by the internal helper \texttt{matrix\_metadata\_put}:
\begin{enumerate}
  \item \textbf{Zero Matrix \texttt{Matrix(m, ... )}}: Allocates an $m \times n$ matrix filled with \texttt{(0.0\_8, 0.0\_8)}. Only \argname{m} is required; defaults are:
  \begin{itemize}
    \item \texttt{n = m}, \texttt{is = 1}, \texttt{js = 1}, \texttt{ie = m}, \texttt{je = n}, \texttt{mtype = 'ZG'}, \texttt{name = ''}.
    \item If \argname{is} is supplied without \argname{js}, then \texttt{js} automatically defaults to \argname{is} (\texttt{js = is}, \texttt{je = is + n - 1}).
    \item If \argname{js} is supplied, only column bounds change (\texttt{js = js}, \texttt{je = js + n - 1}), leaving \texttt{is, ie} unaffected.
    \item Passing $m \le 0$, $n \le 0$, an unrecognized \argname{mtype} (only \texttt{'ZG'} and \texttt{'ZH'} are valid), or \texttt{mtype = 'ZH'} with $m \neq n$ aborts via \texttt{matrix\_error}.
  \end{itemize}
  \item \textbf{From 2D Array \texttt{Matrix(C, ... )}}: Initializes \texttt{mat \% v} from an assumed-shape \typename{complex(8)} array \argname{C(:,:)} with $m = \texttt{size(C,1)}$ and $n = \texttt{size(C,2)}$. Because \argname{C} is passed as assumed-shape without \texttt{allocatable}, Fortran resets lower bounds to \texttt{is = js = 1} unless \texttt{is=lbound(C,1)} and \texttt{js=lbound(C,2)} are passed explicitly. If \texttt{mtype = 'ZH'}, \cbox{call mat \% hermitian\_set} is invoked automatically without checking whether \argname{C} was already Hermitian.
  \item \textbf{Real Scalar Fill \texttt{Matrix(r, m, ... )}}: Constructs the $m \times n$ matrix and sets all entries to \typename{real(8)} scalar \argname{r}. Note that the scalar argument precedes the dimension \argname{m}.
  \item \textbf{Complex Scalar Fill \texttt{Matrix(c, m, ... )}}: Constructs the $m \times n$ matrix and sets \texttt{mat \% v = c}. If \texttt{mtype = 'ZH'}, \cbox{call mat \% hermitian\_set} is called afterwards. A constant matrix with entries $c$ is Hermitian only if $c$ is real, and the averaging $\frac{1}{2}(M + M^\dagger)$ replaces every entry (diagonal and off-diagonal) by $\frac{1}{2}(c + c^*) = \Re(c)$, so the result is the real matrix \texttt{mat \% v = REAL(c, kind=8)}.
  \item \textbf{Random or Identity Matrix \texttt{Matrix(rtype, m, ... )}}: Dispatches on the first character \argname{rtype(1:1)}:
  \begin{itemize}
    \item \texttt{'r'}, \texttt{'R'}, \texttt{'u'}, or \texttt{'U'}: Fills with uniform random complex numbers in $[0,1) + i[0,1)$ via \procname{mat \% random}. If \texttt{mtype = 'ZH'}, the lower triangle is conjugated into the upper triangle (so upper-triangle imaginary parts are negative and all entries remain uniformly distributed).
    \item \texttt{'g'} or \texttt{'G'}: Fills with Gaussian random complex numbers of standard deviation \argname{sigma} (default \texttt{1.0\_8}) via \procname{mat \% gaussian}. If \texttt{mtype = 'ZH'}, variance is preserved by scaling off-diagonals and the diagonal properly.
    \item \texttt{'1'}: Constructs the $m \times m$ complex identity matrix (aborts via \texttt{matrix\_error} if $m \neq n$).
    \item Passing an empty string (\texttt{len(rtype) < 1}) or any other leading character aborts via \texttt{matrix\_error}.
  \end{itemize}
\end{enumerate}

\begin{warnbox}[Performance Tip for Array Construction]
As noted in \filename{matrix\_mod\_matrix\_matrix.f90}, repeated calls to \lcodef{mat = Matrix(C, ...)} inside tight loops allocate and copy an extra temporary \typename{Matrix}. When \texttt{mat} already has the right size, prefer constructing it once (\lcodef{mat = Matrix(m, n)}) and assigning \lcodef{mat = C} or \cbox{mat \% v = C} in-place.
\end{warnbox}

\manheading{Dummy Arguments}
\begin{description}[leftmargin=2.2cm, style=nextline]
  \argitem{m}{integer, intent(in)}{Number of rows ($m > 0$). Required in all forms except array construction \texttt{Matrix(C,...)}.}
  \argitem{n}{integer, optional, intent(in)}{Number of columns ($n > 0$; defaults to \argname{m}). Must equal \argname{m} if \texttt{mtype='ZH'} or \texttt{rtype='1'}.}
  \argitem{is, js}{integer, optional, intent(in)}{Starting row and column indices (\texttt{is} defaults to \texttt{1}; \texttt{js} defaults to \texttt{is}).}
  \argitem{mtype}{character(*), optional, intent(in)}{Symmetry tag: \texttt{'ZG'} (default) or \texttt{'ZH'} (Hermitian; requires $m = n$).}
  \argitem{name}{character(*), optional, intent(in)}{Optional matrix name (up to 20 characters; defaults to \texttt{''}).}
  \argitem{C}{complex(8), dimension(:,:), intent(in)}{Source 2D complex array for \texttt{Matrix(C, ...)}.}
  \argitem{r / c}{real(8) / complex(8), intent(in)}{Constant fill value (passed as the first argument before \argname{m}).}
  \argitem{rtype}{character(*), intent(in)}{Generator type: \texttt{'uniform'}/\texttt{'random'} (\texttt{'u'}/\texttt{'r'}), \texttt{'gaussian'} (\texttt{'g'}), or identity (\texttt{'1'}).}
  \argitem{sigma}{real(8), optional, intent(in)}{Standard deviation for \texttt{'gaussian'} matrices (defaults to \texttt{1.0\_8}).}
  \argitem{MAT}{type(Matrix), return value}{Newly constructed and allocated \typename{Matrix} object.}
\end{description}

\manheading{Example}
\lstset{language=[95]Fortran}
\begin{lstlisting}
use :: matrix_mod_matrix
complex(8)   :: A(0:3, 0:3) = (1.0_8, 2.0_8)
type(Matrix) :: m0, m1, m2, m3, m4

m0 = Matrix(4)                                  ! 4x4 zero matrix, mtype='ZG'
m1 = Matrix(4, 3, js=5, name='Rect')            ! 4x3 zeros, v(1:4, 5:7)
m2 = Matrix(A, is=lbound(A,1), js=lbound(A,2), &
            mtype='ZH', name='HermArr')  ! preserves bounds, Hermitianizes
m3 = Matrix('uniform', 4, mtype='ZH')           ! 4x4 random Hermitian matrix
m4 = Matrix('1', 4)                             ! 4x4 complex identity matrix
\end{lstlisting}

% -------------------------------------------------------------------
\subsection{Component Procedures}
\label{subsec:mat_comp_procs}

All objects of \typename{type(Matrix)} inherit five type-bound procedures from \texttt{MatrixClass} (\procname{save}, \procname{read}, \procname{print}, \procname{random}, \procname{gaussian}) and define eight additional \typename{Matrix}-specific type-bound generics (\procname{hermitian\_set}, \procname{traceless\_set}, \procname{conjg}, \procname{transpose}, \procname{hermitian}, \procname{dagger}, \procname{re}, \procname{im}).

% ===================================================================
% MAT % save
% ===================================================================
\manentry{Type-Bound Subroutine}{MAT \% save}{Matrix \% save}{Writes matrix elements in plain text column-major format to a Fortran unit}{mat:save}

\manheading{Synopsis}
\lstset{language=[95]Fortran}
\begin{lstlisting}
call mat % save([unit], [fmt])
\end{lstlisting}

\manheading{Brief}
Writes all raw complex entries of \texttt{mat \% v} (without headers or metadata) to \argname{unit} in column-major order.

\manheading{Description}
If \texttt{mat \% v} is allocated, \procname{mat \% save} loops in \emph{column-major} order (\texttt{do j = js, je; do i = is, ie}) and writes each \typename{complex(8)} entry \texttt{mat \% v(i,j)} using the format descriptor \texttt{'(2'//fmt//')'}. If \texttt{mat \% v} is not allocated, it returns silently without error. Output matches \procname{mat \% read} in loop order.

\manheading{Dummy Arguments}
\begin{description}[leftmargin=2.2cm, style=nextline]
  \argitem{MAT}{class(Matrix), intent(in)}{Passed-object matrix whose data are saved.}
  \argitem{unit}{integer, optional, intent(in)}{Fortran I/O unit number (defaults to \procname{f\_mout}).}
  \argitem{fmt}{character(*), optional, intent(in)}{Single-real Fortran edit descriptor (defaults to \texttt{'G28.17'}; doubled via \texttt{2} for complex numbers).}
\end{description}

\manheading{Example}
\lstset{language=[95]Fortran}
\begin{lstlisting}
use :: matrix_mod_matrix
type(Matrix) :: m1
integer      :: un

m1 = Matrix('gaussian', 4)
open(newunit=un, file='matrix.dat', status='replace')
call m1 % save(unit=un, fmt='ES24.15E3')
close(un)
\end{lstlisting}

% ===================================================================
% MAT % read
% ===================================================================
\manentry{Type-Bound Subroutine}{MAT \% read}{Matrix \% read}{Reads matrix elements in plain text column-major order from a Fortran unit}{mat:read}

\manheading{Synopsis}
\lstset{language=[95]Fortran}
\begin{lstlisting}
call mat % read([unit])
\end{lstlisting}

\manheading{Brief}
Reads raw complex entries from \argname{unit} into an already-constructed \texttt{mat \% v} in column-major order.

\manheading{Description}
\procname{mat \% read} reads real/imaginary pairs \texttt{(x, y)} using list-directed I/O (\texttt{read(un,*) x, y}) in column-major order (\texttt{do j = js, je; do i = is, ie}) and assigns \texttt{CMPLX(x, y, kind=dp)} to \texttt{mat \% v(i,j)}. The matrix must already be constructed and allocated before calling \procname{\% read}; if \texttt{mat \% v} is unallocated, the call returns silently. Because \texttt{f\_minput} is not publicly exported by \modname{matrix\_mod\_matrix}, omitting \argname{unit} reads from standard input (\texttt{stdin}).

\manheading{Dummy Arguments}
\begin{description}[leftmargin=2.2cm, style=nextline]
  \argitem{MAT}{class(Matrix), intent(inout)}{Pre-constructed matrix whose \texttt{v} array receives the values.}
  \argitem{unit}{integer, optional, intent(in)}{Fortran input unit number (defaults to internal \texttt{f\_minput = stdin}).}
\end{description}

\manheading{Example}
\lstset{language=[95]Fortran}
\begin{lstlisting}
use :: matrix_mod_matrix
type(Matrix) :: m2
integer      :: un

m2 = Matrix(4)
open(newunit=un, file='matrix.dat', status='old')
call m2 % read(unit=un)
close(un)
\end{lstlisting}

% ===================================================================
% MAT % print
% ===================================================================
\manentry{Type-Bound Subroutine}{MAT \% print}{Matrix \% print}{Pretty-prints a matrix, metadata summary, or Mathematica expression}{mat:print}

\manheading{Synopsis}
\lstset{language=[95]Fortran}
\begin{lstlisting}
call mat % print([unit], [fmt], [form], [ips], [ipe], [jps], [jpe])
\end{lstlisting}

\manheading{Brief}
Displays a matrix (or a sub-block \texttt{ips:ipe, jps:jpe}) in human-readable table format (\texttt{'Default'}), metadata-only mode (\texttt{'Info'}), or Wolfram Mathematica syntax (\texttt{'Mathematica'}).

\manheading{Description}
The procedure aborts via \texttt{matrix\_error} when the matrix is not initialized (\texttt{mat \% m == 0}) and, in the \texttt{'Default'} and \texttt{'Mathematica'} forms, when \texttt{mat \% v} is not allocated. It clamps any optional slice bounds \argname{ips, ipe, jps, jpe} inside \texttt{[is,ie]} and \texttt{[js,je]}, and formats output according to \argname{form(1:1)}:
\begin{itemize}
  \item \texttt{'d'} or \texttt{'D'} (\textbf{Default}): Writes a delimited header block (\texttt{\#pmstart: ...}) showing dimensions, \texttt{mtype}, and index bounds, prints rows \texttt{is..ie} using format \texttt{(2*n//fmt)}, and closes with \texttt{\#pmend: ]}.
  \item \texttt{'i'} or \texttt{'I'} (\textbf{Info}): Prints a single \texttt{\#info:} summary line with \texttt{name}, dimensions, \texttt{mtype}, metadata bounds, and actual Fortran array bounds \texttt{lbound/ubound(mat \% v)} (or \texttt{(v not allocated)} when \texttt{mat \% v} is not allocated) without printing elements.
  \item \texttt{'m'} or \texttt{'M'} (\textbf{Mathematica}): Verifies that \texttt{mat \% v} is allocated and rejects unknown subclasses via \texttt{matrix\_error}, then emits a ready-to-paste Mathematica assignment \texttt{<name> = \{\{...\}\};} wrapping every real/imaginary coefficient in \texttt{Interpreter["Number"]["..."]} so Fortran exponential notation (\texttt{E}/\texttt{D}) parses losslessly. If \texttt{mat \% name} is empty, the variable name defaults to \texttt{m}.
\end{itemize}
Passing any other leading character in \argname{form} calls \texttt{matrix\_error}. Unlike vector \procname{\% print}, matrix \procname{\% print} does not take a \argname{name} argument; set \texttt{mat \% name} directly if needed.

\manheading{Dummy Arguments}
\begin{description}[leftmargin=2.2cm, style=nextline]
  \argitem{MAT}{class(Matrix), intent(in)}{Initialized matrix to display.}
  \argitem{unit}{integer, optional, intent(in)}{Destination Fortran unit (defaults to \procname{f\_mout}).}
  \argitem{fmt}{character(*), optional, intent(in)}{Single-real numeric format string (defaults to \texttt{'G28.17'}).}
  \argitem{form}{character(*), optional, intent(in)}{Output mode: \texttt{'Default'} (default), \texttt{'Info'}, or \texttt{'Mathematica'} (case-insensitive first character).}
  \argitem{ips, ipe}{integer, optional, intent(in)}{Optional starting/ending row bounds for printing a sub-slice.}
  \argitem{jps, jpe}{integer, optional, intent(in)}{Optional starting/ending column bounds for printing a sub-slice.}
\end{description}

\manheading{Example}
\lstset{language=[95]Fortran}
\begin{lstlisting}
use :: matrix_mod_matrix
type(Matrix) :: m

m = Matrix('gaussian', 8, name='H', mtype='ZH')
call m % print(form='Info')
call m % print(ips=1, ipe=3, jps=1, jpe=3, fmt='F8.3')
call m % print(form='Mathematica')
\end{lstlisting}

% ===================================================================
% MAT % random
% ===================================================================
\manentry{Type-Bound Subroutine}{MAT \% random}{Matrix \% random}{Fills a matrix with uniform random complex numbers, respecting \texttt{mtype='ZH'}}{mat:random}

\manheading{Synopsis}
\lstset{language=[95]Fortran}
\begin{lstlisting}
call mat % random()
\end{lstlisting}

\manheading{Brief}
Populates an allocated \typename{Matrix} with independent uniform random numbers in $[0,1) + i[0,1)$, preserving Hermiticity when \texttt{mat \% mtype = 'ZH'}.

\manheading{Description}
\procname{mat \% random} fills \texttt{mat \% v} using \texttt{array2\_random\_set}. If \texttt{mat \% mtype = 'ZH'}, instead of averaging $(M + M^\dagger)/2$ (which would produce a triangular distribution $(X_1+X_2)/2$ on off-diagonals), it makes the diagonal real (\texttt{REAL(mat\%v(j,j))}) and copies \texttt{CONJG(mat\%v(j,i))} from the lower triangle into the upper triangle. Consequently, every entry retains a true uniform distribution, and the upper triangle has \emph{negative} imaginary parts in $(-1,0]$. Aborts via \texttt{matrix\_error} if \texttt{mat \% v} is not allocated. Equivalent to \lcodef{call random_number(mat)}.

\manheading{Dummy Arguments}
\begin{description}[leftmargin=2.2cm, style=nextline]
  \argitem{MAT}{class(Matrix), intent(inout)}{Pre-allocated matrix to fill with uniform random entries.}
\end{description}

\manheading{Example}
\lstset{language=[95]Fortran}
\begin{lstlisting}
use :: matrix_mod_matrix
type(Matrix) :: m

call matrix_random_init
m = Matrix(4, mtype='ZH')
call m % random()
\end{lstlisting}

% ===================================================================
% MAT % gaussian
% ===================================================================
\manentry{Type-Bound Subroutine}{MAT \% gaussian}{Matrix \% gaussian}{Fills a matrix with Gaussian random complex numbers, respecting \texttt{mtype='ZH'}}{mat:gaussian}

\manheading{Synopsis}
\lstset{language=[95]Fortran}
\begin{lstlisting}
call mat % gaussian([sigma])
\end{lstlisting}

\manheading{Brief}
Populates an allocated \typename{Matrix} with zero-mean Gaussian random complex entries of standard deviation \argname{sigma}, Hermitianizing when \texttt{mtype = 'ZH'}.

\manheading{Description}
Using Box--Muller Gaussian sampling (\texttt{array2\_gauss\_set}), \procname{mat \% gaussian} fills each real and imaginary part with normal random variates of standard deviation \argname{sigma} (default \texttt{1.0\_8}). When \texttt{mat \% mtype = 'ZH'}, \procname{mat \% gaussian} generates entries with standard deviation $s = \sigma \sqrt{2}$ and invokes \cbox{call mat \% hermitian\_set}, which replaces $X_{ij}$ by $\frac{1}{2}(X_{ij} + X_{ji}^*)$. Off the diagonal this averages two independent variates, so the real and imaginary parts end up with standard deviation $s/\sqrt{2} = \sigma$. On the diagonal the two terms are the same variate, so the result is $\Re X_{ii}$ with standard deviation $s = \sigma\sqrt{2}$; the diagonal is therefore rescaled by $1/\sqrt{2}$. All real parts, imaginary parts, and the (real) diagonal then have standard deviation $\sigma$. Aborts if \texttt{mat \% v} is unallocated.

\manheading{Dummy Arguments}
\begin{description}[leftmargin=2.2cm, style=nextline]
  \argitem{MAT}{class(Matrix), intent(inout)}{Pre-allocated matrix to overwrite with Gaussian random numbers.}
  \argitem{sigma}{real(8), optional, intent(in)}{Target standard deviation $\sigma$ (defaults to \texttt{1.0\_8}).}
\end{description}

\manheading{Example}
\lstset{language=[95]Fortran}
\begin{lstlisting}
use :: matrix_mod_matrix
type(Matrix) :: m

m = Matrix(6, mtype='ZH')
call m % gaussian(sigma=0.5_8)
\end{lstlisting}

% ===================================================================
% MAT % hermitian_set
% ===================================================================
\manentry{Type-Bound Subroutine}{MAT \% hermitian\_set}{Matrix \% hermitian\textunderscore set}{Makes a square matrix Hermitian in-place and sets \texttt{mtype = 'ZH'}}{mat:comp_hermitian_set}

\manheading{Synopsis}
\lstset{language=[95]Fortran}
\begin{lstlisting}
call mat % hermitian_set([uplo])
\end{lstlisting}

\manheading{Brief}
Transforms an allocated square \typename{Matrix} into a Hermitian matrix ($M = M^\dagger$) in-place and updates \texttt{mat \% mtype = 'ZH'}.

\manheading{Description}
\procname{mat \% hermitian\_set} delegates to \procname{hermitian\_set(mat \% v, uplo)}: if \argname{uplo(1:1)} is \texttt{'U'}/\texttt{'u'} (or \texttt{'L'}/\texttt{'l'}), the upper (or lower) triangle is preserved (with diagonal imaginary parts zeroed) and conjugated into the opposite triangle; otherwise (including when \argname{uplo} is omitted), $M \leftarrow \frac{1}{2}(M + M^\dagger)$ is computed. It then sets \texttt{mat \% mtype = 'ZH'}. Unlike the standalone \modname{array\_mod} routine (which writes \texttt{NaN} on non-square inputs), this type-bound method explicitly aborts via \texttt{matrix\_error} if \texttt{mat \% v} is not allocated or if \texttt{mat \% m /= mat \% n}.

\manheading{Dummy Arguments}
\begin{description}[leftmargin=2.2cm, style=nextline]
  \argitem{MAT}{class(Matrix), intent(inout)}{Allocated square matrix to Hermitianize in-place.}
  \argitem{uplo}{character(*), optional, intent(in)}{Triangle selector: \texttt{'U'} preserves upper triangle, \texttt{'L'} preserves lower triangle, omitted/other averages $(M+M^\dagger)/2$.}
\end{description}

\manheading{Example}
\lstset{language=[95]Fortran}
\begin{lstlisting}
use :: matrix_mod_matrix
type(Matrix) :: m

m = Matrix('uniform', 4)
call m % hermitian_set()          ! average (M + M^dagger)/2, mtype='ZH'
call m % hermitian_set(uplo='U')  ! mirror upper triangle into lower
\end{lstlisting}

% ===================================================================
% MAT % traceless_set
% ===================================================================
\manentry{Type-Bound Pure Subroutine}{MAT \% traceless\_set}{Matrix \% traceless\textunderscore set}{Projects a square matrix onto its traceless part in-place}{mat:comp_traceless_set}

\manheading{Synopsis}
\lstset{language=[95]Fortran}
\begin{lstlisting}
call mat % traceless_set()
\end{lstlisting}

\manheading{Brief}
Subtracts $\frac{1}{N}(\Tr M)\, I_N$ from the diagonal of \texttt{mat \% v} in-place so that $\Tr M = 0$.

\manheading{Description}
Calling \cbox{call mat \% traceless\_set()} executes \procname{traceless\_set(mat \% v)}, replacing $M \leftarrow M - \frac{\Tr M}{N} I_N$ for an $N\times N$ matrix while preserving \texttt{mat \% mtype} and \texttt{mat \% name}. Note that if $M$ is Hermitian, subtracting a scalar multiple of the identity keeps $M$ Hermitian iff $\Im(\Tr M)=0$ (always true up to rounding error), so \texttt{mtype = 'ZH'} remains valid. Non-square \texttt{mat \% v} causes the underlying \texttt{traceless\_set} to fill \texttt{mat \% v} with \texttt{NaN}.

\manheading{Dummy Arguments}
\begin{description}[leftmargin=2.2cm, style=nextline]
  \argitem{MAT}{class(Matrix), intent(inout)}{Square matrix modified in-place to have zero trace.}
\end{description}

\manheading{Example}
\lstset{language=[95]Fortran}
\begin{lstlisting}
use :: matrix_mod_matrix
type(Matrix) :: m

m = Matrix('gaussian', 6, mtype='ZH')
call m % traceless_set()
print *, 'Trace after projection:', trace(m)
\end{lstlisting}

% ===================================================================
% MAT % conjg
% ===================================================================
\manentry{Type-Bound Function}{MAT \% conjg}{Matrix \% conjg}{Returns the element-wise complex conjugate $M^*$ as a new \texttt{Matrix}}{mat:comp_conjg}

\manheading{Synopsis}
\lstset{language=[95]Fortran}
\begin{lstlisting}
m2 = m1 % conjg()
\end{lstlisting}

\manheading{Brief}
Allocates and returns a \typename{type(Matrix)} containing the complex conjugate $\texttt{CONJG(m1 \% v)}$ with all metadata preserved.

\manheading{Description}
\procname{m1 \% conjg()} copies all metadata (\texttt{m, n, is, ie, js, je, mtype, name}) from \argname{m1} to the result \argname{m2}, allocates \texttt{m2 \% v(is:ie, js:je)}, and sets \texttt{m2 \% v = CONJG(m1 \% v)}. Because the conjugate of a Hermitian matrix is also Hermitian (and equal to $M^T$), preserving \texttt{mtype} is exact. Equivalent to \lcodef{conjg(m1)}.

\manheading{Dummy Arguments}
\begin{description}[leftmargin=2.2cm, style=nextline]
  \argitem{MATA}{class(Matrix), intent(in)}{Passed-object input matrix.}
  \argitem{MATB}{type(Matrix), return value}{Complex conjugate matrix with identical shape, bounds, \texttt{mtype}, and \texttt{name}.}
\end{description}

\manheading{Example}
\lstset{language=[95]Fortran}
\begin{lstlisting}
use :: matrix_mod_matrix
type(Matrix) :: m1, m2

m1 = Matrix('gaussian', 3, name='M1')
m2 = m1 % conjg()
\end{lstlisting}

% ===================================================================
% MAT % transpose
% ===================================================================
\manentry{Type-Bound Function}{MAT \% transpose}{Matrix \% transpose}{Returns the transpose $M^T$ as a new \texttt{Matrix}, swapping $m\leftrightarrow n$ and index bounds}{mat:comp_transpose}

\manheading{Synopsis}
\lstset{language=[95]Fortran}
\begin{lstlisting}
m2 = m1 % transpose()
\end{lstlisting}

\manheading{Brief}
Returns an $n\times m$ \typename{type(Matrix)} equal to \texttt{TRANSPOSE(m1 \% v)} with row and column bounds swapped.

\manheading{Description}
For an $m \times n$ matrix with bounds \texttt{is:ie, js:je}, \procname{m1 \% transpose()} constructs an $n \times m$ result \argname{m2} with swapped dimensions and bounds: \texttt{m2 \% m = m1 \% n}, \texttt{m2 \% n = m1 \% m}, \texttt{m2 \% is = m1 \% js}, \texttt{m2 \% js = m1 \% is} (so \texttt{m2 \% v} is allocated over \texttt{m1\%js:m1\%je, m1\%is:m1\%ie}), while preserving \texttt{mtype} and \texttt{name}. Equivalent to \lcodef{transpose(m1)}.

\manheading{Dummy Arguments}
\begin{description}[leftmargin=2.2cm, style=nextline]
  \argitem{MATA}{class(Matrix), intent(in)}{Passed-object $m \times n$ input matrix.}
  \argitem{MATB}{type(Matrix), return value}{Transposed $n \times m$ matrix with swapped index ranges.}
\end{description}

\manheading{Example}
\lstset{language=[95]Fortran}
\begin{lstlisting}
use :: matrix_mod_matrix
type(Matrix) :: m1, m2

m1 = Matrix('uniform', 3, 5, is=0, js=10, name='Rect')
m2 = m1 % transpose()   ! 5x3 matrix with is=10, js=0
\end{lstlisting}

% ===================================================================
% MAT % hermitian and MAT % dagger
% ===================================================================
\manentry{Type-Bound Functions}{MAT \% hermitian, MAT \% dagger}{Matrix \% hermitian, Matrix \% dagger}{Return the Hermitian conjugate $M^\dagger$ as a new \texttt{Matrix}}{mat:comp_hermitian}
\addcontentsline{toc}{subsubsection}{\texorpdfstring{\texttt{\textbf{Matrix \% dagger}} \textnormal{\footnotesize--- Alias for \texttt{Matrix \% hermitian}}}{Matrix \% dagger}}

\manheading{Synopsis}
\lstset{language=[95]Fortran}
\begin{lstlisting}
m2 = m1 % hermitian()
m2 = m1 % dagger()     ! identical alias
\end{lstlisting}

\manheading{Brief}
Return the conjugate transpose $M^\dagger = \texttt{TRANSPOSE(CONJG(m1 \% v))}$ as a new \typename{type(Matrix)}, swapping dimensions and row/column index bounds.

\manheading{Description}
Both \procname{\% hermitian()} and its alias \procname{\% dagger()} map to \texttt{matrix\_return\_hermitian}. They construct an $n \times m$ matrix \argname{m2} with \texttt{m2 \% m = m1 \% n}, \texttt{m2 \% n = m1 \% m}, \texttt{m2 \% is = m1 \% js}, \texttt{m2 \% js = m1 \% is}, preserving \texttt{mtype} and \texttt{name}, and compute \texttt{m2 \% v = TRANSPOSE(CONJG(m1 \% v))}. Equivalent to \lcodef{hermitian(m1)}.

\manheading{Dummy Arguments}
\begin{description}[leftmargin=2.2cm, style=nextline]
  \argitem{MATA}{class(Matrix), intent(in)}{Passed-object $m \times n$ input matrix.}
  \argitem{MATB}{type(Matrix), return value}{Hermitian conjugate $n \times m$ matrix $M^\dagger$.}
\end{description}

\manheading{Example}
\lstset{language=[95]Fortran}
\begin{lstlisting}
use :: matrix_mod_matrix
type(Matrix) :: m1, m2, h

m1 = Matrix('gaussian', 3, 5)
m2 = m1 % dagger()     ! 5x3 conjugate transpose
h  = m1 * m1 % dagger()
\end{lstlisting}

% ===================================================================
% MAT % re and MAT % im
% ===================================================================
\manentry{Type-Bound Functions}{MAT \% re, MAT \% im}{Matrix \% re, Matrix \% im}{Extract the real or imaginary part of a \texttt{Matrix} as a \texttt{DMatrix}}{mat:comp_re_im}
\addcontentsline{toc}{subsubsection}{\texorpdfstring{\texttt{\textbf{Matrix \% im}} \textnormal{\footnotesize--- Imaginary part of a \texttt{Matrix} as a \texttt{DMatrix}}}{Matrix \% im}}

\manheading{Synopsis}
\lstset{language=[95]Fortran}
\begin{lstlisting}
d_re = m1 % re()
d_im = m1 % im()
\end{lstlisting}

\manheading{Brief}
Return a \typename{type(DMatrix)} containing the real part (\procname{\% re()}) or imaginary part (\procname{\% im()}) of \argname{m1}, preserving shape, index bounds, and \texttt{name}.

\manheading{Description}
Both methods copy metadata \texttt{(m, n, is, ie, js, je, name)} from \argname{m1} into a new \typename{type(DMatrix)} \argname{d}:
\begin{itemize}
  \item \procname{m1 \% re()}: Sets \texttt{d \% v = REAL(m1 \% v, kind=8)}. If \texttt{m1 \% mtype == 'ZH'}, the real part of a Hermitian matrix is real symmetric, so it sets \texttt{d \% mtype = 'DS'}; otherwise \texttt{d \% mtype = 'DG'}.
  \item \procname{m1 \% im()}: Sets \texttt{d \% v = AIMAG(m1 \% v)} and always sets \texttt{d \% mtype = 'DG'} (since the imaginary part of a Hermitian matrix is antisymmetric rather than symmetric).
\end{itemize}
Equivalent to the standalone functions \lcodef{real(m1)} and \lcodef{aimag(m1)}.

\manheading{Dummy Arguments}
\begin{description}[leftmargin=2.2cm, style=nextline]
  \argitem{MATA}{class(Matrix), intent(in)}{Passed-object complex matrix.}
  \argitem{MATB}{type(DMatrix), return value}{Real double-precision matrix holding $\Re(\argname{m1})$ or $\Im(\argname{m1})$.}
\end{description}

\manheading{Example}
\lstset{language=[95]Fortran}
\begin{lstlisting}
use :: matrix_mod_matrix
type(Matrix)  :: m1
type(DMatrix) :: dre, dim

m1  = Matrix('gaussian', 4, mtype='ZH')
dre = m1 % re()   ! DMatrix with dre % mtype == 'DS'
dim = m1 % im()   ! DMatrix with dim % mtype == 'DG'
\end{lstlisting}

% -------------------------------------------------------------------
\subsection{Operators}
\label{subsec:mat_operators}

\modname{matrix\_mod\_matrix} overloads Fortran's assignment (\procname{=}), unary/binary arithmetic (\procname{+}, \procname{-}, \procname{*}, \procname{/}, \procname{**}), and BLAS-backed matrix multiplication (\procname{*}) operators for \typename{Matrix}.

\mansubgroup{Assignment}

% ===================================================================
% assignment(=) Matrix
% ===================================================================
\manentry{Overloaded Operator}{assignment(=) (Matrix)}{Matrix: assignment(=)}{Assigns another matrix, array, or scalar to a \texttt{Matrix}}{mat:op_assign}

\manheading{Synopsis}
\lstset{language=[95]Fortran}
\begin{lstlisting}
mat = m1     ! type(Matrix)           -> copies metadata and values
mat = d1     ! type(DMatrix)          -> 'ZH' if 'DS' else 'ZG', keeps name
mat = C      ! complex(8) :: C(:,:)   -> copies array C into mat % v
mat = r      ! real(8)    :: r        -> fills pre-allocated mat % v = r
mat = z      ! complex(8) :: z        -> fills pre-allocated mat % v = z
\end{lstlisting}

\manheading{Brief}
Overloads \texttt{=} for assigning a \typename{Matrix}, \typename{DMatrix}, 2D \typename{complex(8)} array, \typename{real(8)} scalar, or \typename{complex(8)} scalar to a \typename{Matrix}.

\manheading{Description}
Depending on the right-hand side type, \texttt{mat = rhs} behaves as follows:
\begin{itemize}
  \item \texttt{mat = m1} (\typename{type(Matrix)}): Intrinsic Fortran~2003 derived-type assignment; copies all metadata (\texttt{m, n, is, ie, js, je, mtype, name}) and automatically allocates or reallocates \texttt{mat \% v} to match \texttt{m1 \% v}, whether or not \argname{mat} was previously constructed.
  \item \texttt{mat = d1} (\typename{type(DMatrix)}): Initializes \argname{mat} with \texttt{(d1 \% m, d1 \% n, d1 \% is, d1 \% js, mt, name = d1 \% name)}, setting \texttt{mt = 'ZH'} when \texttt{d1 \% mtype = 'DS'} (otherwise \texttt{'ZG'}), preserving \texttt{d1 \% name}, and copying \texttt{CMPLX(d1 \% v, 0.0\_8, kind=8)}.
  \item \texttt{mat = C} (\typename{complex(8)(:,:)}): If \argname{mat} has not been constructed (\texttt{mat \% m == 0}), initializes \texttt{m, n} from \texttt{shape(C)}, sets bounds to \texttt{lbound(C)..ubound(C)} (which for assumed-shape descriptors is \texttt{1..size}), sets \texttt{mtype = 'ZG'}, \texttt{name = ''}, and copies \argname{C}. If \argname{mat} was already constructed (\texttt{mat \% m /= 0}), checks that \texttt{shape(C) == [mat\%m, mat\%n]} (aborting via \texttt{matrix\_error} if non-conformable), keeps existing \texttt{is, js, name} unchanged, resets \texttt{mat \% mtype = 'ZG'}, and copies \argname{C} into \texttt{mat \% v}.
  \item \texttt{mat = r} (\typename{real(8)}) and \texttt{mat = z} (\typename{complex(8)}): Require \texttt{mat \% v} to be already allocated (aborting via \texttt{matrix\_error} otherwise) and broadcast the scalar to all elements of \texttt{mat \% v}. For complex \argname{z}, \texttt{mat \% mtype} is reset to \texttt{'ZG'} only when \texttt{AIMAG(z) /= 0.0\_8}; otherwise \texttt{mtype} is preserved.
\end{itemize}

\begin{warnbox}[Metadata Reset on Array Assignment and Missing \texttt{real(8)(:,:)}]
\begin{enumerate}
  \item Assigning a 2D array \lcodef{mat = C} to an already-constructed \typename{Matrix} resets \texttt{mat \% mtype = 'ZG'} (since an arbitrary array is not Hermitian in general); call \cbox{call mat \% hermitian\_set} afterwards if \argname{C} is Hermitian.
  \item Direct assignment from a 2D \typename{real(8)} array (\texttt{mat = A}) is \emph{not} overloaded; wrap it in \lcodef{DMatrix(A)} or assign \cbox{mat \% v = A} after construction.
\end{enumerate}
\end{warnbox}

\manheading{Dummy Arguments}
\begin{description}[leftmargin=2.2cm, style=nextline]
  \argitem{MATB}{type(Matrix), intent(inout)}{Destination matrix on the left-hand side.}
  \argitem{rhs}{type(Matrix), type(DMatrix), complex(8)(:,:), real(8), or complex(8), intent(in)}{Source value or matrix on the right-hand side.}
\end{description}

\manheading{Example}
\lstset{language=[95]Fortran}
\begin{lstlisting}
use :: matrix_mod_matrix
type(Matrix)  :: m1, m2
type(DMatrix) :: d1
complex(8)    :: C(4,4)

m1 = Matrix(4, is=2, name='Target')
C  = (1.0_8, 0.5_8)
m1 = C                    ! keeps m1 % is=2 and m1 % name='Target'
! m2 = (0.0_8, 0.0_8)     ! would abort: m2 is not constructed yet
m2 = Matrix(4)            ! construct m2 first, then scalar assignment works
m2 = 2.5_8                ! sets every element of m2 % v to (2.5, 0.0)
d1 = DMatrix('gaussian', 4)
m2 = d1                   ! converts DMatrix to Matrix with mtype='ZG'
\end{lstlisting}

\mansubgroup{Addition, Subtraction, Scalar Multiplication, Division, and Power}

% ===================================================================
% operator(+) Matrix
% ===================================================================
\manentry{Overloaded Operator}{operator(+) (Matrix)}{Matrix: operator(+)}{Adds a \texttt{Matrix} with a scalar, complex 2D array, \texttt{Matrix}, or \texttt{DMatrix}}{mat:op_plus}

\manheading{Synopsis}
\lstset{language=[95]Fortran}
\begin{lstlisting}
m2 = m1 + r;   m2 = r + m1     ! real(8) scalar + Matrix (pure)
m2 = m1 + z;   m2 = z + m1     ! complex(8) scalar + Matrix (pure)
m2 = m1 + C;   m2 = C + m1     ! complex(8)(:,:) array + Matrix
m3 = m1 + m2                   ! Matrix + Matrix
m3 = m1 + d2;  m3 = d2 + m1    ! mixed Matrix + DMatrix -> Matrix
\end{lstlisting}

\manheading{Brief}
Performs element-wise addition (or scalar broadcast addition) returning a new \typename{type(Matrix)} with propagated metadata.

\manheading{Description}
All nine overloads construct the result by copying the \typename{Matrix} operand's bounds and \texttt{name} (for \texttt{d2 + m1}, metadata are inherited from \argname{m1}), then compute the element-wise sum. Matrix/array operands must have identical shape $(m, n)$, otherwise execution aborts via \texttt{matrix\_error}. The resulting \texttt{mtype} follows these rules:
\begin{itemize}
  \item \texttt{m1 + m2}: \texttt{'ZH'} if both \argname{m1} and \argname{m2} have \texttt{mtype = 'ZH'}, otherwise \texttt{'ZG'}.
  \item \texttt{m1 + d2} and \texttt{d2 + m1}: \texttt{'ZH'} if the \typename{Matrix} is \texttt{'ZH'} and the \typename{DMatrix} is \texttt{'DS'}, otherwise \texttt{'ZG'}.
  \item \texttt{m1 + C} and \texttt{C + m1}: Always sets \texttt{mtype = 'ZG'}.
  \item \texttt{m1 + r}, \texttt{r + m1}: Preserve \texttt{m1 \% mtype}. For complex scalar addition \texttt{m1 + z} and \texttt{z + m1}, \texttt{'ZH'} is preserved only when \texttt{AIMAG(z) == 0.0\_8}; adding a complex scalar with nonzero imaginary part resets \texttt{mtype = 'ZG'}.
\end{itemize}

\manheading{Dummy Arguments}
\begin{description}[leftmargin=2.2cm, style=nextline]
  \argitem{lhs, rhs}{Matrix, DMatrix, complex(8)(:,:), real(8), or complex(8), intent(in)}{Operands (at least one must be \typename{Matrix}).}
  \argitem{result}{type(Matrix), return value}{Allocated sum matrix.}
\end{description}

\manheading{Example}
\lstset{language=[95]Fortran}
\begin{lstlisting}
use :: matrix_mod_matrix
type(Matrix)  :: m1, m2, m3
type(DMatrix) :: d1

m1 = Matrix('gaussian', 4, mtype='ZH')
d1 = DMatrix('gaussian', 4, mtype='DS')
m2 = m1 + 2.0_8
m3 = m1 + d1      ! Hermitian ('ZH') because both m1 ('ZH') and d1 ('DS')
\end{lstlisting}

% ===================================================================
% operator(-) Matrix
% ===================================================================
\manentry{Overloaded Operator}{operator(-) (Matrix)}{Matrix: operator(-)}{Unary negation and binary subtraction involving a \texttt{Matrix}}{mat:op_minus}

\manheading{Synopsis}
\lstset{language=[95]Fortran}
\begin{lstlisting}
m2 = -m1                       ! unary negation
m2 = m1 - r;   m2 = r - m1     ! real(8) scalar and Matrix (pure)
m2 = m1 - z;   m2 = z - m1     ! complex(8) scalar and Matrix (pure)
m2 = m1 - C;   m2 = C - m1     ! complex(8)(:,:) array and Matrix
m3 = m1 - m2                   ! Matrix - Matrix
m3 = m1 - d2;  m3 = d2 - m1    ! mixed Matrix and DMatrix -> Matrix
\end{lstlisting}

\manheading{Brief}
Computes unary element-wise negation \texttt{-m1} or binary element-wise difference, returning a new \typename{type(Matrix)}.

\manheading{Description}
Unary \texttt{-m1} preserves all metadata (including \texttt{'ZH'}) and negates \texttt{v}. Binary \texttt{m1 - m2}, \texttt{m1 - d2}, and \texttt{d2 - m1} enforce shape conformability (calling \texttt{matrix\_error} on mismatch) and preserve \texttt{'ZH'} only when both matrices are Hermitian/symmetric (\texttt{'ZH'}/\texttt{'DS'}); otherwise they set \texttt{mtype = 'ZG'}. Real scalar differences \texttt{m1 - r} and \texttt{r - m1} preserve \texttt{m1 \% mtype}; complex scalar differences \texttt{m1 - z} and \texttt{z - m1} preserve \texttt{'ZH'} iff \texttt{AIMAG(z) == 0.0\_8} and reset \texttt{mtype = 'ZG'} otherwise; and array differences \texttt{m1 - C} and \texttt{C - m1} always reset \texttt{mtype = 'ZG'}.

\manheading{Dummy Arguments}
\begin{description}[leftmargin=2.2cm, style=nextline]
  \argitem{lhs, rhs}{Matrix, DMatrix, complex(8)(:,:), real(8), or complex(8), intent(in)}{Operands for unary or binary subtraction.}
  \argitem{result}{type(Matrix), return value}{Allocated difference matrix.}
\end{description}

\manheading{Example}
\lstset{language=[95]Fortran}
\begin{lstlisting}
use :: matrix_mod_matrix
type(Matrix) :: m1, m2, m3

m1 = Matrix('gaussian', 4, mtype='ZH')
m2 = Matrix('gaussian', 4, mtype='ZH')
m3 = -m1 - m2     ! result retains mtype='ZH'
\end{lstlisting}

% ===================================================================
% Scalar * and / Matrix
% ===================================================================
\manentry{Overloaded Operators}{operator(*) (Scalar), operator(/) (Matrix)}{Matrix: scalar * and /}{Scalar multiplication and division for \texttt{Matrix} and complex-scaled \texttt{DMatrix}}{mat:op_scale}

\manheading{Synopsis}
\lstset{language=[95]Fortran}
\begin{lstlisting}
m2 = r * m1;  m2 = m1 * r;  m2 = m1 / r    ! real(8) scaling (pure)
m2 = z * m1;  m2 = m1 * z;  m2 = m1 / z    ! complex(8) scaling (pure)
m2 = z * d1;  m2 = d1 * z                  ! z * DMatrix -> Matrix (pure)
\end{lstlisting}

\manheading{Brief}
Multiplies or divides every element of a \typename{Matrix} by a real or complex scalar, or promotes a \typename{DMatrix} to a \typename{Matrix} via complex scalar multiplication.

\manheading{Description}
All eight overloads are \texttt{pure} functions:
\begin{itemize}
  \item \texttt{r * m1}, \texttt{m1 * r}, \texttt{m1 / r}: Copy all metadata from \argname{m1} (preserving \texttt{mtype = 'ZH'}, since a real multiple of a Hermitian matrix stays Hermitian) and scale \texttt{m1 \% v}.
  \item \texttt{z * m1}, \texttt{m1 * z}, \texttt{m1 / z}: Copy metadata from \argname{m1}, scale \texttt{m1 \% v}, and reset \texttt{m2 \% mtype = 'ZG'} whenever \texttt{AIMAG(z) /= 0.0\_8}.
  \item \texttt{z * d1}, \texttt{d1 * z}: Convert a \typename{DMatrix} \argname{d1} into a \typename{type(Matrix)} \argname{m2} copying dimensions and index bounds from \argname{d1}, setting \texttt{m2 \% mtype = 'ZG'} and \texttt{m2 \% name = d1 \% name}, with fast special-cased paths for $z = i, 1, x, iy$ to avoid unnecessary complex multiplications. This provides an idiomatic way to build \texttt{cmplx(d1, 0)} via \lcodef{(1.0_8, 0.0_8) * d1} or \texttt{cmplx(0, d1)} via \lcodef{IMU * d1}.
\end{itemize}

\manheading{Dummy Arguments}
\begin{description}[leftmargin=2.2cm, style=nextline]
  \argitem{r / z}{real(8) or complex(8), scalar, intent(in)}{Real or complex scaling factor.}
  \argitem{MATA}{class(Matrix) or class(DMatrix), intent(in)}{Matrix to scale.}
  \argitem{MATB}{type(Matrix), return value}{Scaled \typename{Matrix} result.}
\end{description}

\manheading{Example}
\lstset{language=[95]Fortran}
\begin{lstlisting}
use :: matrix_mod_matrix
type(Matrix)  :: m1, m2
type(DMatrix) :: d1

m1 = Matrix('gaussian', 4, mtype='ZH')
m2 = 2.0_8 * m1 / 3.0_8            ! preserves mtype='ZH'
m2 = (0.0_8, 1.0_8) * m1           ! AIMAG /= 0 -> mtype becomes 'ZG'
d1 = DMatrix('gaussian', 4)
m2 = (1.0_8, 0.0_8) * d1           ! promote DMatrix to complex Matrix
\end{lstlisting}

% ===================================================================
% operator(**) Matrix
% ===================================================================
\manentry{Overloaded Operator}{operator(**) (Matrix)}{Matrix: operator(**)}{Computes element-wise exponentiation $(M_{ij})^x$ of a \texttt{Matrix}}{mat:op_pow}

\manheading{Synopsis}
\lstset{language=[95]Fortran}
\begin{lstlisting}
m2 = m1 ** k    ! integer k (pure)
m2 = m1 ** x    ! real(8) x (pure)
m2 = m1 ** z    ! complex(8) z (pure)
\end{lstlisting}

\manheading{Brief}
Raises each individual entry \texttt{m1 \% v(i,j)} to an integer, real, or complex power, returning a \typename{Matrix} with \texttt{mtype = 'ZG'}.

\manheading{Description}
This \texttt{pure} operator copies dimensions, bounds, and \texttt{name} from \argname{m1}, sets \texttt{m2 \% mtype = 'ZG'}, and evaluates \texttt{(m1 \% v) ** n} using Fortran's element-wise power intrinsic.

\begin{warnbox}[Element-Wise Exponentiation, Not Matrix Power]
\lcodef{m1 ** 2} computes the Hadamard (entry-wise) square $(M_{ij})^2$, \emph{not} the matrix product $M \cdot M$. Use \lcodef{m1 * m1} or \lcodef{call mmmult(m1, m1, m2)} for matrix multiplication.
\end{warnbox}

\manheading{Dummy Arguments}
\begin{description}[leftmargin=2.2cm, style=nextline]
  \argitem{MAT}{class(Matrix), intent(in)}{Base matrix.}
  \argitem{n}{integer, real(8), or complex(8), scalar, intent(in)}{Exponent applied element-wise.}
  \argitem{MATB}{type(Matrix), return value}{Allocated \typename{Matrix} with \texttt{mtype = 'ZG'}.}
\end{description}

\manheading{Example}
\lstset{language=[95]Fortran}
\begin{lstlisting}
use :: matrix_mod_matrix
type(Matrix) :: m1, m2

m1 = Matrix('uniform', 4)
m2 = m1 ** 2       ! m2 % v(i,j) == m1 % v(i,j)**2
\end{lstlisting}

\mansubgroup{Matrix Multiplication}

% ===================================================================
% operator(*) Matrix Multiplication
% ===================================================================
\manentry{Overloaded Operator}{operator(*) (Matrix Multiplication)}{Matrix: matrix *}{BLAS-backed matrix--matrix and matrix--vector multiplication for \texttt{Matrix}}{mat:op_matmul}

\manheading{Synopsis}
\lstset{language=[95]Fortran}
\begin{lstlisting}
m3 = m1 * m2;   m3 = m1 * C;   m3 = C * m1   ! Matrix-Matrix / array (pure)
m3 = m1 * d2;   m3 = d2 * m1                 ! Mixed Matrix * DMatrix (pure)
v2 = m1 * v1;   v2 = v1 * m1                 ! Matrix-Vector -> Vector
\end{lstlisting}

\manheading{Brief}
Computes matrix--matrix products ($A B$) and matrix--vector products ($M v$ or $v^T M$) using \modname{matrix\_mod\_array}'s \texttt{BLAS} wrappers (\texttt{LMATMUL}).

\manheading{Description}
Overloaded multiplication supports matrix--matrix products between \typename{Matrix}, \typename{DMatrix}, and 2D \typename{complex(8)} arrays, as well as matrix--vector products between \typename{Matrix} and \typename{Vector} (also documented in Section~\ref{sec:class_vector}):
\begin{itemize}
  \item \textbf{Matrix--Matrix (\texttt{m1*m2}, \texttt{m1*C}, \texttt{C*m1}, \texttt{m1*d2}, \texttt{d2*m1})}: Any prior metadata or allocation of the LHS target is replaced. Row bounds \texttt{is:ie} come from the left factor (\texttt{lbound(C,1):ubound(C,1)} for array \argname{C}, i.e.\ \texttt{1:size(C,1)}); column bounds \texttt{js:je} come from the right factor (\texttt{1:size(C,2)} for \argname{C}). The result always gets \texttt{mtype = 'ZG'}, and \texttt{name} is copied from the \typename{Matrix} operand (\argname{m1} in all five forms, also in \texttt{C * m1} and \texttt{d2 * m1}). When the left factor is a \typename{Matrix} with \texttt{mtype = 'ZH'}, \texttt{LMATMUL} uses \texttt{BLAS} \texttt{zhemm}; otherwise \texttt{zgemm} is called. If the inner dimensions differ (\texttt{m1 \% n /= m2 \% m}), the program stops with an \texttt{error stop} message naming the procedure (exit status 1).
  \item \textbf{Matrix--Vector (\texttt{m1*v1}, \texttt{v1*m1})}: \texttt{m1 * v1} returns a \typename{Vector} with \texttt{v2 \% v(m1\%is:m1\%ie)} and \texttt{v2 \% name = v1 \% name} (using \texttt{zhemv} if \texttt{m1 \% mtype = 'ZH'}, else \texttt{zgemv}). Left multiplication \texttt{v1 * m1} computes $M_1^T v_1$ with bounds \texttt{m1\%js:m1\%je} (calling \texttt{LMATMUL(CONJG(m1\%v), v1\%v, type='H')} when \texttt{'ZH'}, else \texttt{zgemv('T',...)}; note that for Hermitian $M_1$, $M_1^T = M_1^*$ so this computes $M_1^T v_1 = (v_1^T M_1)^T$ identically). On a dimension mismatch (\texttt{m1\%n /= v1\%n} or \texttt{m1\%m /= v1\%n}), the program aborts via \texttt{matrix\_error}.
\end{itemize}

\manheading{Dummy Arguments}
\begin{description}[leftmargin=2.2cm, style=nextline]
  \argitem{lhs, rhs}{Matrix, DMatrix, Vector, or complex(8)(:,:), intent(in)}{Conformable factors to multiply.}
  \argitem{result}{type(Matrix) or type(Vector), return value}{Allocated matrix or vector product.}
\end{description}

\manheading{Example}
\lstset{language=[95]Fortran}
\begin{lstlisting}
use :: matrix_mod_matrix
type(Matrix) :: m1, m2, m3
type(Vector) :: v1, v2

m1 = Matrix('gaussian', 4, is=1, name='M1')
m2 = Matrix('gaussian', 4, js=10)
m3 = m1 * m2       ! m3 % v(1:4, 10:13), m3 % name = 'M1', m3 % mtype = 'ZG'
v1 = Vector('gaussian', 4, name='V1')
v2 = m1 * v1       ! v2 % v(1:4), v2 % name = 'V1'
\end{lstlisting}

% ===================================================================
% mmmult (Matrix)
% ===================================================================
\manentry{Pure Generic Subroutine}{mmmult (Matrix)}{Matrix: mmmult}{High-performance zero-allocation BLAS subroutine for \texttt{Matrix} multiplication}{mat:mmmult}

\manheading{Synopsis}
\lstset{language=[95]Fortran}
\begin{lstlisting}
call mmmult(m1, m2, m3)   ! computes m3 % v = m1 % v . m2 % v
\end{lstlisting}

\manheading{Brief}
Computes the matrix product $M_3 = M_1 M_2$ into \argname{m3} (auto-allocating/resizing \argname{m3} if needed, or in-place with zero allocations when \argname{m3} already has matching dimensions).

\manheading{Description}
If \texttt{m3 \% v} is unallocated or its output dimensions do not match (\texttt{m3 \% m /= m1 \% m .or. m3 \% n /= m2 \% n}), \procname{mmmult} automatically allocates or reconstructs \argname{m3} with bounds \texttt{m1 \% is : m1 \% ie, m2 \% js : m2 \% je}; if \argname{m3} is already allocated with matching \texttt{(m1\%m, m2\%n)} shape, its existing bounds and \texttt{name} are preserved in-place with zero allocations (important for small $N \le 500$ loops vs.\ \lcodef{m3 = m1 * m2}). It then invokes \modname{matrix\_mod\_array}'s \texttt{mmmult(m1 \% v, m2 \% v, m3 \% v, mtype = m1 \% mtype)} (selecting \texttt{BLAS} \texttt{zhemm} for \texttt{m1 \% mtype = 'ZH'} and \texttt{zgemm} otherwise, after checking that \texttt{m1 \% n == m2 \% m}; otherwise the program stops with an \texttt{error stop} message naming the procedure (exit status 1)) and sets \texttt{m3 \% mtype = 'ZG'} (since $M_1 M_2$ is not Hermitian in general). Note that aliasing an input argument as the output (e.g.\ \lcodef{call mmmult(m1, m2, m1)}) is invalid because \argname{MATA}/\argname{MATB} are \texttt{intent(in)} and \argname{MATC} is \texttt{intent(inout)}.

\manheading{Dummy Arguments}
\begin{description}[leftmargin=2.2cm, style=nextline]
  \argitem{MATA}{class(Matrix), intent(in)}{Left $m \times k$ factor; its \texttt{mtype} selects \texttt{zhemm} (\texttt{'ZH'}) or \texttt{zgemm}.}
  \argitem{MATB}{class(Matrix), intent(in)}{Right $k \times n$ factor.}
  \argitem{MATC}{type(Matrix), intent(inout)}{Destination $m \times n$ matrix receiving the product (auto-allocated/resized if needed; sets \texttt{mtype = 'ZG'}).}
\end{description}

\manheading{Example}
\lstset{language=[95]Fortran}
\begin{lstlisting}
use :: matrix_mod_matrix
type(Matrix) :: m1, m2, m3

m1 = Matrix('gaussian', 128, mtype='ZH')
m2 = Matrix('gaussian', 128)
m3 = Matrix(128)
call mmmult(m1, m2, m3)   ! fast in-place BLAS zhemm into m3 % v
\end{lstlisting}

% ===================================================================
% Section 2 (cont.): Class Matrix - Procedures
% ===================================================================
\subsection{Procedures}
\label{subsec:mat_procs}

This subsection documents all standalone generic and module procedures operating on \typename{type(Matrix)} arguments.

\mansubgroup{Random Numbers}

% ===================================================================
% random_number (Matrix)
% ===================================================================
\manentry{Generic Subroutine}{random\_number (Matrix)}{Matrix: random\textunderscore number}{Overloaded uniform and Gaussian random matrix generator for \texttt{Matrix}}{mat:random_number}

\manheading{Synopsis}
\lstset{language=[95]Fortran}
\begin{lstlisting}
call random_number(mat)          ! uniform random numbers in [0,1) + i[0,1)
call random_number(mat, sigma)   ! Gaussian random numbers with std dev sigma
\end{lstlisting}

\manheading{Brief}
Fills a pre-allocated \typename{type(Matrix)} with uniformly or normally distributed random complex numbers via Fortran's overloaded \procname{random\_number} intrinsic.

\manheading{Description}
\modname{matrix\_mod\_matrix} extends Fortran's \procname{random\_number} with two overloads for \typename{type(Matrix)}:
\begin{itemize}
  \item \lcodef{call random_number(mat)}: Calls \procname{mat \% random()}, filling \texttt{mat \% v} with uniform $[0,1) + i[0,1)$ values while preserving Hermiticity if \texttt{mat \% mtype = 'ZH'} (negative imaginary parts in upper triangle; see Section~\ref{mat:random}).
  \item \lcodef{call random_number(mat, sigma)}: Calls \procname{mat \% gaussian(sigma=sigma)}, filling \texttt{mat \% v} with Gaussian variates of standard deviation \argname{sigma} and properly Hermitianizing if \texttt{mat \% mtype = 'ZH'} (see Section~\ref{mat:gaussian}). Note that \argname{sigma} is mandatory in this overload (to disambiguate it from \texttt{random\_number(mat)}); pass \texttt{sigma} positionally or by keyword.
\end{itemize}
Both forms abort via \texttt{matrix\_error} if \texttt{mat \% v} has not been allocated.

\manheading{Dummy Arguments}
\begin{description}[leftmargin=2.2cm, style=nextline]
  \argitem{MAT}{class(Matrix), intent(inout)}{Pre-allocated matrix to fill with random entries.}
  \argitem{sigma}{real(8)}{Target Gaussian standard deviation (non-optional in the Gaussian overload).}
\end{description}

\manheading{Example}
\lstset{language=[95]Fortran}
\begin{lstlisting}
use :: matrix_mod_matrix
type(Matrix) :: mat

call matrix_random_init
mat = Matrix(8, mtype='ZH')
call random_number(mat)              ! uniform Hermitian
call random_number(mat, sigma=1.0_8) ! Gaussian Hermitian
\end{lstlisting}

\mansubgroup{Matrix Operations}

% ===================================================================
% trace (Matrix)
% ===================================================================
\manentry{Pure Generic Function}{trace (Matrix)}{Matrix: trace}{Computes the complex trace $\Tr M = \sum_i M_{ii}$ of a \texttt{Matrix}}{mat:trace}

\manheading{Synopsis}
\lstset{language=[95]Fortran}
\begin{lstlisting}
z = trace(m1)
\end{lstlisting}

\manheading{Brief}
Returns the sum of the main-diagonal elements of a square \typename{type(Matrix)} \argname{m1} as a \typename{complex(8)} scalar.

\manheading{Description}
Calls \modname{matrix\_mod\_array}'s \procname{trace(m1 \% v)}, returning $\sum_{i=1}^{N} M_{ii}$ independently of custom index bounds \texttt{is, js}. Returns \procname{NaN} if \argname{m1} is non-square ($m \neq n$).

\manheading{Dummy Arguments}
\begin{description}[leftmargin=2.2cm, style=nextline]
  \argitem{MAT}{class(Matrix), intent(in)}{Square complex matrix.}
  \argitem{r}{complex(8), scalar, return value}{Trace $\Tr M$, or \procname{NaN} if non-square.}
\end{description}

\manheading{Example}
\lstset{language=[95]Fortran}
\begin{lstlisting}
use :: matrix_mod_matrix
type(Matrix) :: m1
complex(8)   :: z

m1 = Matrix('gaussian', 8)
z  = trace(m1)
\end{lstlisting}

% ===================================================================
% trace2 (Matrix)
% ===================================================================
\manentry{Pure Generic Function}{trace2 (Matrix)}{Matrix: trace2}{Computes $\Tr(M^2) = \sum_{i,j} M_{ij} M_{ji}$ in $O(N^2)$ operations}{mat:trace2}

\manheading{Synopsis}
\lstset{language=[95]Fortran}
\begin{lstlisting}
z = trace2(m1)
\end{lstlisting}

\manheading{Brief}
Evaluates $\Tr(M^2)$ directly from matrix elements in $O(N^2)$ operations without matrix multiplication, accelerating Hermitian matrices automatically via \texttt{m1 \% mtype}.

\manheading{Description}
Calls \procname{trace2(m1 \% v, mtype = m1 \% mtype)}. When \texttt{m1 \% mtype = 'ZH'}, unit-stride evaluation $\sum_{i,j} |M_{ij}|^2$ is used instead of strided indexing $\sum_{i,j} M_{ij} M_{ji}$. Note that unlike \modname{array\_mod}, \modname{matrix\_mod\_matrix} overloads only the single-matrix \texttt{trace2(m1)} form for \typename{Matrix} (not two-matrix \texttt{trace2(m1, m2)}). Returns \procname{NaN} if \argname{m1} is non-square.

\manheading{Dummy Arguments}
\begin{description}[leftmargin=2.2cm, style=nextline]
  \argitem{MAT}{class(Matrix), intent(in)}{Square complex matrix whose \texttt{mtype} selects the fast Hermitian path if \texttt{'ZH'}.}
  \argitem{r}{complex(8), scalar, return value}{Quadratic trace $\Tr(M^2)$ (or \procname{NaN} if non-square).}
\end{description}

\manheading{Example}
\lstset{language=[95]Fortran}
\begin{lstlisting}
use :: matrix_mod_matrix
type(Matrix) :: m1
complex(8)   :: z

m1 = Matrix('gaussian', 16, mtype='ZH')
z  = trace2(m1)     ! uses m1 % mtype = 'ZH' automatically
\end{lstlisting}

% ===================================================================
% trace2c (Matrix)
% ===================================================================
\manentry{Pure Generic Function}{trace2c (Matrix)}{Matrix: trace2c}{Computes the connected quadratic trace $\Tr(\tilde{M}^2)$ of the traceless part}{mat:trace2c}

\manheading{Synopsis}
\lstset{language=[95]Fortran}
\begin{lstlisting}
z = trace2c(m1)
\end{lstlisting}

\manheading{Brief}
Computes $\Tr\!\left[\left(M - \frac{\Tr M}{N} I_N\right)^2\right]$, subtracting the center-of-mass trace mode.

\manheading{Description}
Delegates to \procname{trace2c(m1 \% v, mtype = m1 \% mtype)}, creating a local traceless copy $\tilde{M}$ and evaluating $\Tr(\tilde{M}^2)$ using the fast Hermitian formula whenever \texttt{m1 \% mtype = 'ZH'}. Returns \procname{NaN} if \argname{m1} is non-square.

\manheading{Dummy Arguments}
\begin{description}[leftmargin=2.2cm, style=nextline]
  \argitem{MAT}{class(Matrix), intent(in)}{Square complex matrix.}
  \argitem{r}{complex(8), scalar, return value}{Connected quadratic trace $\Tr(\tilde{M}^2)$.}
\end{description}

\manheading{Example}
\lstset{language=[95]Fortran}
\begin{lstlisting}
use :: matrix_mod_matrix
type(Matrix) :: m1
complex(8)   :: zc

m1 = Matrix('gaussian', 16, mtype='ZH')
zc = trace2c(m1)
\end{lstlisting}

% ===================================================================
% norm (Matrix)
% ===================================================================
\manentry{Generic Function}{norm (Matrix)}{Matrix: norm}{Computes the normalized $L_1$ norm $\frac{1}{MN}\sum_{i,j}|M_{ij}|$}{mat:norm}

\manheading{Synopsis}
\lstset{language=[95]Fortran}
\begin{lstlisting}
r = norm(m1)
\end{lstlisting}

\manheading{Brief}
Returns the mean element-wise modulus \texttt{SUM(ABS(m1 \% v)) / SIZE(m1 \% v)} as a \typename{real(8)} scalar.

\manheading{Description}
Invokes the underlying rank-2 norm \texttt{array2\_norm(m1 \% v)}, giving a size-independent measure of matrix magnitude useful for convergence and residual checks.

\manheading{Dummy Arguments}
\begin{description}[leftmargin=2.2cm, style=nextline]
  \argitem{MAT}{class(Matrix), intent(in)}{Allocated complex matrix.}
  \argitem{r}{real(8), scalar, return value}{Mean modulus $\ge 0$.}
\end{description}

\manheading{Example}
\lstset{language=[95]Fortran}
\begin{lstlisting}
use :: matrix_mod_matrix
type(Matrix) :: m1
real(8)      :: err

m1  = Matrix('gaussian', 8)
err = norm(m1 - transpose(transpose(m1)))
\end{lstlisting}

% ===================================================================
% conjg, transpose, hermitian, symmetric (Matrix)
% ===================================================================
\manentry{Generic Functions}{conjg, transpose, hermitian, symmetric (Matrix)}{Matrix: conjg, transpose, hermitian, symmetric}{Return the conjugate, transpose, or Hermitian conjugate of a \texttt{Matrix}}{mat:conj_trans_herm}
\addcontentsline{toc}{subsubsection}{\texorpdfstring{\texttt{\textbf{Matrix: transpose}}}{Matrix: transpose}}
\addcontentsline{toc}{subsubsection}{\texorpdfstring{\texttt{\textbf{Matrix: hermitian}}}{Matrix: hermitian}}
\addcontentsline{toc}{subsubsection}{\texorpdfstring{\texttt{\textbf{Matrix: symmetric}}}{Matrix: symmetric}}

\manheading{Synopsis}
\lstset{language=[95]Fortran}
\begin{lstlisting}
m2 = conjg    (m1)   ! element-wise complex conjugate M^* (same m,n,is,js)
m2 = transpose(m1)   ! transpose M^T (swaps m <-> n and is <-> js)
m2 = hermitian(m1)   ! conjugate transpose M^dagger (swaps m <-> n, is <-> js)
m2 = symmetric(m1)  ! alias for transpose(m1): returns M^T (no symmetrization)
\end{lstlisting}

\manheading{Brief}
Functional equivalents of \procname{m1 \% conjg()}, \procname{m1 \% transpose()}, and \procname{m1 \% hermitian()} returning a new \typename{type(Matrix)}.

\manheading{Description}
\begin{itemize}
  \item \procname{conjg(m1)}: Allocates \argname{m2} with identical metadata and computes \texttt{CONJG(m1 \% v)}.
  \item \procname{transpose(m1)} and \procname{symmetric(m1)}: Both bind to \texttt{matrix\_return\_transpose}, allocating an $n \times m$ matrix with swapped dimensions and bounds (\texttt{m <-> n}, \texttt{is <-> js}, \texttt{ie <-> je}), preserving \texttt{mtype} and \texttt{name}, and computing \texttt{TRANSPOSE(m1 \% v)}. Note that \procname{symmetric(m1)} returns $M^T$ without symmetrizing $M$; use \procname{symmetric\_set} for in-place symmetrization.
  \item \procname{hermitian(m1)}: Allocates an $n \times m$ matrix with swapped dimensions and bounds (\texttt{m <-> n}, \texttt{is <-> js}) and computes \texttt{TRANSPOSE(CONJG(m1 \% v))}.
\end{itemize}

\manheading{Dummy Arguments}
\begin{description}[leftmargin=2.2cm, style=nextline]
  \argitem{MATA}{class(Matrix), intent(in)}{Input matrix.}
  \argitem{MATB}{type(Matrix), return value}{Transformed \typename{Matrix} with preserved \texttt{mtype} and \texttt{name}.}
\end{description}

\manheading{Example}
\lstset{language=[95]Fortran}
\begin{lstlisting}
use :: matrix_mod_matrix
type(Matrix) :: m1, m2, m3

m1 = Matrix('gaussian', 3, 5)
m2 = hermitian(m1)    ! 5x3 conjugate transpose
m3 = conjg(m1)        ! 3x5 complex conjugate
\end{lstlisting}

% ===================================================================
% traceless and traceless_set (Matrix)
% ===================================================================
\manentry{Pure Generic Function / Subroutine}{traceless, traceless\_set (Matrix)}{Matrix: traceless, traceless\textunderscore set}{Extracts or projects onto the traceless part $M - \frac{1}{N}(\Tr M) I_N$}{mat:traceless}
\addcontentsline{toc}{subsubsection}{\texorpdfstring{\texttt{\textbf{Matrix: traceless\textunderscore set}}}{Matrix: traceless\textunderscore set}}

\manheading{Synopsis}
\lstset{language=[95]Fortran}
\begin{lstlisting}
m2 = traceless(m1)     ! pure function: returns traceless copy of m1
call traceless_set(m1) ! pure subroutine: makes m1 traceless in-place
\end{lstlisting}

\manheading{Brief}
Returns a traceless copy (\procname{traceless}) or modifies a square \typename{Matrix} in-place (\procname{traceless\_set}) so that $\Tr M = 0$.

\manheading{Description}
\procname{traceless(m1)} copies all metadata from \argname{m1} (including \texttt{mtype} and \texttt{name}) to \argname{m2} and sets \texttt{m2 \% v = traceless(m1 \% v)}. \procname{traceless\_set(m1)} mutates \texttt{m1 \% v} in-place via \texttt{call traceless\_set(m1 \% v)}. Both are \texttt{pure} and fill the data array with \procname{NaN} if \argname{m1} is non-square.

\manheading{Dummy Arguments}
\begin{description}[leftmargin=2.2cm, style=nextline]
  \argitem{MAT}{class(Matrix), intent(in) or intent(inout)}{Input matrix (\texttt{intent(in)} for \procname{traceless}; \texttt{intent(inout)} for \procname{traceless\_set}).}
  \argitem{MATB}{type(Matrix), return value}{Allocated traceless \typename{Matrix} (for \procname{traceless} only).}
\end{description}

\manheading{Example}
\lstset{language=[95]Fortran}
\begin{lstlisting}
use :: matrix_mod_matrix
type(Matrix) :: m1, m2

m1 = Matrix('gaussian', 6, mtype='ZH')
m2 = traceless(m1)         ! m2 has zero trace, mtype='ZH'
call traceless_set(m1)     ! m1 modified in-place
\end{lstlisting}

% ===================================================================
% hermitian_set, symmetric_set, antisymmetric_set (Matrix)
% ===================================================================
\manentry{Generic Subroutines}{hermitian\_set, symmetric\_set, antisymmetric\_set (Matrix)}{Matrix: hermitian\textunderscore set, symmetric\textunderscore set, antisymmetric\textunderscore set}{Enforce Hermiticity, symmetry, or antisymmetry in-place}{mat:sym_set}
\addcontentsline{toc}{subsubsection}{\texorpdfstring{\texttt{\textbf{Matrix: symmetric\textunderscore set}}}{Matrix: symmetric\textunderscore set}}
\addcontentsline{toc}{subsubsection}{\texorpdfstring{\texttt{\textbf{Matrix: antisymmetric\textunderscore set}}}{Matrix: antisymmetric\textunderscore set}}

\manheading{Synopsis}
\lstset{language=[95]Fortran}
\begin{lstlisting}
call hermitian_set    (m1, [uplo])   ! sets M = M^dagger and m1 % mtype = 'ZH'
call symmetric_set    (m1, [uplo])   ! sets M = M^T
call antisymmetric_set(m1, [uplo])   ! sets M = -M^T (zeros diagonal)
\end{lstlisting}

\manheading{Brief}
Modify a square \typename{Matrix} in-place to be Hermitian ($M = M^\dagger$), complex symmetric ($M = M^T$), or complex antisymmetric ($M = -M^T$).

\manheading{Description}
All three routines accept an optional \argname{uplo} selector controlling how the matrix is symmetrized:
\begin{itemize}
  \item If \argname{uplo(1:1)} is \texttt{'U'}/\texttt{'u'} or \texttt{'L'}/\texttt{'l'}, the upper or lower triangle of \texttt{m1 \% v} is preserved and mirrored into the opposite triangle (conjugating for \procname{hermitian\_set}, negating and zeroing the diagonal for \procname{antisymmetric\_set}).
  \item Otherwise (including when \argname{uplo} is omitted), the projection $\frac{1}{2}(M + \mathrm{op}(M^T))$ is applied ($\mathrm{op} = \mathrm{CONJG}$ for Hermitian, $+$ for symmetric, $-$ for antisymmetric).
\end{itemize}
\procname{hermitian\_set}, \procname{symmetric\_set}, and \procname{antisymmetric\_set} all validate allocation and squareness (aborting via \texttt{matrix\_error} on failure) and update \texttt{m1 \% mtype}: \procname{hermitian\_set} sets \texttt{m1 \% mtype = 'ZH'}, whereas \procname{symmetric\_set} and \procname{antisymmetric\_set} set \texttt{m1 \% mtype = 'ZG'} (since complex symmetric and antisymmetric matrices are not Hermitian in general).

\manheading{Dummy Arguments}
\begin{description}[leftmargin=2.2cm, style=nextline]
  \argitem{MAT}{class(Matrix), intent(inout)}{Square \typename{Matrix} transformed in-place.}
  \argitem{uplo}{character(*), optional, intent(in)}{\texttt{'U'} (keep upper), \texttt{'L'} (keep lower), or omitted/other (average with transpose/adjoint).}
\end{description}

\manheading{Example}
\lstset{language=[95]Fortran}
\begin{lstlisting}
use :: matrix_mod_matrix
type(Matrix) :: m1

m1 = Matrix('gaussian', 6)
call hermitian_set(m1)            ! m1 becomes Hermitian, mtype='ZH'
call antisymmetric_set(m1, 'L')   ! mirror lower triangle with minus sign
\end{lstlisting}

% ===================================================================
% isHermitian, isSymmetric, isAntiSymmetric (Matrix)
% ===================================================================
\manentry{Generic Functions}{isHermitian, isSymmetric, isAntiSymmetric (Matrix)}{Matrix: isHermitian, isSymmetric, isAntiSymmetric}{Measure deviation of a \texttt{Matrix} from Hermiticity, symmetry, or antisymmetry}{mat:is_sym}
\addcontentsline{toc}{subsubsection}{\texorpdfstring{\texttt{\textbf{Matrix: isSymmetric}}}{Matrix: isSymmetric}}
\addcontentsline{toc}{subsubsection}{\texorpdfstring{\texttt{\textbf{Matrix: isAntiSymmetric}}}{Matrix: isAntiSymmetric}}

\manheading{Synopsis}
\lstset{language=[95]Fortran}
\begin{lstlisting}
r = isHermitian    (m1)   ! norm( m1 - hermitian(m1) )
r = isSymmetric    (m1)   ! norm( m1 - transpose(m1) )
r = isAntiSymmetric(m1)   ! norm( m1 + transpose(m1) )
\end{lstlisting}

\manheading{Brief}
Return a non-negative \typename{real(8)} residual norm quantifying how close \argname{m1} is to being Hermitian, symmetric, or antisymmetric.

\manheading{Description}
Calling \modname{matrix\_mod\_array}'s underlying checks on \texttt{m1 \% v}, these functions return $0.0$ (up to floating-point rounding) when \argname{m1} satisfies the corresponding symmetry, or \procname{NaN} if \argname{m1} is non-square:
\begin{align}
  \texttt{isHermitian}(M)     &= \texttt{norm}(M - M^\dagger)\, , \qquad
  \texttt{isSymmetric}(M)      = \texttt{norm}(M - M^T)\, , \nonumber \\
  \texttt{isAntiSymmetric}(M) &= \texttt{norm}(M + M^T)\, .
\end{align}

\manheading{Dummy Arguments}
\begin{description}[leftmargin=2.2cm, style=nextline]
  \argitem{MATA}{class(Matrix), intent(in)}{Input square matrix.}
  \argitem{r}{real(8), scalar, return value}{Normalized $L_1$ deviation ($\ge 0$, or \procname{NaN} if non-square).}
\end{description}

\manheading{Example}
\lstset{language=[95]Fortran}
\begin{lstlisting}
use :: matrix_mod_matrix
type(Matrix) :: m1

m1 = Matrix('gaussian', 8, mtype='ZH')
print *, 'Hermiticity residual:', isHermitian(m1)
\end{lstlisting}

% ===================================================================
% diagonal and diagonalMatrix (Matrix)
% ===================================================================
\manentry{Generic Functions}{diagonal, diagonalMatrix (Matrix)}{Matrix: diagonal, diagonalMatrix}{Extract matrix diagonals, construct diagonal matrices from vectors, or zero off-diagonals}{mat:diagonal}
\addcontentsline{toc}{subsubsection}{\texorpdfstring{\texttt{\textbf{Matrix: diagonalMatrix}}}{Matrix: diagonalMatrix}}

\manheading{Synopsis}
\lstset{language=[95]Fortran}
\begin{lstlisting}
v1 = diagonal      (m1)   ! pure: extract main diagonal into Vector v1
m2 = diagonal      (v1)   ! construct square diagonal Matrix from Vector v1
m2 = diagonalMatrix(m1)   ! keep only main diagonal of m1, zero off-diagonals
\end{lstlisting}

\manheading{Brief}
\procname{diagonal} extracts the main diagonal of a \typename{Matrix} into a \typename{Vector} or builds a diagonal \typename{Matrix} from a \typename{Vector}; \procname{diagonalMatrix} zeroes all off-diagonal entries of an existing \typename{Matrix}.

\manheading{Description}
\begin{itemize}
  \item \lcodef{v1 = diagonal(m1)} (\texttt{pure}): Constructs a \typename{Vector} of length \texttt{min(m1 \% m, m1 \% n)} with lower bound \texttt{v1 \% is = m1 \% is} and copies the diagonal entries \texttt{diagonal(m1 \% v)}.
  \item \lcodef{m2 = diagonal(v1)}: Constructs a square \typename{Matrix} of size \texttt{v1 \% n} with \texttt{is = js = v1 \% is}, \texttt{mtype = 'ZG'}, \texttt{name = ''}, sets \texttt{m2 \% v(i,i) = v1 \% v(i)}, and zeroes off-diagonals. Useful for checking the eigen-equation $A V = V \operatorname{diag}(\lambda)$ after \lcodef{V = eigenvectors(A, lambda)}.
  \item \lcodef{m2 = diagonalMatrix(m1)}: Copies \argname{m1} (including its bounds, \texttt{mtype}, and \texttt{name}) to \argname{m2} and sets \texttt{m2 \% v = diagonalMatrix(m1 \% v)} (keeping the $\min(m,n)$ diagonal elements and zeroing all others; non-square input is allowed and does \emph{not} produce \procname{NaN}).
\end{itemize}

\manheading{Dummy Arguments}
\begin{description}[leftmargin=2.2cm, style=nextline]
  \argitem{MAT / MATA}{class(Matrix), intent(in)}{Input matrix.}
  \argitem{vec}{class(Vector), intent(in)}{Input vector of diagonal entries for \texttt{diagonal(vec)}.}
  \argitem{result}{type(Vector) or type(Matrix), return value}{Extracted diagonal vector or constructed diagonal matrix.}
\end{description}

\manheading{Example}
\lstset{language=[95]Fortran}
\begin{lstlisting}
use :: matrix_mod_matrix
type(Matrix) :: m1, m2, m3
type(Vector) :: v1

m1 = Matrix('gaussian', 4, mtype='ZH')
v1 = diagonal(m1)          ! extract 4-element Vector
m2 = diagonal(v1)          ! rebuild 4x4 diagonal Matrix ('ZG')
m3 = diagonalMatrix(m1)    ! diagonal Matrix keeping m1 % mtype ('ZH')
\end{lstlisting}

% ===================================================================
% real, aimag, mcmplx (Matrix)
% ===================================================================
\manentry{Generic Functions}{real, aimag, mcmplx (Matrix / DMatrix)}{Matrix: real, aimag, mcmplx}{Convert between complex \texttt{Matrix} and real \texttt{DMatrix} parts}{mat:real_aimag_mcmplx}
\addcontentsline{toc}{subsubsection}{\texorpdfstring{\texttt{\textbf{Matrix: aimag}}}{Matrix: aimag}}
\addcontentsline{toc}{subsubsection}{\texorpdfstring{\texttt{\textbf{Matrix: mcmplx}}}{Matrix: mcmplx}}

\manheading{Synopsis}
\lstset{language=[95]Fortran}
\begin{lstlisting}
d1 = real  (m1)       ! DMatrix containing Re(m1 % v); 'ZH' -> 'DS'
d1 = aimag (m1)       ! DMatrix containing Im(m1 % v); mtype = 'DG'
m1 = mcmplx(d1, d2)   ! Matrix  containing CMPLX(d1 % v, d2 % v, kind=8)
\end{lstlisting}

\manheading{Brief}
\procname{real} and \procname{aimag} extract the real or imaginary parts of a \typename{Matrix} as a \typename{DMatrix}; \procname{mcmplx} combines two conformable \typename{DMatrix} objects into a complex \typename{Matrix}.

\manheading{Description}
\begin{itemize}
  \item \procname{real(m1)}: Copies dimensions, bounds, and \texttt{name} from \argname{m1} into \typename{DMatrix} \argname{d1}, setting \texttt{d1 \% mtype = 'DS'} if \texttt{m1 \% mtype == 'ZH'} and \texttt{'DG'} otherwise, and fills \texttt{d1 \% v = REAL(m1 \% v, kind=8)}. Equivalent to \cbox{m1 \% re()}.
  \item \procname{aimag(m1)}: Copies dimensions, bounds, and \texttt{name} into \typename{DMatrix} \argname{d1} with \texttt{d1 \% mtype = 'DG'} and \texttt{d1 \% v = AIMAG(m1 \% v)}. Equivalent to \cbox{m1 \% im()}.
  \item \procname{mcmplx(d1, d2)}: Checks that \argname{d1} and \argname{d2} have the same shape $(m, n)$ (aborting via \texttt{matrix\_error} otherwise), inherits \texttt{m, n, is, js, name} from \argname{d1}, sets \texttt{m1 \% mtype = 'ZG'}, and fills \texttt{CMPLX(d1 \% v, d2 \% v, kind=8)}. To promote a single \typename{DMatrix} to pure real or pure imaginary \typename{Matrix}, multiply by \texttt{(1.0\_8, 0.0\_8)} or \texttt{(0.0\_8, 1.0\_8)} using \procname{*}.
\end{itemize}

\manheading{Dummy Arguments}
\begin{description}[leftmargin=2.2cm, style=nextline]
  \argitem{MATA}{class(Matrix) or class(DMatrix), intent(in)}{Input complex matrix (for \procname{real}/\procname{aimag}) or real-part \typename{DMatrix} (for \procname{mcmplx}).}
  \argitem{MATB}{class(DMatrix), intent(in)}{Imaginary-part \typename{DMatrix} (for \procname{mcmplx}).}
  \argitem{result}{type(DMatrix) or type(Matrix), return value}{Converted real or complex matrix.}
\end{description}

\manheading{Example}
\lstset{language=[95]Fortran}
\begin{lstlisting}
use :: matrix_mod_matrix
type(Matrix)  :: m1, m2
type(DMatrix) :: d1, d2

d1 = DMatrix('gaussian', 4, name='RePart')
d2 = DMatrix('gaussian', 4)
m1 = mcmplx(d1, d2)   ! m1 % v = d1 % v + i * d2 % v, name='RePart'
d1 = real(m1)
d2 = aimag(m1)
m2 = (0.0_8, 1.0_8) * d2   ! pure imaginary Matrix from DMatrix
\end{lstlisting}

\mansubgroup{Linear Algebra}

% ===================================================================
% inverse (Matrix)
% ===================================================================
\manentry{Generic Function}{inverse (Matrix)}{Matrix: inverse}{Computes the matrix inverse $M^{-1}$ via \texttt{LAPACK} LU factorization}{mat:inverse}

\manheading{Synopsis}
\lstset{language=[95]Fortran}
\begin{lstlisting}
m2 = inverse(m1)
\end{lstlisting}

\manheading{Brief}
Computes and returns the matrix inverse $M_1^{-1}$ as a new \typename{type(Matrix)} using \texttt{LAPACK}'s \texttt{zgetrf} and \texttt{zgetri}.

\manheading{Description}
\procname{inverse(m1)} copies all metadata (bounds, \texttt{mtype}, \texttt{name}) from \argname{m1} to \argname{m2}, allocates \texttt{m2 \% v}, and calls \procname{inverse(m1 \% v)} from \modname{matrix\_mod\_array}. If \argname{m1} is singular or non-square, \texttt{m2 \% v} is filled with \procname{NaN} without aborting (test with \lcodef{isNaN(m2)}).

\manheading{Dummy Arguments}
\begin{description}[leftmargin=2.2cm, style=nextline]
  \argitem{MATA}{class(Matrix), intent(in)}{Square non-singular input matrix.}
  \argitem{MATB}{type(Matrix), return value}{Inverse matrix $M_1^{-1}$ with identical metadata (or \procname{NaN} if singular/non-square).}
\end{description}

\manheading{Example}
\lstset{language=[95]Fortran}
\begin{lstlisting}
use :: matrix_mod_matrix
type(Matrix) :: m1, m2

m1 = Matrix('gaussian', 8)
m2 = inverse(m1)
print *, 'Inverse error:', norm(m1 * m2 - Matrix('1', 8))
\end{lstlisting}

% ===================================================================
% determinant (Matrix)
% ===================================================================
\manentry{Generic Function}{determinant (Matrix)}{Matrix: determinant}{Computes the determinant $\det M$ of a square \texttt{Matrix}}{mat:determinant}

\manheading{Synopsis}
\lstset{language=[95]Fortran}
\begin{lstlisting}
z = determinant(m1)
\end{lstlisting}

\manheading{Brief}
Evaluates $\det M_1$ as a \typename{complex(8)} scalar via \texttt{LAPACK} LU factorization (\texttt{zgetrf}).

\manheading{Description}
Delegates to \procname{determinant(m1 \% v)} in \modname{matrix\_mod\_array}. Returns \texttt{(0.0\_8, 0.0\_8)} for singular matrices, \texttt{(1.0\_8, 0.0\_8)} for $0\times 0$, and \procname{NaN} if non-square.

\begin{warnbox}[Overflow Limit]
For $\sigma = 1$ complex Gaussian matrices (\lcodef{Matrix('g', N)}), $|\det M|$ exceeds the largest IEEE double ($\approx 1.8\times 10^{308}$) at about $N \approx 270$; already at $N = 200$, $|\det M| \sim 10^{217}$. For large matrices, always use \procname{lndet}.
\end{warnbox}

\manheading{Dummy Arguments}
\begin{description}[leftmargin=2.2cm, style=nextline]
  \argitem{MAT}{class(Matrix), intent(in)}{Square complex matrix.}
  \argitem{z}{complex(8), scalar, return value}{Complex determinant $\det M$.}
\end{description}

\manheading{Example}
\lstset{language=[95]Fortran}
\begin{lstlisting}
use :: matrix_mod_matrix
type(Matrix) :: m1
complex(8)   :: z

m1 = Matrix('gaussian', 16)
z  = determinant(m1)
\end{lstlisting}

% ===================================================================
% lndet (Matrix)
% ===================================================================
\manentry{Generic Function}{lndet (Matrix)}{Matrix: lndet}{Computes $\ln|\det M|$ and complex phase $\e{i\arg(\det M)}$ without overflow}{mat:lndet}

\manheading{Synopsis}
\lstset{language=[95]Fortran}
\begin{lstlisting}
lz = lndet(m1)       ! lz is complex(8) :: lz(2)
! det(m1) = exp(lz(1)) * lz(2)
\end{lstlisting}

\manheading{Brief}
Returns a 2-element \typename{complex(8)} array containing $\ln|\det M|$ in \texttt{lz(1)} and the unit-modulus complex phase in \texttt{lz(2)}.

\manheading{Description}
Factorizing $M = P L U$ via \texttt{zgetrf}, \procname{lndet(m1)} accumulates $\texttt{lz(1)} = \sum_i \ln|U_{ii}|$ and $\texttt{lz(2)} = (-1)^S \prod_i U_{ii}/|U_{ii}|$ so that $\det M = \rho\, \e{i\phi}$ with $\texttt{lz(1)} = \ln\rho$ and $\texttt{lz(2)} = \e{i\phi}$. When $M$ is singular, \texttt{lz(1) = (}$-\infty$\texttt{, 0)} and \texttt{lz(2) = (0, 0)}; when non-square, both elements are \procname{NaN}.

\manheading{Dummy Arguments}
\begin{description}[leftmargin=2.2cm, style=nextline]
  \argitem{MAT}{class(Matrix), intent(in)}{Square complex matrix.}
  \argitem{z}{complex(8), dimension(2), return value}{\texttt{z(1)} $= \ln|\det M|$ and \texttt{z(2)} $= \e{i\arg(\det M)}$.}
\end{description}

\manheading{Example}
\lstset{language=[95]Fortran}
\begin{lstlisting}
use :: matrix_mod_matrix
type(Matrix) :: m1
complex(8)   :: lz(2), z

m1 = Matrix('gaussian', 256)
lz = lndet(m1)
! For smaller matrices: z = exp(lz(1)) * lz(2) == determinant(m1)
\end{lstlisting}

% ===================================================================
% Pfaffian and lnPfaffian (Matrix)
% ===================================================================
\manentry{Generic Functions}{Pfaffian, lnPfaffian (Matrix)}{Matrix: Pfaffian, lnPfaffian}{Compute the Pfaffian $\Pf M$ or its log magnitude and phase for antisymmetric matrices}{mat:pfaffian}
\addcontentsline{toc}{subsubsection}{\texorpdfstring{\texttt{\textbf{Matrix: lnPfaffian}}}{Matrix: lnPfaffian}}

\manheading{Synopsis}
\lstset{language=[95]Fortran}
\begin{lstlisting}
z  = Pfaffian  (m1)     ! returns complex(8) scalar Pf(m1)
lz = lnPfaffian(m1)     ! returns complex(8) :: lz(2); Pf = exp(lz(1))*lz(2)
\end{lstlisting}

\manheading{Brief}
Compute the Pfaffian $\Pf M$ (satisfying $(\Pf M)^2 = \det M$ for even-dimensional skew-symmetric matrices) and its overflow-safe logarithmic representation using \texttt{PFAPACK}~\cite{Wimmer_2012,pfapack}.

\manheading{Description}
Calling \procname{Pfaffian(m1 \% v)} (\texttt{skpfa}) and \procname{lnPfaffian(m1 \% v)} (\texttt{skpf10}) from \modname{matrix\_mod\_array}:
\begin{itemize}
  \item Input \argname{m1} must be an antisymmetric ($M^T = -M$) complex matrix (e.g.\ prepared with \lcodef{call antisymmetric_set(m1)}).
  \item For odd dimension $N$ or singular matrices, \procname{Pfaffian} returns \texttt{(0.0\_8, 0.0\_8)}, while \procname{lnPfaffian} returns \texttt{lz(1) = (}$-\infty$\texttt{, 0)} and \texttt{lz(2) = (0, 0)}.
  \item For Gaussian antisymmetric matrices (built with \lcodef{Matrix('g', N)} and then made antisymmetric with \procname{antisymmetric\_set}), \procname{Pfaffian} works up to about $N \approx 540$; beyond that it overflows and returns \procname{NaN} (not \texttt{Inf}). Use \procname{lnPfaffian} for larger dimensions.
  \item Note that \procname{Pfaffian} and \procname{lnPfaffian} are defined only for complex \typename{Matrix}, not \typename{DMatrix}.
\end{itemize}

\manheading{Dummy Arguments}
\begin{description}[leftmargin=2.2cm, style=nextline]
  \argitem{MAT}{class(Matrix), intent(in)}{Antisymmetric square complex matrix.}
  \argitem{z}{complex(8), return value}{Scalar Pfaffian $\Pf M$ (for \procname{Pfaffian}) or 2-element array \texttt{[ln|Pf M|, exp(i arg Pf M)]} (for \procname{lnPfaffian}).}
\end{description}

\manheading{Example}
\lstset{language=[95]Fortran}
\begin{lstlisting}
use :: matrix_mod_matrix
type(Matrix) :: m1
complex(8)   :: z, lz(2)

m1 = Matrix('gaussian', 16)
call antisymmetric_set(m1)
z  = Pfaffian(m1)
lz = lnPfaffian(m1)     ! exp(lz(1)) * lz(2) == z
\end{lstlisting}

% ===================================================================
% eigenvalues (Matrix)
% ===================================================================
\manentry{Generic Function}{eigenvalues (Matrix)}{Matrix: eigenvalues}{Computes all eigenvalues of a square \texttt{Matrix} as a \texttt{Vector}}{mat:eigenvalues}

\manheading{Synopsis}
\lstset{language=[95]Fortran}
\begin{lstlisting}
v1 = eigenvalues(m1)
\end{lstlisting}

\manheading{Brief}
Computes the eigenvalues of a square \typename{type(Matrix)} and returns them in a \typename{type(Vector)}, automatically dispatching to Hermitian \texttt{zheev} when \texttt{m1 \% mtype = 'ZH'}.

\manheading{Description}
\procname{eigenvalues(m1)} allocates a \typename{Vector} \argname{v1} of length \texttt{m1 \% n} with lower bound \texttt{v1 \% is = m1 \% is} and invokes \texttt{eigenvalues(m1 \% v, mtype = m1 \% mtype)}:
\begin{itemize}
  \item If \texttt{m1 \% mtype = 'ZH'}, \texttt{LAPACK}'s fast Hermitian eigensolver \texttt{zheev} is called, returning real eigenvalues (with zero imaginary part) sorted in ascending algebraic order.
  \item Otherwise (\texttt{'ZG'}), \texttt{LAPACK}'s general complex eigensolver \texttt{zgeev} is called.
\end{itemize}
Aborts via \texttt{matrix\_error} if \argname{m1} is non-square (\texttt{m1 \% m /= m1 \% n}), and fills \texttt{v1 \% v} with \procname{NaN} if \texttt{LAPACK} fails to converge.

\manheading{Dummy Arguments}
\begin{description}[leftmargin=2.2cm, style=nextline]
  \argitem{MAT}{class(Matrix), intent(in)}{Square complex matrix whose \texttt{mtype} selects \texttt{zheev} (\texttt{'ZH'}) or \texttt{zgeev}.}
  \argitem{vec}{type(Vector), return value}{Allocated vector of $N$ complex eigenvalues with \texttt{is = MAT \% is}.}
\end{description}

\manheading{Example}
\lstset{language=[95]Fortran}
\begin{lstlisting}
use :: matrix_mod_matrix
type(Matrix) :: m1
type(Vector) :: evals

m1    = Matrix('gaussian', 8, mtype='ZH')
evals = eigenvalues(m1)   ! uses zheev automatically
call evals % print()
\end{lstlisting}

% ===================================================================
% eigenvectors (Matrix)
% ===================================================================
\manentry{Generic Function}{eigenvectors (Matrix)}{Matrix: eigenvectors}{Computes eigenvalues and right eigenvectors of a square \texttt{Matrix}}{mat:eigenvectors}

\manheading{Synopsis}
\lstset{language=[95]Fortran}
\begin{lstlisting}
m2 = eigenvectors(m1, v2)   ! v2: eigenvalues, m2 cols: right eigenvectors
m3 = diagonal(v2)           ! diagonal matrix: m1 * m2 == m2 * m3
\end{lstlisting}

\manheading{Brief}
Computes the normalized right eigenvectors of an $N \times N$ \typename{Matrix} \argname{m1} (returned as columns of \argname{m2}) and stores the corresponding eigenvalues in \argname{v2}.

\manheading{Description}
Given a square $N \times N$ \typename{Matrix} \argname{m1} (aborting via \texttt{matrix\_error} if \texttt{m1 \% m /= m1 \% n}), \procname{eigenvectors(m1, v2)} invokes \modname{matrix\_mod\_array}'s \procname{eigenvectors} with \texttt{mtype = m1 \% mtype}:
\begin{itemize}
  \item If \texttt{v2 \% n /= m1 \% n} (e.g.\ if \argname{v2} is uninitialized), \argname{v2} is automatically constructed via \cbox{Vector(m1 \% n, is=m1 \% is)} and populated with the $N$ eigenvalues $\lambda_j$.
  \item Return matrix \argname{m2} copies bounds and \texttt{name} from \argname{m1}, sets \texttt{m2 \% mtype = 'ZG'} (since the eigenvector matrix is unitary, not Hermitian), and stores the normalized right eigenvector corresponding to \texttt{v2 \% v(j)} in column \texttt{j} (\texttt{m2 \% v(:,j)}), satisfying $M_1\, M_2 = M_2\, \operatorname{diag}(v_2)$.
\end{itemize}

\begin{notebox}[Eigenvalue Ordering and Hermitian \texttt{zheev} Dispatch]
\begin{enumerate}
  \item For general complex matrices ($N \gtrsim 100$, \texttt{'ZG'}), \texttt{eigenvalues(m1)} and \texttt{eigenvectors(m1, v2)} may return eigenvalues in \emph{different} orders. Always pair eigenvector matrix \argname{m2} with the eigenvalue vector \argname{v2} produced by the same \procname{eigenvectors} call (never with \texttt{eigenvalues(m1)}).
  \item When \texttt{m1 \% mtype = 'ZH'}, \procname{eigenvectors} forwards \texttt{mtype} to \modname{matrix\_mod\_array}, dispatching to \texttt{LAPACK} \texttt{zheev} (reading only the upper triangle) and returning real eigenvalues in ascending order matching \texttt{eigenvalues(m1)} along with orthonormal eigenvectors.
\end{enumerate}
\end{notebox}

\manheading{Dummy Arguments}
\begin{description}[leftmargin=2.2cm, style=nextline]
  \argitem{MATA}{class(Matrix), intent(in)}{Square $N \times N$ complex matrix.}
  \argitem{vec}{type(Vector), intent(inout)}{Receives the $N$ eigenvalues (auto-constructed if \texttt{vec \% n /= N}).}
  \argitem{MATB}{type(Matrix), return value}{Matrix whose $j$-th column is the right eigenvector corresponding to \texttt{vec \% v(j)}.}
\end{description}

\manheading{Example}
\lstset{language=[95]Fortran}
\begin{lstlisting}
use :: matrix_mod_matrix
type(Matrix) :: m1, m2, m3
type(Vector) :: v2

m1 = Matrix('gaussian', 8)
m2 = eigenvectors(m1, v2)
m3 = diagonal(v2)
print *, 'Eigenproblem residual:', norm(m1 * m2 - m2 * m3)
\end{lstlisting}

\mansubgroup{Math Functions}

% ===================================================================
% abs (Matrix)
% ===================================================================
\manentry{Pure Generic Function}{abs (Matrix)}{Matrix: abs}{Computes the element-wise complex modulus $|M_{ij}|$, returning a \texttt{DMatrix}}{mat:abs}

\manheading{Synopsis}
\lstset{language=[95]Fortran}
\begin{lstlisting}
d1 = abs(m1)
\end{lstlisting}

\manheading{Brief}
Applies Fortran's element-wise \texttt{abs} to \texttt{m1 \% v} and returns a \typename{type(DMatrix)} with preserved dimensions, bounds, and \texttt{name}.

\manheading{Description}
This \texttt{pure} function copies metadata from \argname{m1} into a \typename{DMatrix} \argname{d1}, setting \texttt{d1 \% mtype = 'DS'} if \texttt{m1 \% mtype == 'ZH'} (since $|M_{ij}| = |M_{ji}^*|$ makes $|M|$ real symmetric when $M$ is Hermitian) and \texttt{'DG'} otherwise.

\manheading{Dummy Arguments}
\begin{description}[leftmargin=2.2cm, style=nextline]
  \argitem{MAT}{class(Matrix), intent(in)}{Input complex matrix.}
  \argitem{MATB}{type(DMatrix), return value}{Real matrix of element-wise moduli (\texttt{'DS'} if \argname{MAT} is \texttt{'ZH'}, else \texttt{'DG'}).}
\end{description}

\manheading{Example}
\lstset{language=[95]Fortran}
\begin{lstlisting}
use :: matrix_mod_matrix
type(Matrix)  :: m1
type(DMatrix) :: d1

m1 = Matrix('gaussian', 4, mtype='ZH')
d1 = abs(m1)     ! d1 % mtype == 'DS'
\end{lstlisting}

% ===================================================================
% sin, cos, exp, log, sqrt (Matrix)
% ===================================================================
\manentry{Pure Generic Functions}{sin, cos, exp, log, sqrt (Matrix)}{Matrix: sin, cos, exp, log, sqrt}{Element-wise complex mathematical functions on \texttt{Matrix}}{mat:math_funcs}
\addcontentsline{toc}{subsubsection}{\texorpdfstring{\texttt{\textbf{Matrix: cos}}}{Matrix: cos}}
\addcontentsline{toc}{subsubsection}{\texorpdfstring{\texttt{\textbf{Matrix: exp}}}{Matrix: exp}}
\addcontentsline{toc}{subsubsection}{\texorpdfstring{\texttt{\textbf{Matrix: log}}}{Matrix: log}}
\addcontentsline{toc}{subsubsection}{\texorpdfstring{\texttt{\textbf{Matrix: sqrt}}}{Matrix: sqrt}}

\manheading{Synopsis}
\lstset{language=[95]Fortran}
\begin{lstlisting}
m2 = sin (m1);   m2 = cos (m1);   m2 = exp (m1)   ! preserves m1 % mtype
m2 = log (m1);   m2 = sqrt(m1)                    ! sets m2 % mtype = 'ZG'
\end{lstlisting}

\manheading{Brief}
Evaluate the standard complex functions $\sin, \cos, \exp, \ln, \sqrt{\cdot}$ element-by-element on \texttt{m1 \% v}, returning a new \typename{type(Matrix)}.

\manheading{Description}
All five functions are \texttt{pure}, preserve shape, index bounds, and \texttt{name}, and operate element-wise (not as matrix functions defined by power series). \procname{sin}, \procname{cos}, and \procname{exp} preserve \texttt{m1 \% mtype}, whereas \procname{log} and \procname{sqrt} always set \texttt{m2 \% mtype = 'ZG'} because branch cuts on negative diagonal or off-diagonal elements can break Hermiticity.

\manheading{Dummy Arguments}
\begin{description}[leftmargin=2.2cm, style=nextline]
  \argitem{MAT}{class(Matrix), intent(in)}{Input complex matrix.}
  \argitem{MATB}{type(Matrix), return value}{Result \typename{Matrix} with element-wise function values.}
\end{description}

\manheading{Example}
\lstset{language=[95]Fortran}
\begin{lstlisting}
use :: matrix_mod_matrix
type(Matrix) :: m1, m2

m1 = Matrix('uniform', 4) + 1.0_8
m2 = exp(m1) + sin(m1) - log(m1)
\end{lstlisting}

\mansubgroup{Misc Procedures}

% ===================================================================
% metadata_copy (Matrix)
% ===================================================================
\manentry{Pure Generic Subroutine}{metadata\_copy (Matrix / DMatrix)}{Matrix: metadata\textunderscore copy}{Copies matrix metadata without allocating or copying elements}{mat:metadata_copy}

\manheading{Synopsis}
\lstset{language=[95]Fortran}
\begin{lstlisting}
call metadata_copy(mat_src, mat_dst)
\end{lstlisting}

\manheading{Brief}
Copies dimensions, index bounds, \texttt{mtype}, and \texttt{name} between any pair of \typename{Matrix} or \typename{DMatrix} objects without allocating or copying \texttt{v}.

\manheading{Description}
\procname{metadata\_copy} accepts any combination of \typename{Matrix} and \typename{DMatrix} for \argname{mat\_src} and \argname{mat\_dst} (via \texttt{class(MatrixClass)}). It copies \texttt{m, n, is, ie, js, je, mtype, name}, and automatically adjusts \texttt{mtype} on cross-type copies:
\begin{itemize}
  \item \typename{DMatrix} $\to$ \typename{Matrix}: if source \texttt{mtype} starts with \texttt{'D'}, sets destination \texttt{mtype = 'ZG'}.
  \item \typename{Matrix} $\to$ \typename{DMatrix}: if source \texttt{mtype} starts with \texttt{'Z'}, sets destination \texttt{mtype = 'DG'}.
\end{itemize}
Note that \texttt{mat\_dst \% v} is \emph{not} allocated or modified; allocate it explicitly using \texttt{is:ie, js:je} afterwards.

\manheading{Dummy Arguments}
\begin{description}[leftmargin=2.2cm, style=nextline]
  \argitem{MATA}{class(MatrixClass), intent(in)}{Source \typename{Matrix} or \typename{DMatrix}.}
  \argitem{MATB}{class(MatrixClass), intent(inout)}{Destination \typename{Matrix} or \typename{DMatrix} receiving metadata.}
\end{description}

\manheading{Example}
\lstset{language=[95]Fortran}
\begin{lstlisting}
use :: matrix_mod_matrix
type(Matrix)  :: m1, m2
type(DMatrix) :: d1

m1 = Matrix(4, 3, is=0, js=2, name='Src')
call metadata_copy(m1, m2)
allocate(m2 % v(m2 % is : m2 % ie, m2 % js : m2 % je))
call metadata_copy(m1, d1)   ! d1 % mtype becomes 'DG'
\end{lstlisting}

% ===================================================================
% isNaN (Matrix)
% ===================================================================
\manentry{Generic Function}{isNaN (Matrix)}{Matrix: isNaN}{Tests whether any entry of a \texttt{Matrix} contains IEEE \texttt{NaN}}{mat:isNaN}

\manheading{Synopsis}
\lstset{language=[95]Fortran}
\begin{lstlisting}
has_nan = isNaN(m1)
\end{lstlisting}

\manheading{Brief}
Returns \texttt{.true.} if any element of \texttt{m1 \% v} has a \texttt{NaN} real or imaginary component, and \texttt{.false.} otherwise.

\manheading{Description}
Calls \procname{isNaN(m1 \% v)} from \modname{matrix\_mod\_array} (which evaluates \texttt{ANY(IEEE\_IS\_NAN(REAL(v)) .OR. IEEE\_IS\_NAN(AIMAG(v)))}). Use this check after \procname{inverse}, \procname{eigenvalues}, or \procname{eigenvectors} to detect singularity or non-convergence without halting execution. See Section~\ref{mod:NaN} for generating \procname{NaN} values.

\manheading{Dummy Arguments}
\begin{description}[leftmargin=2.2cm, style=nextline]
  \argitem{MAT}{class(Matrix), intent(in)}{Matrix to inspect.}
  \argitem{itis}{logical, scalar, return value}{\texttt{.true.} if any entry contains \texttt{NaN}, else \texttt{.false.}.}
\end{description}

\manheading{Example}
\lstset{language=[95]Fortran}
\begin{lstlisting}
use :: matrix_mod_matrix
type(Matrix) :: m1
real(8)      :: r = 0.0_8

m1 = Matrix(4)
m1 % v(2,3) = NaN(r)
print *, 'Contains NaN?', isNaN(m1)   ! prints T
\end{lstlisting}


% Section 3: Class DMatrix (Components, Constructors, Component Procedures, Operators, Procedures)
% ===================================================================
% Section 3: Class DMatrix (Components, Constructors, Component Procedures, Operators)
% ===================================================================
\section{Class \texttt{DMatrix}}
\label{sec:class_dmatrix}

The \typename{DMatrix} derived type provides a structured interface for double-precision real (\typename{real(8)} / \typename{real(real64)}) matrix manipulation. Like \typename{Matrix}, it extends the private base type \texttt{MatrixClass}, sharing dimension, index-bound, symmetry, and label metadata as well as inherited I/O and random-filling methods.

% -------------------------------------------------------------------
\subsection{Components}
\label{subsec:dmat_components}

% ===================================================================
% type(DMatrix)
% ===================================================================
\manentry{Derived Type}{type(DMatrix)}{DMatrix (type)}{Double-precision real matrix class extending the private \texttt{MatrixClass}}{dmat:type}

\manheading{Synopsis}
\lstset{language=[95]Fortran}
\begin{lstlisting}
type(DMatrix) :: mat
! Access elements and metadata:
! mat % m, mat % n, mat % is, mat % ie, mat % js, mat % je
! mat % mtype, mat % name, mat % v(is:ie, js:je)
\end{lstlisting}

\manheading{Brief}
Encapsulates an allocatable $m \times n$ \typename{real(8)} array \texttt{v(is:ie, js:je)} with dimension, index-bound, symmetry-type, and name metadata.

\manheading{Description}
An object of type \typename{DMatrix} holds a real $m \times n$ matrix indexed over rows \texttt{is:ie} (\texttt{ie = is + m - 1}) and columns \texttt{js:je} (\texttt{je = js + n - 1}). Before construction, \texttt{m = n = is = ie = js = je = 0} and \texttt{mtype = 'GG'}. After construction, \texttt{mtype} is \texttt{'DG'} (general double-precision real matrix) or \texttt{'DS'} (real symmetric matrix, $D^T = D$). When \texttt{mtype = 'DS'}, \procname{*} (for \typename{DMatrix} $\times$ \typename{DMatrix} and \typename{DMatrix} $\times$ array), \procname{mmmult}, \procname{mvmult}, \procname{trace2}, \procname{trace2c}, \procname{eigenvalues}, and \procname{eigenvectors} automatically select optimized symmetric \texttt{BLAS}/\texttt{LAPACK} paths (which read only the upper triangle) or $O(N^2)$ formulas. The operator \texttt{d1 * u1} (\typename{DMatrix} $\times$ \typename{DVector}) always uses the general \texttt{dgemv}; use \procname{mvmult} for the symmetric \texttt{dsymv} path. Note that \texttt{name} has length \texttt{name\_len = 20} (inherited from \texttt{MatrixClass}).

\manheading{Data Components}
\begin{description}[leftmargin=2.2cm, style=nextline]
  \argitem{m, n}{integer, default 0}{Number of rows (\texttt{m}) and columns (\texttt{n}); \texttt{m == 0} indicates an uninitialized \typename{DMatrix}.}
  \argitem{is, ie}{integer, default 0}{Lower and upper row index bounds satisfying \texttt{ie = is + m - 1}.}
  \argitem{js, je}{integer, default 0}{Lower and upper column index bounds satisfying \texttt{je = js + n - 1}.}
  \argitem{mtype}{character(4), default \texttt{'GG'}}{\texttt{LAPACK}-style matrix type (\texttt{mtype\_len = 4}): \texttt{'DG'} for general real matrices, \texttt{'DS'} for real symmetric matrices, or \texttt{'GG'} if uninitialized.}
  \argitem{name}{character(20), default \texttt{''}}{User-defined matrix label (\texttt{name\_len = 20}) displayed by \procname{\% print}.}
  \argitem{v(:,:)}{real(8), allocatable}{Allocated rank-2 real array storing matrix elements with bounds \texttt{v(is:ie, js:je)}.}
\end{description}

\manheading{Example}
\lstset{language=[95]Fortran}
\begin{lstlisting}
use :: matrix_mod_matrix
type(DMatrix) :: dmat

dmat = DMatrix(4, 3, is=1, js=5, name='D')
print *, dmat % m, dmat % n       ! 4, 3
print *, dmat % mtype             ! 'DG'
print *, dmat % js, dmat % je     ! 5, 7
\end{lstlisting}

% -------------------------------------------------------------------
\subsection{Constructors}
\label{subsec:dmat_constructors}

% ===================================================================
% DMatrix(...) constructor
% ===================================================================
\manentry{Generic Constructor}{DMatrix(...)}{DMatrix (constructor)}{Constructs and allocates a \texttt{type(DMatrix)} object (zero, constant, random, identity, or from 2D real array)}{dmat:constructor}

\manheading{Synopsis}
\lstset{language=[95]Fortran}
\begin{lstlisting}
dmat = DMatrix(m, [n], [is], [js], [mtype], [name])
dmat = DMatrix(C,      [is], [js], [mtype], [name])
dmat = DMatrix(r, m, [n], [is], [js], [mtype], [name])
dmat = DMatrix(z, m, [n], [is], [js], [mtype], [name])
dmat = DMatrix(rtype, m, [n], [is], [js], [mtype], [name], [sigma])
\end{lstlisting}

\manheading{Brief}
Allocates a \typename{type(DMatrix)}, initializes metadata and index bounds, and fills \texttt{dmat \% v} with zeros, a constant scalar, a 2D \typename{real(8)} array, an identity matrix, or uniform/Gaussian random numbers.

\manheading{Description}
All five constructor forms share the metadata rules in \texttt{matrix\_metadata\_put}:
\begin{enumerate}
  \item \textbf{Zero Matrix \texttt{DMatrix(m, ... )}}: Allocates an $m \times n$ matrix filled with \texttt{0.0\_8}. Defaults: \texttt{n = m}, \texttt{is = 1}, \texttt{js = 1} (or \texttt{js = is} if \argname{is} is supplied without \argname{js}), \texttt{mtype = 'DG'}, \texttt{name = ''}. Specifying \argname{m} or \argname{n} $\le 0$, an \argname{mtype} other than \texttt{'DG'}/\texttt{'DS'}, or \texttt{mtype = 'DS'} with $m \neq n$ aborts via \texttt{matrix\_error}.
  \item \textbf{From 2D Real Array \texttt{DMatrix(C, ... )}}: Copies a \typename{real(8)(:,:)} array \argname{C}. Defaults to \texttt{is = js = 1}; pass \texttt{is=lbound(C,1)} and \texttt{js=lbound(C,2)} to preserve arbitrary bounds. If \texttt{mtype = 'DS'}, \cbox{call dmat \% symmetric\_set} is called automatically without checking symmetry of \argname{C}.
  \item \textbf{Constant Scalar \texttt{DMatrix(r, m, ... )} / \texttt{DMatrix(z, m, ... )}}: Sets every entry to \typename{real(8)} scalar \argname{r}, or to \texttt{REAL(z, kind=8)} for a \typename{complex(8)} scalar \argname{z}. Note that the scalar value comes \emph{before} dimension \argname{m}.
  \item \textbf{Random or Identity Matrix \texttt{DMatrix(rtype, m, ... )}}: Dispatches on \argname{rtype(1:1)}:
  \begin{itemize}
    \item \texttt{'r'}/\texttt{'R'}/\texttt{'u'}/\texttt{'U'}: Uniform random numbers in $[0,1)$ via \procname{dmat \% random} (copies lower triangle into upper triangle when \texttt{mtype = 'DS'}).
    \item \texttt{'g'}/\texttt{'G'}: Gaussian random numbers with standard deviation \argname{sigma} (default \texttt{1.0\_8}) via \procname{dmat \% gaussian} (variance-preserving symmetrization when \texttt{mtype = 'DS'}).
    \item \texttt{'1'}: Real $m \times m$ identity matrix (aborts if $m \neq n$).
    \item Passing an empty string (\texttt{len(rtype) < 1}) or any other leading character aborts via \texttt{matrix\_error}.
  \end{itemize}
\end{enumerate}

\manheading{Dummy Arguments}
\begin{description}[leftmargin=2.2cm, style=nextline]
  \argitem{m}{integer, intent(in)}{Number of rows ($m > 0$). Required except in \texttt{DMatrix(C, ...)}.}
  \argitem{n}{integer, optional, intent(in)}{Number of columns ($n > 0$; defaults to \argname{m}). Must equal \argname{m} for \texttt{mtype='DS'} or \texttt{rtype='1'}.}
  \argitem{is, js}{integer, optional, intent(in)}{Starting row and column indices (\texttt{is} defaults to \texttt{1}; \texttt{js} defaults to \texttt{is}).}
  \argitem{mtype}{character(*), optional, intent(in)}{Symmetry tag: \texttt{'DG'} (default) or \texttt{'DS'} (symmetric; requires $m = n$).}
  \argitem{name}{character(*), optional, intent(in)}{Optional name (up to 20 characters; defaults to \texttt{''}).}
  \argitem{C}{real(8), dimension(:,:), intent(in)}{Source 2D real array for \texttt{DMatrix(C, ...)}.}
  \argitem{r / z}{real(8) / complex(8), intent(in)}{Constant fill value (\texttt{REAL(z, kind=8)} if complex), passed before \argname{m}.}
  \argitem{rtype}{character(*), intent(in)}{Generator selector: \texttt{'uniform'}/\texttt{'random'}, \texttt{'gaussian'}, or \texttt{'1'}.}
  \argitem{sigma}{real(8), optional, intent(in)}{Gaussian standard deviation (defaults to \texttt{1.0\_8}).}
  \argitem{MAT}{type(DMatrix), return value}{Allocated \typename{DMatrix} object.}
\end{description}

\manheading{Example}
\lstset{language=[95]Fortran}
\begin{lstlisting}
use :: matrix_mod_matrix
real(8)       :: A(0:2, 0:2) = 1.0_8
type(DMatrix) :: d0, d1, d2, d3

d0 = DMatrix(4, mtype='DS', name='SymZero')
d1 = DMatrix(3.14159_8, 4)
d2 = DMatrix(A, is=lbound(A,1), js=lbound(A,2), mtype='DS')
d3 = DMatrix('gaussian', 8, mtype='DS', sigma=0.5_8)
\end{lstlisting}

% -------------------------------------------------------------------
\subsection{Component Procedures}
\label{subsec:dmat_comp_procs}

Objects of \typename{type(DMatrix)} inherit five type-bound procedures from \texttt{MatrixClass} (\procname{save}, \procname{read}, \procname{print}, \procname{random}, \procname{gaussian}) and define four \typename{DMatrix}-specific type-bound generics: \procname{symmetric\_set}, \procname{traceless\_set}, \procname{transpose}, and \procname{symmetric}.

% ===================================================================
% DMAT % save
% ===================================================================
\manentry{Type-Bound Subroutine}{DMAT \% save}{DMatrix \% save}{Writes real matrix elements in plain text column-major order}{dmat:save}

\manheading{Synopsis}
\lstset{language=[95]Fortran}
\begin{lstlisting}
call dmat % save([unit], [fmt])
\end{lstlisting}

\manheading{Brief}
Saves all raw \typename{real(8)} elements of \texttt{dmat \% v} in column-major order to \argname{unit}.

\manheading{Description}
Iterates in column-major order (\texttt{do j = js, je; do i = is, ie}) and writes each \texttt{dmat \% v(i,j)} using \texttt{'('//fmt//')'} (default \texttt{fmt = 'G28.17'}) to \argname{unit} (default \procname{f\_mout}). Returns silently if \texttt{dmat \% v} is unallocated.

\manheading{Dummy Arguments}
\begin{description}[leftmargin=2.2cm, style=nextline]
  \argitem{MAT}{class(DMatrix), intent(in)}{Passed-object real matrix.}
  \argitem{unit}{integer, optional, intent(in)}{Fortran output unit (defaults to \procname{f\_mout}).}
  \argitem{fmt}{character(*), optional, intent(in)}{Real edit descriptor (defaults to \texttt{'G28.17'}).}
\end{description}

\manheading{Example}
\lstset{language=[95]Fortran}
\begin{lstlisting}
use :: matrix_mod_matrix
type(DMatrix) :: d1
integer       :: un

d1 = DMatrix('uniform', 4)
open(newunit=un, file='dmat.dat', status='replace')
call d1 % save(unit=un, fmt='ES22.14')
close(un)
\end{lstlisting}

% ===================================================================
% DMAT % read
% ===================================================================
\manentry{Type-Bound Subroutine}{DMAT \% read}{DMatrix \% read}{Reads real matrix elements in column-major order from a Fortran unit}{dmat:read}

\manheading{Synopsis}
\lstset{language=[95]Fortran}
\begin{lstlisting}
call dmat % read([unit])
\end{lstlisting}

\manheading{Brief}
Reads plain-text real numbers into a pre-allocated \texttt{dmat \% v} in column-major order via list-directed \texttt{read(un,*)}.

\manheading{Description}
Reads one \typename{real(8)} scalar per entry in column-major order (\texttt{do j = js, je; do i = is, ie}) from \argname{unit} (default internal \texttt{f\_minput = stdin}). Returns silently if \texttt{dmat \% v} is not allocated.

\manheading{Dummy Arguments}
\begin{description}[leftmargin=2.2cm, style=nextline]
  \argitem{MAT}{class(DMatrix), intent(inout)}{Pre-constructed real matrix receiving data.}
  \argitem{unit}{integer, optional, intent(in)}{Fortran input unit (defaults to \texttt{stdin}).}
\end{description}

\manheading{Example}
\lstset{language=[95]Fortran}
\begin{lstlisting}
use :: matrix_mod_matrix
type(DMatrix) :: d2
integer       :: un

d2 = DMatrix(4)
open(newunit=un, file='dmat.dat', status='old')
call d2 % read(unit=un)
close(un)
\end{lstlisting}

% ===================================================================
% DMAT % print
% ===================================================================
\manentry{Type-Bound Subroutine}{DMAT \% print}{DMatrix \% print}{Pretty-prints a \texttt{DMatrix}, metadata summary, or Mathematica expression}{dmat:print}

\manheading{Synopsis}
\lstset{language=[95]Fortran}
\begin{lstlisting}
call dmat % print([unit], [fmt], [form], [ips], [ipe], [jps], [jpe])
\end{lstlisting}

\manheading{Brief}
Prints a \typename{DMatrix} (or sub-block \texttt{ips:ipe, jps:jpe}) in \texttt{'Default'}, \texttt{'Info'}, or \texttt{'Mathematica'} format.

\manheading{Description}
Shares the same inherited implementation (\texttt{matrix\_print}) as \procname{MAT \% print}, formatting real rows with \texttt{(n//fmt)} in \texttt{'Default'} mode, printing \texttt{(v not allocated)} in \texttt{'Info'} mode when \texttt{dmat \% v} is unallocated, and emitting real \texttt{Interpreter["Number"]} nested lists in \texttt{'Mathematica'} mode. Aborts via \texttt{matrix\_error} if uninitialized (\texttt{m == 0}), unallocated (in \texttt{'Default'} or \texttt{'Mathematica'} mode), of unknown subclass, or given an unknown \argname{form}.

\manheading{Dummy Arguments}
\begin{description}[leftmargin=2.2cm, style=nextline]
  \argitem{MAT}{class(DMatrix), intent(in)}{Initialized real matrix.}
  \argitem{unit}{integer, optional, intent(in)}{Destination Fortran unit (defaults to \procname{f\_mout}).}
  \argitem{fmt}{character(*), optional, intent(in)}{Numeric format string (defaults to \texttt{'G28.17'}).}
  \argitem{form}{character(*), optional, intent(in)}{Output mode: \texttt{'Default'} (default), \texttt{'Info'}, or \texttt{'Mathematica'}.}
  \argitem{ips, ipe, jps, jpe}{integer, optional, intent(in)}{Optional slice bounds clamped to \texttt{[is,ie]} and \texttt{[js,je]}.}
\end{description}

\manheading{Example}
\lstset{language=[95]Fortran}
\begin{lstlisting}
use :: matrix_mod_matrix
type(DMatrix) :: d

d = DMatrix('gaussian', 6, mtype='DS', name='Sym')
call d % print(fmt='F8.4')
call d % print(form='Info')
\end{lstlisting}

% ===================================================================
% DMAT % random
% ===================================================================
\manentry{Type-Bound Subroutine}{DMAT \% random}{DMatrix \% random}{Fills a \texttt{DMatrix} with uniform random numbers in $[0,1)$, respecting \texttt{mtype='DS'}}{dmat:random}

\manheading{Synopsis}
\lstset{language=[95]Fortran}
\begin{lstlisting}
call dmat % random()
\end{lstlisting}

\manheading{Brief}
Fills \texttt{dmat \% v} with uniform $[0,1)$ random numbers, symmetrizing without changing the distribution when \texttt{dmat \% mtype = 'DS'}.

\manheading{Description}
Calls \texttt{array2\_random\_set\_d(dmat \% v)}. If \texttt{dmat \% mtype = 'DS'}, copies the lower triangle into the upper triangle (\texttt{dmat \% v(is+i, js+j) = dmat \% v(is+j, js+i)}) instead of averaging $(X_1+X_2)/2$, keeping all entries uniformly distributed. Aborts if \texttt{dmat \% v} is unallocated.

\manheading{Dummy Arguments}
\begin{description}[leftmargin=2.2cm, style=nextline]
  \argitem{MAT}{class(DMatrix), intent(inout)}{Pre-allocated real matrix.}
\end{description}

\manheading{Example}
\lstset{language=[95]Fortran}
\begin{lstlisting}
use :: matrix_mod_matrix
type(DMatrix) :: d

d = DMatrix(4, mtype='DS')
call d % random()
\end{lstlisting}

% ===================================================================
% DMAT % gaussian
% ===================================================================
\manentry{Type-Bound Subroutine}{DMAT \% gaussian}{DMatrix \% gaussian}{Fills a \texttt{DMatrix} with Gaussian random numbers of std dev $\sigma$, respecting \texttt{'DS'}}{dmat:gaussian}

\manheading{Synopsis}
\lstset{language=[95]Fortran}
\begin{lstlisting}
call dmat % gaussian([sigma])
\end{lstlisting}

\manheading{Brief}
Populates \texttt{dmat \% v} with zero-mean Gaussian variates of standard deviation \argname{sigma} (default \texttt{1.0\_8}), maintaining $\sigma$ when symmetrizing \texttt{mtype = 'DS'}.

\manheading{Description}
Draws normal variates via \texttt{array2\_gauss\_set\_d}. When \texttt{dmat \% mtype = 'DS'}, generates numbers with $s = \sigma\sqrt{2}$, invokes \cbox{call dmat \% symmetric\_set}, and divides the diagonal by $\sqrt{2}$ so both diagonal and off-diagonal entries have standard deviation $\sigma$. Aborts if unallocated.

\manheading{Dummy Arguments}
\begin{description}[leftmargin=2.2cm, style=nextline]
  \argitem{MAT}{class(DMatrix), intent(inout)}{Pre-allocated real matrix.}
  \argitem{sigma}{real(8), optional, intent(in)}{Gaussian standard deviation (defaults to \texttt{1.0\_8}).}
\end{description}

\manheading{Example}
\lstset{language=[95]Fortran}
\begin{lstlisting}
use :: matrix_mod_matrix
type(DMatrix) :: d

d = DMatrix(6, mtype='DS')
call d % gaussian(sigma=2.0_8)
\end{lstlisting}

% ===================================================================
% DMAT % symmetric_set
% ===================================================================
\manentry{Type-Bound Subroutine}{DMAT \% symmetric\_set}{DMatrix \% symmetric\textunderscore set}{Symmetrizes a square \texttt{DMatrix} in-place and sets \texttt{mtype = 'DS'}}{dmat:comp_symmetric_set}

\manheading{Synopsis}
\lstset{language=[95]Fortran}
\begin{lstlisting}
call dmat % symmetric_set([uplo])
\end{lstlisting}

\manheading{Brief}
Makes an allocated square \typename{DMatrix} symmetric ($D = D^T$) in-place and sets \texttt{dmat \% mtype = 'DS'}.

\manheading{Description}
Checks allocation and squareness (aborting via \texttt{matrix\_error} on failure), calls \procname{symmetric\_set} on \texttt{dmat \% v}, and updates \texttt{dmat \% mtype = 'DS'}. If \argname{uplo(1:1)} is \texttt{'U'} or \texttt{'L'}, the chosen triangle is copied to the opposite triangle; otherwise $D \leftarrow \frac{1}{2}(D + D^T)$ is computed.

\manheading{Dummy Arguments}
\begin{description}[leftmargin=2.2cm, style=nextline]
  \argitem{MAT}{class(DMatrix), intent(inout)}{Allocated square real matrix.}
  \argitem{uplo}{character(*), optional, intent(in)}{\texttt{'U'} (keep upper), \texttt{'L'} (keep lower), or omitted/other (average).}
\end{description}

\manheading{Example}
\lstset{language=[95]Fortran}
\begin{lstlisting}
use :: matrix_mod_matrix
type(DMatrix) :: d

d = DMatrix('gaussian', 5)
call d % symmetric_set()     ! d is now symmetric and d % mtype == 'DS'
\end{lstlisting}

% ===================================================================
% DMAT % traceless_set
% ===================================================================
\manentry{Type-Bound Pure Subroutine}{DMAT \% traceless\_set}{DMatrix \% traceless\textunderscore set}{Makes a square \texttt{DMatrix} traceless in-place}{dmat:comp_traceless_set}

\manheading{Synopsis}
\lstset{language=[95]Fortran}
\begin{lstlisting}
call dmat % traceless_set()
\end{lstlisting}

\manheading{Brief}
Subtracts $\frac{1}{N}(\Tr D)\, I_N$ from the diagonal of \texttt{dmat \% v} in-place so that $\Tr D = 0$.

\manheading{Description}
Bound via \texttt{generic :: traceless\_set => dmatrix\_traceless\_set} on \typename{DMatrix}, this \texttt{pure} method calls \procname{traceless\_set(dmat \% v)}, preserving \texttt{mtype} (\texttt{'DG'} or \texttt{'DS'}) and \texttt{name}. Non-square matrices are filled with \procname{NaN}.

\manheading{Dummy Arguments}
\begin{description}[leftmargin=2.2cm, style=nextline]
  \argitem{MAT}{class(DMatrix), intent(inout)}{Square real matrix modified in-place.}
\end{description}

\manheading{Example}
\lstset{language=[95]Fortran}
\begin{lstlisting}
use :: matrix_mod_matrix
type(DMatrix) :: d

d = DMatrix('gaussian', 6, mtype='DS')
call d % traceless_set()
print *, 'Trace:', trace(d)
\end{lstlisting}

% ===================================================================
% DMAT % transpose and DMAT % symmetric
% ===================================================================
\manentry{Type-Bound Functions}{DMAT \% transpose, DMAT \% symmetric}{DMatrix \% transpose, DMatrix \% symmetric}{Return the transpose $D^T$ as a new \texttt{DMatrix}, swapping $m \leftrightarrow n$ and bounds}{dmat:comp_transpose}
\addcontentsline{toc}{subsubsection}{\texorpdfstring{\texttt{\textbf{DMatrix \% symmetric}}}{DMatrix \% symmetric}}

\manheading{Synopsis}
\lstset{language=[95]Fortran}
\begin{lstlisting}
d2 = d1 % transpose()
d2 = d1 % symmetric()   ! identical alias returning D^T
\end{lstlisting}

\manheading{Brief}
Both methods return the transpose \texttt{TRANSPOSE(d1 \% v)} as a new \typename{DMatrix} with dimensions and index ranges swapped ($m \leftrightarrow n$, \texttt{is} $\leftrightarrow$ \texttt{js}).

\manheading{Description}
Both \procname{\% transpose()} and \procname{\% symmetric()} map to \texttt{matrix\_return\_transpose\_d}, allocating an $n \times m$ \typename{DMatrix} with bounds \texttt{js:je, is:ie}, preserving \texttt{mtype} and \texttt{name}, and setting \texttt{d2 \% v = TRANSPOSE(d1 \% v)}. Note that \procname{\% symmetric()} does \emph{not} symmetrize \argname{d1}; use \procname{\% symmetric\_set} for in-place symmetrization.

\manheading{Dummy Arguments}
\begin{description}[leftmargin=2.2cm, style=nextline]
  \argitem{MATA}{class(DMatrix), intent(in)}{Passed-object $m \times n$ real matrix.}
  \argitem{MATB}{type(DMatrix), return value}{Transposed $n \times m$ \typename{DMatrix}.}
\end{description}

\manheading{Example}
\lstset{language=[95]Fortran}
\begin{lstlisting}
use :: matrix_mod_matrix
type(DMatrix) :: d1, d2, d3

d1 = DMatrix('uniform', 3, 5, is=1, js=4)
d2 = d1 % transpose()   ! 5x3 matrix with is=4, js=1
d3 = d1 % symmetric()   ! identical to d1 % transpose()
\end{lstlisting}

% -------------------------------------------------------------------
\subsection{Operators}
\label{subsec:dmat_operators}

\modname{matrix\_mod\_matrix} overloads assignment, arithmetic, element-wise power, and matrix/vector multiplication operators for \typename{DMatrix}.

\mansubgroup{Assignment}

% ===================================================================
% assignment(=) DMatrix
% ===================================================================
\manentry{Overloaded Operator}{assignment(=) (DMatrix)}{DMatrix: assignment(=)}{Assigns another \texttt{DMatrix}, 2D real array, or scalar to a \texttt{DMatrix}}{dmat:op_assign}

\manheading{Synopsis}
\lstset{language=[95]Fortran}
\begin{lstlisting}
dmat = d1     ! type(DMatrix)        -> copies metadata and values
dmat = A      ! real(8)    :: A(:,:) -> copies array A into dmat % v
dmat = r      ! real(8)    :: r      -> fills pre-allocated dmat % v = r
dmat = z      ! complex(8) :: z      -> fills pre-allocated dmat % v = Re(z)
\end{lstlisting}

\manheading{Brief}
Overloads \texttt{=} for assigning a \typename{DMatrix}, 2D \typename{real(8)} array, \typename{real(8)} scalar, or \typename{complex(8)} scalar (taking its real part) to a \typename{DMatrix}.

\manheading{Description}
\begin{itemize}
  \item \texttt{dmat = d1} (\typename{type(DMatrix)}): Copies all metadata and automatically allocates/reallocates \texttt{dmat \% v}.
  \item \texttt{dmat = r} (\typename{real(8)}) and \texttt{dmat = z} (\typename{complex(8)}): Require \texttt{dmat \% v} to be already allocated (else \texttt{matrix\_error}); fill \texttt{dmat \% v = r} or \texttt{REAL(z, kind=8)} while keeping metadata unchanged.
  \item \texttt{dmat = A} (\typename{real(8)(:,:)}): If \texttt{dmat \% m == 0} (uninitialized), initializes \texttt{m, n, is, ie, js, je} from \argname{A}, sets \texttt{mtype = 'DG'}, \texttt{name = ''}, and copies \argname{A}. If already constructed (\texttt{dmat \% m /= 0}), checks shape conformance (aborting via \texttt{matrix\_error} on mismatch), preserves existing \texttt{is, js, name}, resets \texttt{dmat \% mtype = 'DG'}, and copies \argname{A} into \texttt{dmat \% v}.
\end{itemize}

\begin{warnbox}[Metadata Reset and Unsupported Right-Hand Types]
\begin{enumerate}
  \item Assigning a raw real array \lcodef{dmat = A} resets \texttt{dmat \% mtype = 'DG'} (since a general array need not be symmetric). Call \cbox{call dmat \% symmetric\_set} afterwards if \argname{A} is symmetric.
  \item \texttt{dmat = m1} (from \typename{Matrix}) and \texttt{dmat = C} (from \typename{complex(8)(:,:)}) are \emph{not} overloaded for \typename{DMatrix}; use \lcodef{dmat = real(m1)} or \cbox{dmat = m1 \% re()}.
\end{enumerate}
\end{warnbox}

\manheading{Dummy Arguments}
\begin{description}[leftmargin=2.2cm, style=nextline]
  \argitem{MATB}{type(DMatrix), intent(inout)}{Destination \typename{DMatrix}.}
  \argitem{rhs}{type(DMatrix), real(8)(:,:), real(8), or complex(8), intent(in)}{Right-hand value.}
\end{description}

\manheading{Example}
\lstset{language=[95]Fortran}
\begin{lstlisting}
use :: matrix_mod_matrix
type(DMatrix) :: d1, d2
real(8)       :: A(4,4) = 2.0_8

d1 = DMatrix(4, is=0, name='Target')
d1 = A                    ! keeps is=0, name='Target', mtype='DG'
d2 = DMatrix(4)
d2 = (3.5_8, 9.9_8)       ! stores real part 3.5_8 everywhere
\end{lstlisting}

\mansubgroup{Addition, Subtraction, Scalar Multiplication, Division, and Power}

% ===================================================================
% operator(+) DMatrix
% ===================================================================
\manentry{Overloaded Operator}{operator(+) (DMatrix)}{DMatrix: operator(+)}{Adds a \texttt{DMatrix} with a real scalar, 2D real array, or another \texttt{DMatrix}}{dmat:op_plus}

\manheading{Synopsis}
\lstset{language=[95]Fortran}
\begin{lstlisting}
d2 = d1 + r;   d2 = r + d1     ! real(8) scalar + DMatrix (pure)
d2 = d1 + A;   d2 = A + d1     ! real(8)(:,:) array + DMatrix
d3 = d1 + d2                   ! DMatrix + DMatrix
\end{lstlisting}

\manheading{Brief}
Computes element-wise sums involving \typename{DMatrix}, enforcing conformability and propagating symmetry tags.

\manheading{Description}
\texttt{d1 + r} and \texttt{r + d1} preserve all metadata (including \texttt{'DS'}). \texttt{d1 + A} and \texttt{A + d1} verify shape matching (calling \texttt{matrix\_error} otherwise) and set \texttt{mtype = 'DG'}. For \texttt{d1 + d2}, the result gets \texttt{mtype = 'DS'} iff both \argname{d1} and \argname{d2} are \texttt{'DS'}, and \texttt{'DG'} otherwise. Note that \texttt{DMatrix + complex} is not overloaded, while mixed \texttt{DMatrix + Matrix} returns a \typename{Matrix} (Section~\ref{mat:op_plus}).

\manheading{Dummy Arguments}
\begin{description}[leftmargin=2.2cm, style=nextline]
  \argitem{lhs, rhs}{class(DMatrix), real(8)(:,:), or real(8), intent(in)}{Operands.}
  \argitem{result}{type(DMatrix), return value}{Allocated real sum matrix.}
\end{description}

\manheading{Example}
\lstset{language=[95]Fortran}
\begin{lstlisting}
use :: matrix_mod_matrix
type(DMatrix) :: d1, d2, d3

d1 = DMatrix('gaussian', 4, mtype='DS')
d2 = DMatrix('gaussian', 4, mtype='DS')
d3 = d1 + d2 + 1.0_8    ! preserves mtype='DS'
\end{lstlisting}

% ===================================================================
% operator(-) DMatrix
% ===================================================================
\manentry{Overloaded Operator}{operator(-) (DMatrix)}{DMatrix: operator(-)}{Unary negation and binary subtraction for \texttt{DMatrix}}{dmat:op_minus}

\manheading{Synopsis}
\lstset{language=[95]Fortran}
\begin{lstlisting}
d2 = -d1                       ! unary negation
d2 = d1 - r;   d2 = r - d1     ! real(8) scalar and DMatrix (pure)
d2 = d1 - A;   d2 = A - d1     ! real(8)(:,:) array and DMatrix
d3 = d1 - d2                   ! DMatrix - DMatrix
\end{lstlisting}

\manheading{Brief}
Negates a \typename{DMatrix} or subtracts scalars, 2D real arrays, or another \typename{DMatrix} element-wise.

\manheading{Description}
Unary \texttt{-d1} and scalar differences preserve \texttt{d1 \% mtype}; \texttt{d1 - d2} produces \texttt{'DS'} iff both \argname{d1} and \argname{d2} are \texttt{'DS'}, else \texttt{'DG'}, and array differences \texttt{d1 - A} and \texttt{A - d1} reset \texttt{mtype = 'DG'}. Shape mismatches abort via \texttt{matrix\_error}. Mixed \texttt{DMatrix - Matrix} returns a \typename{Matrix} (Section~\ref{mat:op_minus}).

\manheading{Dummy Arguments}
\begin{description}[leftmargin=2.2cm, style=nextline]
  \argitem{lhs, rhs}{class(DMatrix), real(8)(:,:), or real(8), intent(in)}{Operands.}
  \argitem{result}{type(DMatrix), return value}{Allocated difference matrix.}
\end{description}

\manheading{Example}
\lstset{language=[95]Fortran}
\begin{lstlisting}
use :: matrix_mod_matrix
type(DMatrix) :: d1, d2, d3

d1 = DMatrix('gaussian', 4, mtype='DS')
d2 = DMatrix('gaussian', 4, mtype='DS')
d3 = -d1 - d2     ! d3 % mtype == 'DS'
\end{lstlisting}

% ===================================================================
% Scalar * and / DMatrix
% ===================================================================
\manentry{Overloaded Operators}{operator(*) (Scalar), operator(/) (DMatrix)}{DMatrix: scalar * and /}{Scalar multiplication and division for \texttt{DMatrix}}{dmat:op_scale}

\manheading{Synopsis}
\lstset{language=[95]Fortran}
\begin{lstlisting}
d2 = r * d1;   d2 = d1 * r;   d2 = d1 / r   ! returns DMatrix (pure)
m3 = z * d1;   m3 = d1 * z                  ! z * DMatrix -> Matrix (pure)
\end{lstlisting}

\manheading{Brief}
Scales a \typename{DMatrix} by a real scalar (returning a \typename{DMatrix} with preserved metadata) or by a complex scalar (returning a \typename{Matrix}).

\manheading{Description}
The real overloads \texttt{r * d1}, \texttt{d1 * r}, and \texttt{d1 / r} preserve all metadata including \texttt{mtype = 'DS'}. Multiplying a \typename{DMatrix} by a \typename{complex(8)} scalar \argname{z} promotes the result to \typename{type(Matrix)} with \texttt{mtype = 'ZG'} and \texttt{name = d1 \% name} (see also Section~\ref{mat:op_scale}). Note that division \texttt{d1 / z} by a complex scalar is not overloaded for \typename{DMatrix}.

\manheading{Dummy Arguments}
\begin{description}[leftmargin=2.2cm, style=nextline]
  \argitem{r / z}{real(8) or complex(8), scalar, intent(in)}{Scalar multiplier or real divisor.}
  \argitem{MATA}{class(DMatrix), intent(in)}{Input real matrix.}
  \argitem{MATB}{type(DMatrix) or type(Matrix), return value}{Scaled real or promoted complex matrix.}
\end{description}

\manheading{Example}
\lstset{language=[95]Fortran}
\begin{lstlisting}
use :: matrix_mod_matrix
type(DMatrix) :: d1, d2
type(Matrix)  :: m3

d1 = DMatrix('gaussian', 4, mtype='DS')
d2 = 2.5_8 * d1 / 5.0_8
m3 = (0.0_8, 1.0_8) * d1   ! Matrix with pure imaginary entries
\end{lstlisting}

% ===================================================================
% operator(**) DMatrix
% ===================================================================
\manentry{Overloaded Operator}{operator(**) (DMatrix)}{DMatrix: operator(**)}{Computes element-wise exponentiation $(D_{ij})^x$ of a \texttt{DMatrix}}{dmat:op_pow}

\manheading{Synopsis}
\lstset{language=[95]Fortran}
\begin{lstlisting}
d2 = d1 ** k    ! integer k (pure)
d2 = d1 ** x    ! real(8) x (pure)
\end{lstlisting}

\manheading{Brief}
Applies element-wise exponentiation \texttt{(d1 \% v) ** x} for an integer or real exponent, preserving \texttt{d1 \% mtype}.

\manheading{Description}
This \texttt{pure} operator copies all metadata (including \texttt{mtype}) from \argname{d1} and raises each entry to the power \argname{x}. Note that complex exponents are not supported for \typename{DMatrix}, and \texttt{d1 ** k} is element-wise (not matrix power $D^k$).

\manheading{Dummy Arguments}
\begin{description}[leftmargin=2.2cm, style=nextline]
  \argitem{MAT}{class(DMatrix), intent(in)}{Base real matrix.}
  \argitem{n}{integer or real(8), scalar, intent(in)}{Real or integer exponent.}
  \argitem{MATB}{type(DMatrix), return value}{Result \typename{DMatrix} with preserved metadata.}
\end{description}

\manheading{Example}
\lstset{language=[95]Fortran}
\begin{lstlisting}
use :: matrix_mod_matrix
type(DMatrix) :: d1, d2

d1 = DMatrix('uniform', 4)
d2 = d1 ** 0.5_8    ! element-wise square root
\end{lstlisting}

\mansubgroup{Matrix Multiplication}

% ===================================================================
% operator(*) DMatrix Multiplication
% ===================================================================
\manentry{Overloaded Operator}{operator(*) (DMatrix Multiplication)}{DMatrix: matrix *}{BLAS-backed matrix--matrix and matrix--vector multiplication for \texttt{DMatrix}}{dmat:op_matmul}

\manheading{Synopsis}
\lstset{language=[95]Fortran}
\begin{lstlisting}
d3 = d1 * d2;   d3 = d1 * A;   d3 = A * d1   ! DMatrix * DMatrix / array
u2 = d1 * u1;   u2 = u1 * d1                 ! DMatrix * DVector -> DVector
\end{lstlisting}

\manheading{Brief}
Computes real matrix--matrix products ($D_1 D_2$) and real matrix--vector products ($D_1 u_1$ or $u_1^T D_1$) via \texttt{BLAS} (\texttt{dgemm}, \texttt{dsymm}, \texttt{dgemv}).

\manheading{Description}
\begin{itemize}
  \item \textbf{Matrix--Matrix (\texttt{d1*d2}, \texttt{d1*A}, \texttt{A*d1})}: Row bounds \texttt{is:ie} come from the left operand and column bounds \texttt{js:je} from the right operand (\texttt{1:size(A,.)} for raw arrays); \texttt{mtype} is always set to \texttt{'DG'}, and \texttt{name} is copied from the left operand (or from \argname{d1} for \texttt{A * d1}). When the left operand is \texttt{d1} with \texttt{mtype = 'DS'}, \texttt{LMATMUL} uses \texttt{BLAS} \texttt{dsymm}; otherwise \texttt{dgemm}. If the inner dimensions differ (\texttt{d1 \% n /= d2 \% m}), the program aborts via \texttt{matrix\_error} (exit status 1). For mixed \typename{DMatrix} $\times$ \typename{Matrix}, see Section~\ref{mat:op_matmul}.
  \item \textbf{Matrix--Vector (\texttt{d1*u1}, \texttt{u1*d1})}: Returns a \typename{DVector} with \texttt{u2 \% name = u1 \% name} and bounds \texttt{d1\%is:d1\%ie} (for \texttt{d1 * u1}) or \texttt{d1\%js:d1\%je} (for \texttt{u1 * d1}). Calls \texttt{dgemv} (or \texttt{dgemv('T',...)}) and aborts via \texttt{matrix\_error} on a dimension mismatch.
\end{itemize}

\manheading{Dummy Arguments}
\begin{description}[leftmargin=2.2cm, style=nextline]
  \argitem{lhs, rhs}{DMatrix, DVector, or real(8)(:,:), intent(in)}{Conformable real factors.}
  \argitem{result}{type(DMatrix) or type(DVector), return value}{Allocated real matrix or vector product.}
\end{description}

\manheading{Example}
\lstset{language=[95]Fortran}
\begin{lstlisting}
use :: matrix_mod_matrix
type(DMatrix) :: d1, d2, d3
type(DVector) :: u1, u2

d1 = DMatrix('gaussian', 4, is=1)
d2 = DMatrix('gaussian', 4, js=5)
d3 = d1 * d2       ! d3 % v(1:4, 5:8)
u1 = DVector('gaussian', 4)
u2 = d1 * u1       ! u2 % v(1:4)
\end{lstlisting}

% ===================================================================
% mmmult (DMatrix)
% ===================================================================
\manentry{Pure Generic Subroutine}{mmmult (DMatrix)}{DMatrix: mmmult}{High-performance zero-allocation BLAS subroutine for \texttt{DMatrix} multiplication}{dmat:mmmult}

\manheading{Synopsis}
\lstset{language=[95]Fortran}
\begin{lstlisting}
call mmmult(d1, d2, d3)   ! computes d3 % v = d1 % v . d2 % v
\end{lstlisting}

\manheading{Brief}
Computes $D_3 = D_1 D_2$ into \argname{d3} (auto-allocating/resizing \argname{d3} if needed, or in-place with zero allocations when \argname{d3} already has matching shape).

\manheading{Description}
If \texttt{d3 \% v} is unallocated or its dimensions do not match (\texttt{d3 \% m /= d1 \% m .or. d3 \% n /= d2 \% n}), \procname{mmmult} automatically allocates or reconstructs \argname{d3} with bounds \texttt{d1 \% is : d1 \% ie, d2 \% js : d2 \% je}; if \argname{d3} is already allocated with conformable output shape, its existing bounds and \texttt{name} are preserved in-place. It calls \texttt{mmmult(d1 \% v, d2 \% v, d3 \% v, mtype = d1 \% mtype)} (dispatching to \texttt{dsymm} when \texttt{d1 \% mtype = 'DS'} and \texttt{dgemm} otherwise, after checking that \texttt{d1 \% n == d2 \% m}; otherwise the program stops with an \texttt{error stop} message naming the procedure (exit status 1)) and sets \texttt{d3 \% mtype = 'DG'} (since $D_1 D_2$ is not symmetric in general). Aliasing an input as the output (e.g.\ \lcodef{call mmmult(d1, d2, d1)}) is invalid because \argname{MATA}/\argname{MATB} are \texttt{intent(in)} and \argname{MATC} is \texttt{intent(inout)}.

\manheading{Dummy Arguments}
\begin{description}[leftmargin=2.2cm, style=nextline]
  \argitem{MATA}{class(DMatrix), intent(in)}{Left $m \times k$ real factor (\texttt{mtype='DS'} selects \texttt{dsymm}).}
  \argitem{MATB}{class(DMatrix), intent(in)}{Right $k \times n$ real factor.}
  \argitem{MATC}{type(DMatrix), intent(inout)}{Destination $m \times n$ \typename{DMatrix} (auto-allocated/resized if needed; sets \texttt{mtype = 'DG'}).}
\end{description}

\manheading{Example}
\lstset{language=[95]Fortran}
\begin{lstlisting}
use :: matrix_mod_matrix
type(DMatrix) :: d1, d2, d3

d1 = DMatrix('gaussian', 128, mtype='DS')
d2 = DMatrix('gaussian', 128)
d3 = DMatrix(128)
call mmmult(d1, d2, d3)
\end{lstlisting}

% ===================================================================
% Section 3 (cont.): Class DMatrix - Procedures
% ===================================================================
\subsection{Procedures}
\label{subsec:dmat_procs}

This subsection documents all standalone generic procedures accepting \typename{type(DMatrix)} arguments.

\mansubgroup{Random Numbers}

% ===================================================================
% random_number (DMatrix)
% ===================================================================
\manentry{Generic Subroutine}{random\_number (DMatrix)}{DMatrix: random\textunderscore number}{Overloaded uniform and Gaussian random matrix generator for \texttt{DMatrix}}{dmat:random_number}

\manheading{Synopsis}
\lstset{language=[95]Fortran}
\begin{lstlisting}
call random_number(dmat)          ! uniform random numbers in [0,1)
call random_number(dmat, sigma)   ! Gaussian random numbers with std dev sigma
\end{lstlisting}

\manheading{Brief}
Fills a pre-allocated \typename{type(DMatrix)} with uniform $[0,1)$ or Gaussian random numbers via Fortran's overloaded \procname{random\_number}.

\manheading{Description}
Dispatches to \procname{dmat \% random()} (single-argument call) or \procname{dmat \% gaussian(sigma=sigma)} (two-argument call; \argname{sigma} is mandatory to disambiguate the interface). If \texttt{dmat \% mtype = 'DS'}, symmetry is enforced automatically while preserving the uniform distribution or target Gaussian variance $\sigma^2$. Aborts via \texttt{matrix\_error} if \texttt{dmat \% v} is unallocated.

\manheading{Dummy Arguments}
\begin{description}[leftmargin=2.2cm, style=nextline]
  \argitem{MAT}{class(DMatrix), intent(inout)}{Pre-allocated real matrix.}
  \argitem{sigma}{real(8)}{Target Gaussian standard deviation (mandatory when calling the Gaussian form).}
\end{description}

\manheading{Example}
\lstset{language=[95]Fortran}
\begin{lstlisting}
use :: matrix_mod_matrix
type(DMatrix) :: dmat

call matrix_random_init
dmat = DMatrix(8, mtype='DS')
call random_number(dmat)              ! uniform symmetric
call random_number(dmat, sigma=1.0_8) ! Gaussian symmetric
\end{lstlisting}

\mansubgroup{Matrix Operations}

% ===================================================================
% trace, trace2, trace2c (DMatrix)
% ===================================================================
\manentry{Pure Generic Functions}{trace, trace2, trace2c (DMatrix)}{DMatrix: trace, trace2, trace2c}{Compute $\Tr D$, $\Tr(D^2)$, and connected quadratic trace $\Tr(\tilde{D}^2)$}{dmat:traces}
\addcontentsline{toc}{subsubsection}{\texorpdfstring{\texttt{\textbf{DMatrix: trace2}}}{DMatrix: trace2}}
\addcontentsline{toc}{subsubsection}{\texorpdfstring{\texttt{\textbf{DMatrix: trace2c}}}{DMatrix: trace2c}}

\manheading{Synopsis}
\lstset{language=[95]Fortran}
\begin{lstlisting}
r = trace  (d1)   ! sum_i  d1 % v_ii
r = trace2 (d1)   ! sum_ij d1 % v_ij * d1 % v_ji (faster if mtype='DS')
r = trace2c(d1)   ! quadratic trace of traceless part d1 - (Tr d1 / N) I
\end{lstlisting}

\manheading{Brief}
Evaluate scalar traces of a square \typename{type(DMatrix)} \argname{d1} as \typename{real(8)} values, using unit-stride $O(N^2)$ formulas automatically when \texttt{d1 \% mtype = 'DS'}.

\manheading{Description}
All three functions are \texttt{pure} and wrap \modname{matrix\_mod\_array}'s real trace functions:
\begin{itemize}
  \item \procname{trace(d1)}: Computes $\sum_{i=1}^N D_{ii}$.
  \item \procname{trace2(d1)}: Calls \texttt{trace2(d1 \% v, mtype = d1 \% mtype)}. When \texttt{d1 \% mtype = 'DS'}, evaluates $\sum_{i,j} D_{ij}^2$ directly with contiguous access.
  \item \procname{trace2c(d1)}: Calls \texttt{trace2c(d1 \% v, mtype = d1 \% mtype)}, returning $\Tr\!\left[(D - \frac{\Tr D}{N} I_N)^2\right]$.
\end{itemize}
All three return \procname{NaN} if \argname{d1} is non-square.

\manheading{Dummy Arguments}
\begin{description}[leftmargin=2.2cm, style=nextline]
  \argitem{MAT}{class(DMatrix), intent(in)}{Square real matrix.}
  \argitem{r}{real(8), scalar, return value}{Real trace value (or \procname{NaN} if non-square).}
\end{description}

\manheading{Example}
\lstset{language=[95]Fortran}
\begin{lstlisting}
use :: matrix_mod_matrix
type(DMatrix) :: d1
real(8)       :: t1, t2, tc

d1 = DMatrix('gaussian', 8, mtype='DS')
t1 = trace  (d1)
t2 = trace2 (d1)
tc = trace2c(d1)
\end{lstlisting}

% ===================================================================
% norm (DMatrix)
% ===================================================================
\manentry{Generic Function}{norm (DMatrix)}{DMatrix: norm}{Computes the normalized $L_1$ norm $\frac{1}{MN}\sum_{i,j}|D_{ij}|$}{dmat:norm}

\manheading{Synopsis}
\lstset{language=[95]Fortran}
\begin{lstlisting}
r = norm(d1)
\end{lstlisting}

\manheading{Brief}
Returns the mean element-wise absolute value \texttt{SUM(ABS(d1 \% v)) / SIZE(d1 \% v)} as a \typename{real(8)} scalar.

\manheading{Description}
Calls \texttt{array2\_norm\_d(d1 \% v)}, giving a dimension-normalized magnitude convenient for convergence and symmetry residual checks.

\manheading{Dummy Arguments}
\begin{description}[leftmargin=2.2cm, style=nextline]
  \argitem{MAT}{class(DMatrix), intent(in)}{Allocated real matrix.}
  \argitem{r}{real(8), scalar, return value}{Mean absolute element value $\ge 0$.}
\end{description}

\manheading{Example}
\lstset{language=[95]Fortran}
\begin{lstlisting}
use :: matrix_mod_matrix
type(DMatrix) :: d1

d1 = DMatrix('gaussian', 8)
print *, 'Norm(d1):', norm(d1)
\end{lstlisting}

% ===================================================================
% transpose and symmetric (DMatrix)
% ===================================================================
\manentry{Generic Functions}{transpose, symmetric (DMatrix)}{DMatrix: transpose, symmetric}{Return the transpose $D^T$ as a new \texttt{DMatrix}}{dmat:transpose}
\addcontentsline{toc}{subsubsection}{\texorpdfstring{\texttt{\textbf{DMatrix: symmetric}}}{DMatrix: symmetric}}

\manheading{Synopsis}
\lstset{language=[95]Fortran}
\begin{lstlisting}
d2 = transpose(d1)
d2 = symmetric(d1)   ! identical alias returning transpose(d1 % v)
\end{lstlisting}

\manheading{Brief}
Functional equivalents of \procname{d1 \% transpose()} and \procname{d1 \% symmetric()}, returning $D^T$ with swapped dimensions and bounds.

\manheading{Description}
Both \procname{transpose(d1)} and \procname{symmetric(d1)} invoke \texttt{matrix\_return\_transpose\_d}, allocating an $n \times m$ \typename{DMatrix} with bounds \texttt{js:je, is:ie}, preserving \texttt{mtype} and \texttt{name}. Remember that \procname{symmetric(d1)} returns $D^T$ rather than symmetrizing \argname{d1} (use \procname{symmetric\_set} for in-place symmetrization).

\manheading{Dummy Arguments}
\begin{description}[leftmargin=2.2cm, style=nextline]
  \argitem{MATA}{class(DMatrix), intent(in)}{Input $m \times n$ real matrix.}
  \argitem{MATB}{type(DMatrix), return value}{Transposed $n \times m$ \typename{DMatrix}.}
\end{description}

\manheading{Example}
\lstset{language=[95]Fortran}
\begin{lstlisting}
use :: matrix_mod_matrix
type(DMatrix) :: d1, d2

d1 = DMatrix('uniform', 3, 5)
d2 = transpose(d1)   ! 5x3 matrix
\end{lstlisting}

% ===================================================================
% traceless and traceless_set (DMatrix)
% ===================================================================
\manentry{Pure Generic Function / Subroutine}{traceless, traceless\_set (DMatrix)}{DMatrix: traceless, traceless\textunderscore set}{Returns or sets in-place the traceless part $D - \frac{1}{N}(\Tr D) I_N$}{dmat:traceless}
\addcontentsline{toc}{subsubsection}{\texorpdfstring{\texttt{\textbf{DMatrix: traceless\textunderscore set}}}{DMatrix: traceless\textunderscore set}}

\manheading{Synopsis}
\lstset{language=[95]Fortran}
\begin{lstlisting}
d2 = traceless(d1)     ! pure function returning traceless copy of d1
call traceless_set(d1) ! pure subroutine making d1 traceless in-place
\end{lstlisting}

\manheading{Brief}
Subtract $\frac{1}{N}(\Tr D)\, I_N$ from the diagonal of a square \typename{DMatrix}, either returning a new copy (\procname{traceless}) or modifying \argname{d1} in-place (\procname{traceless\_set}).

\manheading{Description}
Both procedures preserve \texttt{d1 \% mtype} and \texttt{d1 \% name} and fill the data array with \procname{NaN} if \argname{d1} is non-square.

\manheading{Dummy Arguments}
\begin{description}[leftmargin=2.2cm, style=nextline]
  \argitem{MAT}{class(DMatrix), intent(in) or intent(inout)}{Input real matrix.}
  \argitem{MATB}{type(DMatrix), return value}{Traceless \typename{DMatrix} copy (for \procname{traceless}).}
\end{description}

\manheading{Example}
\lstset{language=[95]Fortran}
\begin{lstlisting}
use :: matrix_mod_matrix
type(DMatrix) :: d1, d2

d1 = DMatrix('gaussian', 6, mtype='DS')
d2 = traceless(d1)
call traceless_set(d1)
\end{lstlisting}

% ===================================================================
% symmetric_set and antisymmetric_set (DMatrix)
% ===================================================================
\manentry{Generic Subroutines}{symmetric\_set, antisymmetric\_set (DMatrix)}{DMatrix: symmetric\textunderscore set, antisymmetric\textunderscore set}{Make a square \texttt{DMatrix} symmetric or antisymmetric in-place}{dmat:sym_set}
\addcontentsline{toc}{subsubsection}{\texorpdfstring{\texttt{\textbf{DMatrix: antisymmetric\textunderscore set}}}{DMatrix: antisymmetric\textunderscore set}}

\manheading{Synopsis}
\lstset{language=[95]Fortran}
\begin{lstlisting}
call symmetric_set    (d1, [uplo])   ! sets D = D^T and d1 % mtype = 'DS'
call antisymmetric_set(d1, [uplo])   ! sets D = -D^T (zeros diagonal)
\end{lstlisting}

\manheading{Brief}
Project a square \typename{DMatrix} in-place onto its symmetric ($D = D^T$) or antisymmetric ($D = -D^T$) part.

\manheading{Description}
If \argname{uplo(1:1)} is \texttt{'U'} or \texttt{'L'}, the corresponding triangle of \texttt{d1 \% v} is copied (or negated, with zeros on the diagonal for \procname{antisymmetric\_set}) to the opposite triangle; otherwise $\frac{1}{2}(D \pm D^T)$ is formed. Both \procname{symmetric\_set} and \procname{antisymmetric\_set} validate allocation and squareness (aborting via \texttt{matrix\_error} on failure), forward to \modname{matrix\_mod\_array}, and update \texttt{d1 \% mtype} to \texttt{'DS'} (symmetric) or \texttt{'DG'} (antisymmetric), respectively.

\manheading{Dummy Arguments}
\begin{description}[leftmargin=2.2cm, style=nextline]
  \argitem{MAT}{class(DMatrix), intent(inout)}{Square real matrix modified in-place.}
  \argitem{uplo}{character(*), optional, intent(in)}{\texttt{'U'}, \texttt{'L'}, or omitted/other (average).}
\end{description}

\manheading{Example}
\lstset{language=[95]Fortran}
\begin{lstlisting}
use :: matrix_mod_matrix
type(DMatrix) :: d1

d1 = DMatrix('gaussian', 6)
call symmetric_set(d1)             ! d1 % mtype becomes 'DS'
call antisymmetric_set(d1, 'U')    ! keep upper triangle, zero diagonal
\end{lstlisting}

% ===================================================================
% isSymmetric and isAntiSymmetric (DMatrix)
% ===================================================================
\manentry{Generic Functions}{isSymmetric, isAntiSymmetric (DMatrix)}{DMatrix: isSymmetric, isAntiSymmetric}{Measure deviation of a \texttt{DMatrix} from symmetry or antisymmetry}{dmat:is_sym}
\addcontentsline{toc}{subsubsection}{\texorpdfstring{\texttt{\textbf{DMatrix: isAntiSymmetric}}}{DMatrix: isAntiSymmetric}}

\manheading{Synopsis}
\lstset{language=[95]Fortran}
\begin{lstlisting}
r = isSymmetric    (d1)   ! norm( d1 - transpose(d1) )
r = isAntiSymmetric(d1)   ! norm( d1 + transpose(d1) )
\end{lstlisting}

\manheading{Brief}
Return a \typename{real(8)} residual norm measuring how far \argname{d1} is from $D = D^T$ or $D = -D^T$.

\manheading{Description}
Returns $0.0$ (within machine precision) for exact symmetric or skew-symmetric matrices, or \procname{NaN} if \argname{d1} is non-square.

\manheading{Dummy Arguments}
\begin{description}[leftmargin=2.2cm, style=nextline]
  \argitem{MATA}{class(DMatrix), intent(in)}{Input square real matrix.}
  \argitem{r}{real(8), scalar, return value}{Normalized $L_1$ deviation ($\ge 0$, or \procname{NaN} if non-square).}
\end{description}

\manheading{Example}
\lstset{language=[95]Fortran}
\begin{lstlisting}
use :: matrix_mod_matrix
type(DMatrix) :: d1

d1 = DMatrix('gaussian', 8)
call antisymmetric_set(d1)
print *, 'Antisymmetric residual:', isAntiSymmetric(d1)
\end{lstlisting}

% ===================================================================
% diagonal and diagonalMatrix (DMatrix)
% ===================================================================
\manentry{Generic Functions}{diagonal, diagonalMatrix (DMatrix)}{DMatrix: diagonal, diagonalMatrix}{Extract real diagonals, construct diagonal \texttt{DMatrix} from \texttt{DVector}, or zero off-diagonals}{dmat:diagonal}
\addcontentsline{toc}{subsubsection}{\texorpdfstring{\texttt{\textbf{DMatrix: diagonalMatrix}}}{DMatrix: diagonalMatrix}}

\manheading{Synopsis}
\lstset{language=[95]Fortran}
\begin{lstlisting}
u1 = diagonal      (d1)   ! pure: extract main diagonal into DVector u1
d2 = diagonal      (u1)   ! construct square diagonal DMatrix from DVector u1
d2 = diagonalMatrix(d1)   ! keep only main diagonal of d1, zero off-diagonals
\end{lstlisting}

\manheading{Brief}
\procname{diagonal} extracts a \typename{DMatrix} diagonal as a \typename{DVector} or builds a diagonal \typename{DMatrix} from a \typename{DVector}; \procname{diagonalMatrix} zeroes off-diagonal elements of an existing \typename{DMatrix}.

\manheading{Description}
\begin{itemize}
  \item \lcodef{u1 = diagonal(d1)} (\texttt{pure}): Returns a \typename{DVector} of length \texttt{min(d1 \% m, d1 \% n)} with \texttt{is = d1 \% is} containing the main diagonal entries.
  \item \lcodef{d2 = diagonal(u1)}: Returns a square \typename{DMatrix} of size \texttt{u1 \% n} with \texttt{is = js = u1 \% is}, \texttt{mtype = 'DG'}, \texttt{name = ''}, and diagonal entries from \argname{u1}.
  \item \lcodef{d2 = diagonalMatrix(d1)}: Copies \argname{d1}'s metadata (preserving \texttt{mtype}) and zeroes all off-diagonal elements, keeping the $\min(m,n)$ diagonal elements (non-square input does \emph{not} produce \procname{NaN}).
\end{itemize}
To convert between \typename{DMatrix} and \typename{Matrix} via real/imaginary components, use \procname{real}, \procname{aimag}, and \procname{mcmplx} (documented in Section~\ref{mat:real_aimag_mcmplx}).

\manheading{Dummy Arguments}
\begin{description}[leftmargin=2.2cm, style=nextline]
  \argitem{MAT / MATA}{class(DMatrix), intent(in)}{Input real matrix.}
  \argitem{vec}{class(DVector), intent(in)}{Input real vector for \texttt{diagonal(vec)}.}
  \argitem{result}{type(DVector) or type(DMatrix), return value}{Extracted diagonal vector or diagonal real matrix.}
\end{description}

\manheading{Example}
\lstset{language=[95]Fortran}
\begin{lstlisting}
use :: matrix_mod_matrix
type(DMatrix) :: d1, d2, d3
type(DVector) :: u1

d1 = DMatrix('gaussian', 4, mtype='DS')
u1 = diagonal(d1)
d2 = diagonal(u1)          ! d2 % mtype == 'DG'
d3 = diagonalMatrix(d1)    ! d3 % mtype == 'DS'
\end{lstlisting}

\mansubgroup{Linear Algebra}

% ===================================================================
% inverse (DMatrix)
% ===================================================================
\manentry{Generic Function}{inverse (DMatrix)}{DMatrix: inverse}{Computes the matrix inverse $D^{-1}$ via \texttt{LAPACK} \texttt{dgetrf}/\texttt{dgetri}}{dmat:inverse}

\manheading{Synopsis}
\lstset{language=[95]Fortran}
\begin{lstlisting}
d2 = inverse(d1)
\end{lstlisting}

\manheading{Brief}
Computes and returns the real matrix inverse $D_1^{-1}$ as a new \typename{type(DMatrix)} via \texttt{LAPACK} LU factorization.

\manheading{Description}
\procname{inverse(d1)} copies all metadata from \argname{d1} to \argname{d2}, allocates \texttt{d2 \% v}, and calls \procname{inverse(d1 \% v)} (\texttt{dgetrf}/\texttt{dgetri}). Singular or non-square inputs yield a \typename{DMatrix} filled with \procname{NaN} (detectable with \lcodef{isNaN(d2)}) without aborting.

\manheading{Dummy Arguments}
\begin{description}[leftmargin=2.2cm, style=nextline]
  \argitem{MATA}{class(DMatrix), intent(in)}{Square non-singular real matrix.}
  \argitem{MATB}{type(DMatrix), return value}{Inverse real matrix $D_1^{-1}$ with identical metadata (or \procname{NaN} if singular/non-square).}
\end{description}

\manheading{Example}
\lstset{language=[95]Fortran}
\begin{lstlisting}
use :: matrix_mod_matrix
type(DMatrix) :: d1, d2

d1 = DMatrix('gaussian', 8)
d2 = inverse(d1)
print *, 'Residual:', norm(d1 * d2 - DMatrix('1', 8))
\end{lstlisting}

% ===================================================================
% determinant and lndet (DMatrix)
% ===================================================================
\manentry{Generic Functions}{determinant, lndet (DMatrix)}{DMatrix: determinant, lndet}{Compute $\det D$ or $(\ln|\det D|, \operatorname{sgn}(\det D))$ for a square \texttt{DMatrix}}{dmat:det_lndet}
\addcontentsline{toc}{subsubsection}{\texorpdfstring{\texttt{\textbf{DMatrix: lndet}}}{DMatrix: lndet}}

\manheading{Synopsis}
\lstset{language=[95]Fortran}
\begin{lstlisting}
r  = determinant(d1)     ! real(8) scalar determinant
lr = lndet      (d1)     ! real(8) :: lr(2); det = exp(lr(1)) * lr(2)
\end{lstlisting}

\manheading{Brief}
Evaluate the real determinant $\det D$ (\procname{determinant}) or its overflow-safe logarithmic magnitude and sign (\procname{lndet}) via \texttt{LAPACK} \texttt{dgetrf}.

\manheading{Description}
\begin{itemize}
  \item \procname{determinant(d1)}: Returns a \typename{real(8)} scalar (\texttt{0.0\_8} if singular, \texttt{1.0\_8} if $0\times 0$, \procname{NaN} if non-square). Overflows IEEE double precision for $\sigma = 1$ Gaussian matrices (\lcodef{DMatrix('g', N)}) at about $N \approx 300$; use \procname{lndet} for large matrices.
  \item \procname{lndet(d1)}: Returns a 2-element \typename{real(8)} array \texttt{lr(2)} where \texttt{lr(1) = ln|det D|} and \texttt{lr(2) = +1.0\_8} or \texttt{-1.0\_8} (or \texttt{lr(1) = }$-\infty$, \texttt{lr(2) = 0.0\_8} when singular), so that $\det D = \exp(\texttt{lr(1)}) \times \texttt{lr(2)}$ without overflow.
\end{itemize}

\manheading{Dummy Arguments}
\begin{description}[leftmargin=2.2cm, style=nextline]
  \argitem{MAT}{class(DMatrix), intent(in)}{Square real matrix.}
  \argitem{z}{real(8) or real(8)(2), return value}{Scalar determinant (\procname{determinant}) or \texttt{[ln|det D|, sign(det D)]} (\procname{lndet}).}
\end{description}

\manheading{Example}
\lstset{language=[95]Fortran}
\begin{lstlisting}
use :: matrix_mod_matrix
type(DMatrix) :: d1
real(8)       :: r, lr(2)

d1 = DMatrix('gaussian', 32)
r  = determinant(d1)
lr = lndet(d1)          ! exp(lr(1)) * lr(2) == r
\end{lstlisting}

% ===================================================================
% eigenvalues (DMatrix)
% ===================================================================
\manentry{Generic Function}{eigenvalues (DMatrix)}{DMatrix: eigenvalues}{Computes all eigenvalues of a square \texttt{DMatrix}, returning a complex \texttt{Vector}}{dmat:eigenvalues}

\manheading{Synopsis}
\lstset{language=[95]Fortran}
\begin{lstlisting}
v1 = eigenvalues(d1)   ! v1 is type(Vector)
\end{lstlisting}

\manheading{Brief}
Computes the eigenvalues of a square \typename{type(DMatrix)} and returns them as a complex \typename{type(Vector)}, automatically using \texttt{dsyev} when \texttt{d1 \% mtype = 'DS'}.

\manheading{Description}
Because a general real matrix (\texttt{mtype = 'DG'}) can have complex-conjugate eigenvalue pairs, \procname{eigenvalues(d1)} always returns a complex \typename{type(Vector)} of length \texttt{d1 \% n} with \texttt{is = d1 \% is}. It calls \texttt{eigenvalues(d1 \% v, mtype = d1 \% mtype)}:
\begin{itemize}
  \item When \texttt{d1 \% mtype = 'DS'}, \texttt{LAPACK}'s fast symmetric eigensolver \texttt{dsyev} is invoked, yielding real eigenvalues ordered ascending (with zero imaginary parts; extract via \cbox{v1 \% re()} if a \typename{DVector} is desired).
  \item When \texttt{d1 \% mtype = 'DG'}, \texttt{LAPACK}'s general real eigensolver \texttt{dgeev} is invoked.
\end{itemize}
Aborts via \texttt{matrix\_error} if \argname{d1} is non-square (\texttt{d1 \% m /= d1 \% n}), and returns \procname{NaN} if \texttt{LAPACK} fails to converge.

\manheading{Dummy Arguments}
\begin{description}[leftmargin=2.2cm, style=nextline]
  \argitem{MAT}{class(DMatrix), intent(in)}{Square real matrix (\texttt{mtype='DS'} uses \texttt{dsyev}).}
  \argitem{vec}{type(Vector), return value}{Complex \typename{Vector} containing all $N$ eigenvalues.}
\end{description}

\manheading{Example}
\lstset{language=[95]Fortran}
\begin{lstlisting}
use :: matrix_mod_matrix
type(DMatrix) :: d1
type(Vector)  :: v1
type(DVector) :: u1

d1 = DMatrix('gaussian', 8, mtype='DS')
v1 = eigenvalues(d1)   ! fast dsyev path
u1 = v1 % re()         ! real DVector of eigenvalues
\end{lstlisting}

% ===================================================================
% eigenvectors (DMatrix)
% ===================================================================
\manentry{Generic Function}{eigenvectors (DMatrix)}{DMatrix: eigenvectors}{Computes eigenvalues (\texttt{Vector}) and right eigenvectors (\texttt{Matrix}) of a \texttt{DMatrix}}{dmat:eigenvectors}

\manheading{Synopsis}
\lstset{language=[95]Fortran}
\begin{lstlisting}
m2 = eigenvectors(d1, v2)   ! m2 is type(Matrix), v2 is type(Vector)
m1 = diagonal(v2)           ! d1 * m2 == m2 * m1
\end{lstlisting}

\manheading{Brief}
Computes the complex eigenvalues \argname{v2} (\typename{Vector}) and normalized right eigenvectors \argname{m2} (\typename{Matrix}) of a square \typename{DMatrix} \argname{d1}.

\manheading{Description}
Because general real matrices can possess complex eigenvalues and eigenvectors, \procname{eigenvectors(d1, v2)} returns a \typename{type(Matrix)} \argname{m2} (with \texttt{m2 \% mtype = 'ZG'} and bounds/name copied via \procname{metadata\_copy}) whose $j$-th column is the right eigenvector corresponding to \texttt{v2 \% v(j)}. If \texttt{v2 \% n /= d1 \% n}, \argname{v2} is allocated automatically via \cbox{Vector(d1 \% n, is=d1 \% is)}. The eigen-equation satisfies:
\begin{align}
  D_1 M_2 = M_2 \Lambda\, , \qquad \Lambda = \texttt{diagonal}(v_2)\, .
\end{align}

\begin{notebox}[Ordering, Non-Square Check, and Symmetric \texttt{dsyev} Dispatch]
\begin{enumerate}
  \item For general real matrices (\texttt{'DG'}), always use \argname{m2} together with the vector \argname{v2} returned by the same \procname{eigenvectors} call, since \texttt{eigenvalues(d1)} may return eigenvalues in a different order.
  \item When \texttt{d1 \% mtype = 'DS'}, \procname{eigenvectors} forwards \texttt{mtype = 'DS'} to \modname{matrix\_mod\_array}, dispatching to \texttt{LAPACK} \texttt{dsyev} (reading only the upper triangle) and returning real ascending eigenvalues matching \texttt{eigenvalues(d1)} along with orthonormal eigenvectors.
  \item Non-square inputs (\texttt{d1 \% m /= d1 \% n}) abort via \texttt{matrix\_error}.
\end{enumerate}
\end{notebox}

\manheading{Dummy Arguments}
\begin{description}[leftmargin=2.2cm, style=nextline]
  \argitem{MATA}{class(DMatrix), intent(in)}{Square $N \times N$ real matrix.}
  \argitem{vec}{type(Vector), intent(inout)}{Complex eigenvalue vector (auto-constructed if \texttt{vec \% n /= N}).}
  \argitem{MATB}{type(Matrix), return value}{Complex matrix (\texttt{mtype = 'ZG'}) holding right eigenvector columns.}
\end{description}

\manheading{Example}
\lstset{language=[95]Fortran}
\begin{lstlisting}
use :: matrix_mod_matrix
type(DMatrix) :: d1
type(Matrix)  :: m1, m2
type(Vector)  :: v2

d1 = DMatrix('gaussian', 8)
m2 = eigenvectors(d1, v2)
m1 = diagonal(v2)
print *, 'Residual:', norm(d1 * m2 - m2 * m1)
\end{lstlisting}

\mansubgroup{Math Functions}

% ===================================================================
% abs, sin, cos, exp, log, sqrt (DMatrix)
% ===================================================================
\manentry{Pure Generic Functions}{abs, sin, cos, exp, log, sqrt (DMatrix)}{DMatrix: abs, sin, cos, exp, log, sqrt}{Element-wise real mathematical functions on \texttt{DMatrix}}{dmat:math_funcs}
\addcontentsline{toc}{subsubsection}{\texorpdfstring{\texttt{\textbf{DMatrix: sin}}}{DMatrix: sin}}
\addcontentsline{toc}{subsubsection}{\texorpdfstring{\texttt{\textbf{DMatrix: cos}}}{DMatrix: cos}}
\addcontentsline{toc}{subsubsection}{\texorpdfstring{\texttt{\textbf{DMatrix: exp}}}{DMatrix: exp}}
\addcontentsline{toc}{subsubsection}{\texorpdfstring{\texttt{\textbf{DMatrix: log}}}{DMatrix: log}}
\addcontentsline{toc}{subsubsection}{\texorpdfstring{\texttt{\textbf{DMatrix: sqrt}}}{DMatrix: sqrt}}

\manheading{Synopsis}
\lstset{language=[95]Fortran}
\begin{lstlisting}
d2 = abs(d1);  d2 = sin(d1);  d2 = cos(d1);  d2 = exp(d1)
d2 = log(d1);  d2 = sqrt(d1)   ! log: 0 -> -Inf, <0 -> NaN;  sqrt: <0 -> NaN
\end{lstlisting}

\manheading{Brief}
Evaluate $|x|$, $\sin x$, $\cos x$, $\e{x}$, $\ln x$, and $\sqrt{x}$ element-by-element on \texttt{d1 \% v}, returning a new \typename{type(DMatrix)} with identical metadata.

\manheading{Description}
All six functions are \texttt{pure} and copy all metadata---including \texttt{d1 \% mtype} (\texttt{'DG'} or \texttt{'DS'})---to the returned \typename{DMatrix}. Negative entries give IEEE \texttt{NaN} in \procname{log} and \procname{sqrt} (detectable with \procname{isNaN}). A zero entry gives $-\infty$ in \procname{log}, which \procname{isNaN} does \emph{not} detect; test for such entries with \texttt{any(.not.\ ieee\_is\_finite(d2 \% v))} (module \texttt{ieee\_arithmetic}).

\manheading{Dummy Arguments}
\begin{description}[leftmargin=2.2cm, style=nextline]
  \argitem{MAT}{class(DMatrix), intent(in)}{Input real matrix.}
  \argitem{MATB}{type(DMatrix), return value}{Allocated \typename{DMatrix} preserving \texttt{MAT \% mtype}.}
\end{description}

\manheading{Example}
\lstset{language=[95]Fortran}
\begin{lstlisting}
use :: matrix_mod_matrix
type(DMatrix) :: d1, d2

d1 = DMatrix('uniform', 4, mtype='DS') + 0.5_8
d2 = log(d1) + sqrt(d1)   ! d2 % mtype == 'DS'
\end{lstlisting}

\mansubgroup{Misc Procedures}

% ===================================================================
% metadata_copy (DMatrix)
% ===================================================================
\manentry{Pure Generic Subroutine}{metadata\_copy (DMatrix)}{DMatrix: metadata\textunderscore copy}{Copies metadata between \texttt{DMatrix} and \texttt{Matrix} objects without copying data}{dmat:metadata_copy}

\manheading{Synopsis}
\lstset{language=[95]Fortran}
\begin{lstlisting}
call metadata_copy(src, dst)
\end{lstlisting}

\manheading{Brief}
Copies \texttt{m, n, is, ie, js, je, mtype, name} between any combination of \typename{DMatrix} and \typename{Matrix} objects without allocating \texttt{dst \% v}.

\manheading{Description}
When copying from \typename{Matrix} to \typename{DMatrix}, \texttt{dst \% mtype} is converted from \texttt{'Z*'} to \texttt{'DG'}; when copying from \typename{DMatrix} to \typename{Matrix}, \texttt{'D*'} becomes \texttt{'ZG'}; same-type copies preserve \texttt{mtype} verbatim.

\manheading{Dummy Arguments}
\begin{description}[leftmargin=2.2cm, style=nextline]
  \argitem{MATA}{class(MatrixClass), intent(in)}{Source \typename{DMatrix} or \typename{Matrix}.}
  \argitem{MATB}{class(MatrixClass), intent(inout)}{Destination \typename{DMatrix} or \typename{Matrix}.}
\end{description}

\manheading{Example}
\lstset{language=[95]Fortran}
\begin{lstlisting}
use :: matrix_mod_matrix
type(DMatrix) :: d1, d2
type(Matrix)  :: m1

d1 = DMatrix(4, is=2, mtype='DS', name='Sym')
call metadata_copy(d1, d2)   ! d2 % mtype == 'DS'
call metadata_copy(d1, m1)   ! m1 % mtype == 'ZG'
\end{lstlisting}

% ===================================================================
% isNaN (DMatrix)
% ===================================================================
\manentry{Generic Function}{isNaN (DMatrix)}{DMatrix: isNaN}{Tests whether any entry of a \texttt{DMatrix} is IEEE \texttt{NaN}}{dmat:isNaN}

\manheading{Synopsis}
\lstset{language=[95]Fortran}
\begin{lstlisting}
has_nan = isNaN(d1)
\end{lstlisting}

\manheading{Brief}
Returns \texttt{.true.} if any element in \texttt{d1 \% v} satisfies \texttt{IEEE\_IS\_NAN}, and \texttt{.false.} otherwise.

\manheading{Description}
Calls \procname{isNaN(d1 \% v)} from \modname{matrix\_mod\_array}. Use this function after \procname{inverse}, \procname{log}, or \procname{sqrt} to test for singular matrices or invalid domains (note that \procname{log} of a zero entry gives $-\infty$, not \texttt{NaN}, and an overflowing \procname{determinant} gives $\pm\infty$). See Section~\ref{mod:NaN} for generating \procname{NaN}.

\manheading{Dummy Arguments}
\begin{description}[leftmargin=2.2cm, style=nextline]
  \argitem{MAT}{class(DMatrix), intent(in)}{Real matrix to test.}
  \argitem{itis}{logical, scalar, return value}{\texttt{.true.} if at least one entry is \texttt{NaN}.}
\end{description}

\manheading{Example}
\lstset{language=[95]Fortran}
\begin{lstlisting}
use :: matrix_mod_matrix
type(DMatrix) :: d1, d2

d1 = DMatrix(-1.0_8, 4)
d2 = sqrt(d1)
print *, 'Has NaN?', isNaN(d2)   ! prints T
\end{lstlisting}


% Section 4: Class Vector
% ===================================================================
% Section 4: Class Vector
% ===================================================================
\section{Class \texttt{Vector}}
\label{sec:class_vector}

The \typename{Vector} derived type encapsulates a one-dimensional double-precision complex (\typename{complex(8)} / \typename{complex(real64)}) vector with arbitrary lower/upper index bounds and an optional user label. Unlike \typename{Matrix} and \typename{DMatrix}, \typename{Vector} is a standalone type (it does not extend \texttt{MatrixClass} and has no \texttt{mtype} symmetry field).

% -------------------------------------------------------------------
\subsection{Components}
\label{subsec:vec_components}

% ===================================================================
% type(Vector)
% ===================================================================
\manentry{Derived Type}{type(Vector)}{Vector (type)}{Double-precision complex 1D vector class with index bounds and label metadata}{vec:type}

\manheading{Synopsis}
\lstset{language=[95]Fortran}
\begin{lstlisting}
type(Vector) :: vec
! Components: vec % n, vec % is, vec % ie, vec % name, vec % v(is:ie)
\end{lstlisting}

\manheading{Brief}
Stores a 1D allocatable \typename{complex(8)} array \texttt{v(is:ie)} of length \texttt{n} together with its index range \texttt{is:ie} and a 20-character \texttt{name}.

\manheading{Description}
Before construction, \texttt{vec \% n = vec \% is = vec \% ie = 0} and \texttt{vec \% v} is unallocated. Once constructed, elements are indexed over \texttt{is:ie} where \texttt{ie = is + n - 1}.

\manheading{Data Components}
\begin{description}[leftmargin=2.2cm, style=nextline]
  \argitem{n}{integer, default 0}{Number of vector elements (\texttt{n == 0} indicates uninitialized state).}
  \argitem{is, ie}{integer, default 0}{Lower and upper index bounds of \texttt{v}, satisfying \texttt{ie = is + n - 1}.}
  \argitem{name}{character(20), default \texttt{''}}{User label (\texttt{name\_len = 20}) displayed by \procname{\% print}.}
  \argitem{v(:)}{complex(8), allocatable}{Allocated 1D complex array \texttt{v(is:ie)}.}
\end{description}

\manheading{Example}
\lstset{language=[95]Fortran}
\begin{lstlisting}
use :: matrix_mod_matrix
type(Vector) :: vec

vec = Vector(8, is=0, name='psi')
print *, vec % n, vec % is, vec % ie   ! 8, 0, 7
\end{lstlisting}

% -------------------------------------------------------------------
\subsection{Constructors}
\label{subsec:vec_constructors}

% ===================================================================
% Vector(...) constructor
% ===================================================================
\manentry{Generic Constructor}{Vector(...)}{Vector (constructor)}{Constructs and allocates a \texttt{type(Vector)} (zero, constant, random, or from 1D complex array)}{vec:constructor}

\manheading{Synopsis}
\lstset{language=[95]Fortran}
\begin{lstlisting}
vec = Vector(n,      [is], [name])                ! pure: zero vector
vec = Vector(n, r,   [is], [name])                ! pure: filled with real r
vec = Vector(n, z,   [is], [name])                ! pure: filled with cmplx z
vec = Vector(C,      [is], [name])                ! pure: copied from C(:)
vec = Vector(rtype, n, [is], [name], [sigma])     ! uniform or Gaussian random
\end{lstlisting}

\manheading{Brief}
Allocates a \typename{type(Vector)} of length \argname{n} (or \texttt{size(C)}), configures index bounds \texttt{is:ie} and \texttt{name}, and fills \texttt{vec \% v}.

\manheading{Description}
\begin{enumerate}
  \item \textbf{\texttt{Vector(n, [is], [name])}} (\texttt{pure}): Allocates \texttt{v(is:is+n-1)} and sets all elements to \texttt{(0.0\_8, 0.0\_8)}. Defaults: \texttt{is = 1}, \texttt{name = ''}.
  \item \textbf{\texttt{Vector(n, r, [is], [name])}} and \textbf{\texttt{Vector(n, z, [is], [name])}} (\texttt{pure}): Allocate a vector of length \argname{n} filled with a \typename{real(8)} scalar \argname{r} or \typename{complex(8)} scalar \argname{z}.
  \item \textbf{\texttt{Vector(C, [is], [name])}} (\texttt{pure}): Copies a 1D \typename{complex(8)} array \argname{C(:)}. Pass \texttt{is = lbound(C,1)} to preserve non-unit bounds.
  \item \textbf{\texttt{Vector(rtype, n, [is], [name], [sigma])}}: Dispatches on \argname{rtype(1:1)}: \texttt{'r'}/\texttt{'R'}/\texttt{'u'}/\texttt{'U'} fills with uniform $[0,1)+i[0,1)$ numbers; \texttt{'g'}/\texttt{'G'} fills with Gaussian numbers of standard deviation \argname{sigma} (default \texttt{1.0\_8}). Passing an empty string or any other leading character aborts via \texttt{matrix\_error}.
\end{enumerate}

\begin{warnbox}[Argument Order Differs from \texttt{Matrix}]
In \typename{Vector} constant constructors, the length \argname{n} comes \emph{first} and the scalar comes \emph{second}: \lcodef{Vector(n, r)} and \lcodef{Vector(n, z)} (whereas \lcodef{Matrix(r, m)} takes the scalar first). Also note that \texttt{rtype = '1'} is only supported by \typename{Matrix}/\typename{DMatrix}, not by \typename{Vector}.
\end{warnbox}

\manheading{Dummy Arguments}
\begin{description}[leftmargin=2.2cm, style=nextline]
  \argitem{n}{integer, intent(in)}{Length of the vector.}
  \argitem{r / z}{real(8) / complex(8), intent(in)}{Constant fill value (passed \emph{after} \argname{n}).}
  \argitem{C}{complex(8), dimension(:), intent(in)}{Source 1D complex array for \texttt{Vector(C, ...)}.}
  \argitem{is}{integer, optional, intent(in)}{Lower index bound (defaults to \texttt{1}).}
  \argitem{name}{character(*), optional, intent(in)}{Optional vector label (up to 20 characters; defaults to \texttt{''}).}
  \argitem{rtype}{character(*), intent(in)}{Random generator: \texttt{'uniform'}/\texttt{'random'} (\texttt{'u'}/\texttt{'r'}) or \texttt{'gaussian'} (\texttt{'g'}).}
  \argitem{sigma}{real(8), optional, intent(in)}{Gaussian standard deviation (defaults to \texttt{1.0\_8}).}
  \argitem{vec}{type(Vector), return value}{Allocated \typename{Vector} object.}
\end{description}

\manheading{Example}
\lstset{language=[95]Fortran}
\begin{lstlisting}
use :: matrix_mod_matrix
complex(8)   :: c(0:3) = (1.0_8, -1.0_8)
type(Vector) :: v0, v1, v2, v3, v4

v0 = Vector(4)
v1 = Vector(4, 2.5_8, is=0, name='ConstR')
v2 = Vector(4, (1.0_8, 2.0_8))
v3 = Vector('gaussian', 8, sigma=0.5_8)
v4 = Vector(c, is=lbound(c,1))
\end{lstlisting}

% -------------------------------------------------------------------
\subsection{Component Procedures}
\label{subsec:vec_comp_procs}

\typename{type(Vector)} defines eight public type-bound generics: \procname{save}, \procname{read}, \procname{print}, \procname{random}, \procname{gaussian}, \procname{conjg}, \procname{re}, and \procname{im}. (Note that the internal type-bound procedure \texttt{vector\_norm} is private; compute norms with the standalone function \procname{norm(vec)}.)

% ===================================================================
% VEC % save and VEC % read
% ===================================================================
\manentry{Type-Bound Subroutines}{VEC \% save, VEC \% read}{Vector \% save, Vector \% read}{Plain-text sequential I/O for \texttt{Vector} elements}{vec:save_read}
\addcontentsline{toc}{subsubsection}{\texorpdfstring{\texttt{\textbf{Vector \% read}}}{Vector \% read}}

\manheading{Synopsis}
\lstset{language=[95]Fortran}
\begin{lstlisting}
call vec % save([unit], [fmt])
call vec % read([unit])
\end{lstlisting}

\manheading{Brief}
Write (\procname{\% save}) or read (\procname{\% read}) raw \typename{complex(8)} elements \texttt{vec \% v(is:ie)} sequentially from a Fortran unit.

\manheading{Description}
\procname{vec \% save} writes each element \texttt{vec \% v(i)} from \texttt{i = is} to \texttt{ie} using \texttt{'(2'//fmt//')'} (defaults: \texttt{unit = f\_mout}, \texttt{fmt = 'G28.17'}). \procname{vec \% read} reads \texttt{x, y} pairs via list-directed \texttt{read(un,*)} (default internal \texttt{f\_minput = stdin}) into pre-allocated \texttt{vec \% v(is:ie)}. Both return silently if \texttt{vec \% v} is unallocated.

\manheading{Dummy Arguments}
\begin{description}[leftmargin=2.2cm, style=nextline]
  \argitem{vec}{class(Vector), intent(in) / intent(inout)}{Target vector (\texttt{in} for \procname{\% save}; \texttt{inout} for \procname{\% read}).}
  \argitem{unit}{integer, optional, intent(in)}{Fortran I/O unit (defaults to \procname{f\_mout} for save, \texttt{stdin} for read).}
  \argitem{fmt}{character(*), optional, intent(in)}{Real edit descriptor for \procname{\% save} (defaults to \texttt{'G28.17'}).}
\end{description}

\manheading{Example}
\lstset{language=[95]Fortran}
\begin{lstlisting}
use :: matrix_mod_matrix
type(Vector) :: v1, v2
integer      :: un

v1 = Vector('uniform', 4)
open(newunit=un, file='vec.dat', status='replace')
call v1 % save(unit=un)
rewind(un)
v2 = Vector(4)
call v2 % read(unit=un)
close(un)
\end{lstlisting}

% ===================================================================
% VEC % print
% ===================================================================
\manentry{Type-Bound Subroutine}{VEC \% print}{Vector \% print}{Pretty-prints a \texttt{Vector}, metadata info, or Mathematica list}{vec:print}

\manheading{Synopsis}
\lstset{language=[95]Fortran}
\begin{lstlisting}
call vec % print([unit], [fmt], [form], [name], [ips], [ipe])
\end{lstlisting}

\manheading{Brief}
Displays \texttt{vec \% v} (or sub-range \texttt{ips:ipe}) in \texttt{'Default'}, \texttt{'Info'}, or \texttt{'Mathematica'} format, with optional label override \argname{name}.

\manheading{Description}
Checks that \texttt{vec \% n /= 0} and \texttt{vec \% v} is allocated (aborting via \texttt{matrix\_error} otherwise), clamps \argname{ips, ipe} inside \texttt{is:ie}, and prints:
\begin{itemize}
  \item \texttt{form = 'Default'} (default): \texttt{\#pvstart: <name>} header, one complex element per line with \texttt{(2//fmt)} (default \texttt{fmt = 'G28.17'}), and \texttt{\#pvend: ]} footer.
  \item \texttt{form = 'Info'}: Single-line metadata and allocation bounds summary.
  \item \texttt{form = 'Mathematica'}: 1D list \texttt{<name> = \{...\};} using \texttt{Interpreter["Number"]}.
\end{itemize}

\manheading{Dummy Arguments}
\begin{description}[leftmargin=2.2cm, style=nextline]
  \argitem{vec}{class(Vector), intent(in)}{Initialized complex vector.}
  \argitem{unit}{integer, optional, intent(in)}{Output Fortran unit (defaults to \procname{f\_mout}).}
  \argitem{fmt}{character(*), optional, intent(in)}{Real format descriptor (defaults to \texttt{'G28.17'}).}
  \argitem{form}{character(*), optional, intent(in)}{\texttt{'Default'}, \texttt{'Info'}, or \texttt{'Mathematica'} (checked by first character).}
  \argitem{name}{character(*), optional, intent(in)}{Optional override for \texttt{vec \% name} in output headers.}
  \argitem{ips, ipe}{integer, optional, intent(in)}{Optional starting and ending index slice.}
\end{description}

\manheading{Example}
\lstset{language=[95]Fortran}
\begin{lstlisting}
use :: matrix_mod_matrix
type(Vector) :: v1

v1 = Vector('gaussian', 20, name='eval')
call v1 % print(form='Info')
call v1 % print(ips=1, ipe=5, form='Mathematica')
\end{lstlisting}

% ===================================================================
% VEC % random and VEC % gaussian
% ===================================================================
\manentry{Type-Bound Subroutines}{VEC \% random, VEC \% gaussian}{Vector \% random, Vector \% gaussian}{Fill a pre-allocated \texttt{Vector} with uniform or Gaussian complex random entries}{vec:rand_gauss}
\addcontentsline{toc}{subsubsection}{\texorpdfstring{\texttt{\textbf{Vector \% gaussian}}}{Vector \% gaussian}}

\manheading{Synopsis}
\lstset{language=[95]Fortran}
\begin{lstlisting}
call vec % random()
call vec % gaussian([sigma])
\end{lstlisting}

\manheading{Brief}
Overwrite \texttt{vec \% v} with uniform $[0,1)+i[0,1)$ (\procname{\% random}) or Box--Muller Gaussian (\procname{\% gaussian}) complex random numbers.

\manheading{Description}
Delegate to \cbox{call random\_number(vec \% v)} and \cbox{call random\_number(vec \% v, sigma=s)} (where \texttt{s} defaults to \texttt{1.0\_8}). Return silently if \texttt{vec \% v} is unallocated.

\manheading{Dummy Arguments}
\begin{description}[leftmargin=2.2cm, style=nextline]
  \argitem{vec}{class(Vector), intent(inout)}{Pre-allocated complex vector.}
  \argitem{sigma}{real(8), optional}{Gaussian standard deviation for \procname{\% gaussian} (defaults to \texttt{1.0\_8}).}
\end{description}

\manheading{Example}
\lstset{language=[95]Fortran}
\begin{lstlisting}
use :: matrix_mod_matrix
type(Vector) :: v1

v1 = Vector(8)
call v1 % random()
call v1 % gaussian(sigma=0.25_8)
\end{lstlisting}

% ===================================================================
% VEC % conjg, VEC % re, VEC % im
% ===================================================================
\manentry{Type-Bound Pure Functions}{VEC \% conjg, VEC \% re, VEC \% im}{Vector \% conjg, Vector \% re, Vector \% im}{Return the complex conjugate, real part, or imaginary part of a \texttt{Vector}}{vec:comp_re_im_conjg}
\addcontentsline{toc}{subsubsection}{\texorpdfstring{\texttt{\textbf{Vector \% re}}}{Vector \% re}}
\addcontentsline{toc}{subsubsection}{\texorpdfstring{\texttt{\textbf{Vector \% im}}}{Vector \% im}}

\manheading{Synopsis}
\lstset{language=[95]Fortran}
\begin{lstlisting}
v2 = v1 % conjg()   ! pure: returns type(Vector)  with CONJG(v1 % v)
u1 = v1 % re()      ! pure: returns type(DVector) with REAL(v1 % v, kind=8)
u2 = v1 % im()      ! pure: returns type(DVector) with AIMAG(v1 % v)
\end{lstlisting}

\manheading{Brief}
Construct a new \typename{Vector} (\procname{\% conjg}) or \typename{DVector} (\procname{\% re}, \procname{\% im}) copying all metadata (\texttt{n, is, ie, name}) from \argname{v1}.

\manheading{Description}
All three methods are \texttt{pure} type-bound functions. Note that the standalone \procname{conjg}, \procname{real}, and \procname{aimag} generics are \emph{not} overloaded for \typename{Vector}; use these type-bound methods instead.

\manheading{Dummy Arguments}
\begin{description}[leftmargin=2.2cm, style=nextline]
  \argitem{vec}{class(Vector), intent(in)}{Passed-object complex vector.}
  \argitem{vecb}{type(Vector) or type(DVector), return value}{Transformed vector with copied metadata.}
\end{description}

\manheading{Example}
\lstset{language=[95]Fortran}
\begin{lstlisting}
use :: matrix_mod_matrix
type(Vector)  :: v1, v2
type(DVector) :: u_re, u_im

v1   = Vector('gaussian', 6, name='psi')
v2   = v1 % conjg()
u_re = v1 % re()
u_im = v1 % im()
\end{lstlisting}

% -------------------------------------------------------------------
\subsection{Operators}
\label{subsec:vec_operators}

\mansubgroup{Assignment}

% ===================================================================
% assignment(=) Vector
% ===================================================================
\manentry{Overloaded Operator}{assignment(=) (Vector)}{Vector: assignment(=)}{Assigns a \texttt{Vector}, \texttt{DVector}, 1D complex array, or scalar to a \texttt{Vector}}{vec:op_assign}

\manheading{Synopsis}
\lstset{language=[95]Fortran}
\begin{lstlisting}
v1 = v2    ! type(Vector)        -> intrinsic copy (metadata + values)
v1 = u1    ! type(DVector)       -> copies metadata if uninit; aborts on size mismatch
v1 = a     ! complex(8) :: a(:)  -> keeps bounds/name; aborts on size mismatch
v1 = b     ! real(8)    :: b(:)  -> keeps bounds/name; aborts on size mismatch
v1 = r     ! real(8)    :: r     -> fills pre-allocated v1 % v = r
v1 = z     ! complex(8) :: z     -> fills pre-allocated v1 % v = z
\end{lstlisting}

\manheading{Brief}
Overloads \texttt{=} for assigning a \typename{Vector}, \typename{DVector}, 1D \typename{complex(8)} array, or real/complex scalar to a \typename{Vector}.

\manheading{Description}
\begin{itemize}
  \item \texttt{v1 = v2}: Copies all metadata and data cleanly whether or not \argname{v1} was initialized.
  \item \texttt{v1 = u1}, \texttt{v1 = a}, and \texttt{v1 = b}: If \texttt{v1 \% n == 0} (uninitialized), initializes metadata (from \argname{u1} or \texttt{is = 1..size}, \texttt{name = ''}) and copies values (promoted to \typename{complex(8)} for real sources). If \argname{v1} was already initialized, preserves its \texttt{is, ie, name} and copies values; if \texttt{v1 \% n /= size}, the program stops with an \texttt{error stop} message naming the procedure (exit status 1).
  \item \texttt{v1 = r} and \texttt{v1 = z}: Require \texttt{v1 \% v} to be allocated (aborting via Fortran \texttt{error stop} if unconstructed) and assign \texttt{v1 \% v = r} or \texttt{z}.
\end{itemize}

\begin{notebox}[Construct \texttt{v1} Before Scalar Assignment]
Always construct \argname{v1} before executing \lcodef{v1 = r} or \lcodef{v1 = z}; assigning a scalar to an unallocated \typename{Vector} aborts via Fortran \texttt{error stop}.
\end{notebox}

\manheading{Dummy Arguments}
\begin{description}[leftmargin=2.2cm, style=nextline]
  \argitem{vecb}{type(Vector), intent(inout)}{Destination complex vector.}
  \argitem{rhs}{type(Vector), type(DVector), complex(8)(:), real(8)(:), real(8), or complex(8), intent(in)}{Source value or array.}
\end{description}

\manheading{Example}
\lstset{language=[95]Fortran}
\begin{lstlisting}
use :: matrix_mod_matrix
type(Vector)  :: v1
type(DVector) :: u1
complex(8)    :: a(4) = (1.0_8, 1.0_8)

v1 = Vector(4, is=0, name='Target')
v1 = a                ! preserves is=0, name='Target'
v1 = (2.0_8, -1.0_8)  ! OK because v1 is allocated
u1 = DVector('uniform', 4)
v1 = u1
\end{lstlisting}

\mansubgroup{Addition, Subtraction, Scalar Multiplication, Division, and Power}

% ===================================================================
% operator(+) and operator(-) Vector
% ===================================================================
\manentry{Overloaded Operators}{operator(+), operator(-) (Vector)}{Vector: operator(+), operator(-)}{Element-wise addition, unary negation, and subtraction for \texttt{Vector}}{vec:op_plus_minus}
\addcontentsline{toc}{subsubsection}{\texorpdfstring{\texttt{\textbf{Vector: operator(-)}}}{Vector: operator(-)}}

\manheading{Synopsis}
\lstset{language=[95]Fortran}
\begin{lstlisting}
v3 = v1 + v2;   v3 = v1 - v2;   v3 = -v1        ! Vector +/- Vector, unary -
v3 = v1 + u1;   v3 = u1 + v1;   v3 = v1 - u1    ! Vector and DVector (pure)
v3 = v1 + a;    v3 = a  + v1;   v3 = v1 - a     ! complex(8)(:) array (pure)
v3 = v1 + b;    v3 = b  + v1;   v3 = v1 - b     ! real(8)(:) array (pure)
v3 = v1 + r;    v3 = r  + v1;   v3 = v1 - r     ! real(8) scalar (pure)
v3 = v1 + z;    v3 = z  + v1;   v3 = v1 - z     ! complex(8) scalar (pure)
! Also reversed subtractions: u1 - v1, a - v1, b - v1, r - v1, z - v1
\end{lstlisting}

\manheading{Brief}
Perform \texttt{pure} element-wise vector addition, subtraction, or unary negation, inheriting metadata from the \typename{Vector} operand (operands of different length abort the program).

\manheading{Description}
All unary and binary \procname{+} and \procname{-} overloads for \typename{Vector} are \texttt{pure} functions returning a \typename{type(Vector)} whose metadata (\texttt{n, is, ie, name}) are copied from the \typename{Vector} operand (or \argname{v1} in \texttt{v1 +/- v2}). If the array/vector lengths differ, the program stops with an \texttt{error stop} message naming the procedure (exit status 1). Note that both \typename{complex(8)(:)} arrays \argname{a} and \typename{real(8)(:)} arrays \argname{b} are supported in \procname{+} and \procname{-}.


\manheading{Dummy Arguments}
\begin{description}[leftmargin=2.2cm, style=nextline]
  \argitem{lhs, rhs}{class(Vector), class(DVector), complex(8)(:), real(8)(:), complex(8), or real(8), intent(in)}{Operands.}
  \argitem{result}{type(Vector), return value}{Allocated sum or difference vector.}
\end{description}

\manheading{Example}
\lstset{language=[95]Fortran}
\begin{lstlisting}
use :: matrix_mod_matrix
type(Vector)  :: v1, v2, v3
type(DVector) :: u1

v1 = Vector('gaussian', 4)
u1 = DVector('uniform', 4)
v2 = v1 + u1                 ! correct mixed addition
v3 = -v1 - (0.0_8, 1.0_8)    ! unary negation and complex scalar shift
\end{lstlisting}

% ===================================================================
% Scalar *, /, ** Vector
% ===================================================================
\manentry{Overloaded Operators}{operator(*), operator(/), operator(**) (Vector)}{Vector: scalar *, /, **}{Scalar multiplication, scalar division, and element-wise power for \texttt{Vector}}{vec:op_scale_pow}
\addcontentsline{toc}{subsubsection}{\texorpdfstring{\texttt{\textbf{Vector: operator(/)}}}{Vector: operator(/)}}
\addcontentsline{toc}{subsubsection}{\texorpdfstring{\texttt{\textbf{Vector: operator(**)}}}{Vector: operator(**)}}

\manheading{Synopsis}
\lstset{language=[95]Fortran}
\begin{lstlisting}
v2 = r * v1;   v2 = v1 * r;   v2 = v1 / r      ! real(8) scaling (pure)
v2 = z * v1;   v2 = v1 * z;   v2 = v1 / z      ! complex(8) scaling (pure)
v2 = v1 ** k;  v2 = v1 ** x;  v2 = v1 ** z     ! element-wise power (pure)
\end{lstlisting}

\manheading{Brief}
Scale a \typename{Vector} by a real or complex scalar or raise each element \texttt{v1 \% v(i)} to an integer, real, or complex power.

\manheading{Description}
All overloads are \texttt{pure} functions copying metadata (\texttt{n, is, ie, name}) from \argname{v1} to \argname{v2}. For \procname{**}, \texttt{k} is \typename{integer}, \texttt{x} is \typename{real(8)}, and \texttt{z} is \typename{complex(8)}.

\manheading{Dummy Arguments}
\begin{description}[leftmargin=2.2cm, style=nextline]
  \argitem{veca}{class(Vector), intent(in)}{Input complex vector.}
  \argitem{scalar / n}{real(8), complex(8), or integer, intent(in)}{Scalar multiplier/divisor or element-wise exponent.}
  \argitem{vecb}{type(Vector), return value}{Result \typename{Vector} with copied metadata.}
\end{description}

\manheading{Example}
\lstset{language=[95]Fortran}
\begin{lstlisting}
use :: matrix_mod_matrix
type(Vector) :: v1, v2

v1 = Vector('gaussian', 4, name='psi')
v2 = (0.0_8, 1.0_8) * v1 / norm(v1)
v2 = v1 ** 2
\end{lstlisting}

\mansubgroup{Matrix--Vector Multiplication}

% ===================================================================
% operator(*) Matrix * Vector
% ===================================================================
\manentry{Overloaded Operator}{operator(*) (Matrix--Vector)}{Vector: Matrix * Vector}{Computes matrix--vector products $M v$ and $v^T M$ via \texttt{BLAS}}{vec:op_matvec}

\manheading{Synopsis}
\lstset{language=[95]Fortran}
\begin{lstlisting}
v2 = m1 * v1   ! m1(is:ie, js:je) * v1 -> v2 % v(is:ie), v2 % name = v1 % name
v2 = v1 * m1   ! v1 * m1(is:ie, js:je) -> v2 % v(js:je), v2 % name = v1 % name
\end{lstlisting}

\manheading{Brief}
Multiplies a \typename{Matrix} \argname{m1} and a \typename{Vector} \argname{v1} via \texttt{LMATMUL} (\texttt{zhemv} or \texttt{zgemv}), returning a \typename{Vector}.

\manheading{Description}
\begin{itemize}
  \item \lcodef{v2 = m1 * v1}: Allocates \argname{v2} with length \texttt{m1 \% m}, bounds \texttt{m1\%is:m1\%ie}, and \texttt{v2 \% name = v1 \% name}. Calls \texttt{zhemv} when \texttt{m1 \% mtype = 'ZH'} (since Hermitian \texttt{m1} has second char \texttt{'H'}) and \texttt{zgemv('N', ...)} otherwise.
  \item \lcodef{v2 = v1 * m1}: Allocates \argname{v2} with length \texttt{m1 \% n}, bounds \texttt{m1\%js:m1\%je}, and \texttt{v2 \% name = v1 \% name}. Uses \texttt{LMATMUL(CONJG(m1 \% v), v1 \% v, type='H')} when \texttt{m1 \% mtype = 'ZH'}, and \texttt{zgemv('T', ...)} otherwise.
\end{itemize}
If matrix and vector dimensions are incompatible (\texttt{m1 \% n /= v1 \% n} or \texttt{m1 \% m /= v1 \% n}), the program aborts via \texttt{matrix\_error}.

\manheading{Dummy Arguments}
\begin{description}[leftmargin=2.2cm, style=nextline]
  \argitem{MATA}{class(Matrix), intent(in)}{Matrix factor.}
  \argitem{vecb}{class(Vector), intent(in)}{Vector factor.}
  \argitem{vecc}{type(Vector), return value}{Product vector (dimension mismatches abort).}
\end{description}

\manheading{Example}
\lstset{language=[95]Fortran}
\begin{lstlisting}
use :: matrix_mod_matrix
type(Matrix) :: m1
type(Vector) :: v1, v2, v3

m1 = Matrix('gaussian', 3, 5, is=0, js=10)
v1 = Vector('gaussian', 5, name='colvec')
v2 = m1 * v1      ! length 3, bounds 0:2
v3 = Vector('gaussian', 3) * m1   ! length 5, bounds 10:14
\end{lstlisting}

% ===================================================================
% mvmult and vmmult (Vector)
% ===================================================================
\manentry{Pure Generic Subroutines}{mvmult, vmmult (Vector)}{Vector: mvmult, vmmult}{Zero-allocation \texttt{BLAS} subroutines for \texttt{Matrix}--\texttt{Vector} multiplication}{vec:mvmult_vmmult}
\addcontentsline{toc}{subsubsection}{\texorpdfstring{\texttt{\textbf{Vector: vmmult}}}{Vector: vmmult}}

\manheading{Synopsis}
\lstset{language=[95]Fortran}
\begin{lstlisting}
call mvmult(m1, v2, v3)   ! pure: v3 % v = m1 % v . v2 % v
call vmmult(v2, m1, v3)   ! pure: v3 % v = transpose(m1 % v) . v2 % v
\end{lstlisting}

\manheading{Brief}
Execute matrix--vector multiplication in-place into \argname{v3} (auto-allocating/resizing \argname{v3} if needed) without operator-overloading temporaries.

\manheading{Description}
Automatically allocates or reconstructs \argname{v3} whenever \texttt{v3 \% v} is not allocated or its length mismatches (\texttt{v3 \% n /= m1 \% m} for \procname{mvmult}, \texttt{v3 \% n /= m1 \% n} for \procname{vmmult}); when \argname{v3} is already allocated with matching length, it writes in-place preserving existing bounds and \texttt{name}. Delegates to \modname{matrix\_mod\_array}, passing \texttt{type = 'H'} (\texttt{zhemv}) in \procname{mvmult} when \texttt{m1 \% mtype = 'ZH'}, whereas \procname{vmmult} calls \texttt{vmmult(v2 \% v, m1 \% v, v3 \% v)} (\texttt{zgemv}). Inner-dimension mismatches (\texttt{m1 \% n /= v2 \% n} for \procname{mvmult}, \texttt{m1 \% m /= v2 \% n} for \procname{vmmult}) stop the program with an \texttt{error stop} message.

\manheading{Dummy Arguments}
\begin{description}[leftmargin=2.2cm, style=nextline]
  \argitem{MATA}{class(Matrix), intent(in)}{Input complex matrix.}
  \argitem{vecb}{class(Vector), intent(in)}{Input complex vector.}
  \argitem{vecc}{type(Vector), intent(inout)}{Destination \typename{Vector} (auto-allocated/resized if needed).}
\end{description}

\manheading{Example}
\lstset{language=[95]Fortran}
\begin{lstlisting}
use :: matrix_mod_matrix
type(Matrix) :: m1
type(Vector) :: v1, v2, v3

m1 = Matrix('gaussian', 64)
v1 = Vector('gaussian', 64); v2 = Vector(64); v3 = Vector(64)
call mvmult(m1, v1, v2)   ! v2 = m1 * v1 without temporary allocation
call vmmult(v1, m1, v3)   ! v3 = v1 * m1 without temporary allocation
\end{lstlisting}

% -------------------------------------------------------------------
\subsection{Procedures}
\label{subsec:vec_procs}

\mansubgroup{Random Numbers}

% ===================================================================
% random_number (Vector)
% ===================================================================
\manentry{Generic Subroutine}{random\_number (Vector)}{Vector: random\textunderscore number}{Overloaded uniform and Gaussian random generator for \texttt{Vector}}{vec:random_number}

\manheading{Synopsis}
\lstset{language=[95]Fortran}
\begin{lstlisting}
call random_number(vec)          ! uniform random numbers in [0,1) + i[0,1)
call random_number(vec, sigma)   ! Gaussian random numbers with std dev sigma
\end{lstlisting}

\manheading{Brief}
Populates a pre-allocated \typename{Vector} with uniform or Gaussian complex random variates (passing \argname{sigma} non-optionally to select the Gaussian form).

\manheading{Description}
Calls \texttt{vector\_random\_set(vec)} or \texttt{random\_number(vec \% v, sigma=sigma)}. Returns silently if \texttt{vec \% v} is not allocated.

\manheading{Dummy Arguments}
\begin{description}[leftmargin=2.2cm, style=nextline]
  \argitem{vec}{class(Vector), intent(inout)}{Pre-allocated \typename{Vector}.}
  \argitem{sigma}{real(8)}{Target Gaussian standard deviation (non-optional in the Gaussian overload).}
\end{description}

\manheading{Example}
\lstset{language=[95]Fortran}
\begin{lstlisting}
use :: matrix_mod_matrix
type(Vector) :: v

call matrix_random_init
v = Vector(10)
call random_number(v)
call random_number(v, sigma=1.0_8)
\end{lstlisting}

\mansubgroup{Math Functions}

% ===================================================================
% abs, sin, cos, exp, log, sqrt (Vector)
% ===================================================================
\manentry{Pure Generic Functions}{abs, sin, cos, exp, log, sqrt (Vector)}{Vector: abs, sin, cos, exp, log, sqrt}{Element-wise mathematical functions on \texttt{Vector}}{vec:math_funcs}
\addcontentsline{toc}{subsubsection}{\texorpdfstring{\texttt{\textbf{Vector: sin}}}{Vector: sin}}
\addcontentsline{toc}{subsubsection}{\texorpdfstring{\texttt{\textbf{Vector: cos}}}{Vector: cos}}
\addcontentsline{toc}{subsubsection}{\texorpdfstring{\texttt{\textbf{Vector: exp}}}{Vector: exp}}
\addcontentsline{toc}{subsubsection}{\texorpdfstring{\texttt{\textbf{Vector: log}}}{Vector: log}}
\addcontentsline{toc}{subsubsection}{\texorpdfstring{\texttt{\textbf{Vector: sqrt}}}{Vector: sqrt}}

\manheading{Synopsis}
\lstset{language=[95]Fortran}
\begin{lstlisting}
u1 = abs(v1)                                       ! returns type(DVector)
v2 = sin(v1);  v2 = cos(v1);  v2 = exp(v1)         ! returns type(Vector)
v2 = log(v1);  v2 = sqrt(v1)                       ! returns type(Vector)
\end{lstlisting}

\manheading{Brief}
Evaluate $|v_i|, \sin v_i, \cos v_i, \e{v_i}, \ln v_i, \sqrt{v_i}$ element-wise on \argname{v1}, preserving \texttt{n, is, ie, name}.

\manheading{Description}
All six functions are \texttt{pure}. \procname{abs(v1)} returns a \typename{type(DVector)} with \texttt{u1 \% v = abs(v1 \% v)}; the remaining five functions return a \typename{type(Vector)}.

\manheading{Dummy Arguments}
\begin{description}[leftmargin=2.2cm, style=nextline]
  \argitem{vec}{class(Vector), intent(in)}{Input complex vector.}
  \argitem{vecb}{type(DVector) or type(Vector), return value}{Element-wise result vector.}
\end{description}

\manheading{Example}
\lstset{language=[95]Fortran}
\begin{lstlisting}
use :: matrix_mod_matrix
type(Vector)  :: v1, v2
type(DVector) :: mags

v1   = Vector('gaussian', 8)
mags = abs(v1)
v2   = exp(v1) + sqrt(v1)
\end{lstlisting}

\mansubgroup{Vector and Misc Procedures}

% ===================================================================
% dot_product, sort, norm (Vector)
% ===================================================================
\manentry{Generic Functions}{dot\_product, sort, norm (Vector)}{Vector: dot\textunderscore product, sort, norm}{Inner product, Quicksort ordering, and normalized $L_1$ norm for \texttt{Vector}}{vec:dot_sort_norm}
\addcontentsline{toc}{subsubsection}{\texorpdfstring{\texttt{\textbf{Vector: sort}}}{Vector: sort}}
\addcontentsline{toc}{subsubsection}{\texorpdfstring{\texttt{\textbf{Vector: norm}}}{Vector: norm}}

\manheading{Synopsis}
\lstset{language=[95]Fortran}
\begin{lstlisting}
z  = dot_product(v1, v2)     ! pure: sum( conjg(v1 % v) * v2 % v )
v2 = sort(v1, [by])          ! Quicksort; by='MF','MR','RF','RR','IF','IR'
r  = norm(v1)                ! pure: sum( abs(v1 % v) ) / v1 % n
\end{lstlisting}

\manheading{Brief}
Compute the Hermitian inner product $\langle v_1, v_2 \rangle$, sort complex vector elements, and evaluate the mean element modulus.

\manheading{Description}
\begin{itemize}
  \item \procname{dot\_product(v1, v2)} (\texttt{pure}): Calls Fortran intrinsic \texttt{dot\_product(v1 \% v, v2 \% v)}, computing $\sum_i (v_1)_i^* (v_2)_i$ as a \typename{complex(8)} scalar. Vectors of different length abort the program.
  \item \procname{sort(v1, [by])}: Copies \argname{v1} to \argname{v2} (preserving metadata) and sets \texttt{v2 \% v = sort(v1 \% v, by=by)} via \modname{matrix\_mod\_array}. The two-character code \argname{by} selects the key (\texttt{'M'} modulus, \texttt{'R'} real part, \texttt{'I'} imaginary part) and direction (\texttt{'F'} ascending, \texttt{'R'} descending), defaulting to \texttt{'MF'}.
  \item \procname{norm(v1)} (\texttt{pure}): Returns $\frac{1}{N} \sum_{i=1}^{N} |v_i|$ as a \typename{real(8)} scalar.
\end{itemize}

\manheading{Dummy Arguments}
\begin{description}[leftmargin=2.2cm, style=nextline]
  \argitem{veca, vecb}{class(Vector), intent(in)}{Input complex vector(s).}
  \argitem{by}{character(*), optional, intent(in)}{Sort order code: \texttt{'MF'} (default), \texttt{'MR'}, \texttt{'RF'}, \texttt{'RR'}, \texttt{'IF'}, or \texttt{'IR'}.}
  \argitem{result}{complex(8), type(Vector), or real(8)}{Inner product, sorted vector, or normalized $L_1$ norm.}
\end{description}

\manheading{Example}
\lstset{language=[95]Fortran}
\begin{lstlisting}
use :: matrix_mod_matrix
type(Vector) :: v1, v2
complex(8)   :: ip
real(8)      :: nrm

v1  = Vector('gaussian', 8)
v2  = sort(v1, by='RF')       ! sort ascending by real part
ip  = dot_product(v1, v2)
nrm = norm(v1)
\end{lstlisting}

% ===================================================================
% mcmplx, metadata_copy, isNaN (Vector)
% ===================================================================
\manentry{Generic Procedures}{mcmplx, metadata\_copy, isNaN (Vector)}{Vector: mcmplx, metadata\textunderscore copy, isNaN}{Complex vector construction, metadata copy, and \texttt{NaN} detection}{vec:misc}
\addcontentsline{toc}{subsubsection}{\texorpdfstring{\texttt{\textbf{Vector: metadata\textunderscore copy}}}{Vector: metadata\textunderscore copy}}
\addcontentsline{toc}{subsubsection}{\texorpdfstring{\texttt{\textbf{Vector: isNaN}}}{Vector: isNaN}}

\manheading{Synopsis}
\lstset{language=[95]Fortran}
\begin{lstlisting}
v1 = mcmplx(u1, u2)            ! Vector from two conformable DVectors
call metadata_copy(src, dst)   ! pure: copies n, is, ie, name (Vector/DVector)
has_nan = isNaN(v1)            ! .true. if any entry in v1 % v is NaN
\end{lstlisting}

\manheading{Brief}
Construct a complex \typename{Vector} from two real \typename{DVector} objects (\procname{mcmplx}), copy vector metadata without allocating data (\procname{metadata\_copy}), or test for \texttt{NaN} (\procname{isNaN}).

\manheading{Description}
\begin{itemize}
  \item \procname{mcmplx(u1, u2)}: Checks that \texttt{u1 \% n == u2 \% n} (aborting via \texttt{matrix\_error} otherwise), copies metadata \texttt{(n, is, ie, name)} from \argname{u1}, and sets \texttt{v1 \% v = CMPLX(u1 \% v, u2 \% v, kind=8)}. Note the module procedure name is \procname{mcmplx}, \emph{not} Fortran's intrinsic \texttt{cmplx}.
  \item \procname{metadata\_copy(src, dst)} (\texttt{pure}): Works across any combination of \typename{Vector} and \typename{DVector}, copying \texttt{n, is, ie, name} without allocating or modifying \texttt{dst \% v}.
  \item \procname{isNaN(v1)}: Returns \texttt{.true.} if any entry of \texttt{v1 \% v} has a \texttt{NaN} real or imaginary part. See Section~\ref{mod:NaN} for generating \texttt{NaN}.
\end{itemize}

\manheading{Dummy Arguments}
\begin{description}[leftmargin=2.2cm, style=nextline]
  \argitem{u1, u2}{class(DVector), intent(in)}{Conformable real vectors supplying real and imaginary parts to \procname{mcmplx}.}
  \argitem{src, dst}{class(Vector) or class(DVector)}{Source (\texttt{intent(in)}) and destination (\texttt{intent(inout)}) for \procname{metadata\_copy}.}
  \argitem{vec}{class(Vector), intent(in)}{Vector inspected by \procname{isNaN}.}
\end{description}

\manheading{Example}
\lstset{language=[95]Fortran}
\begin{lstlisting}
use :: matrix_mod_matrix
type(DVector) :: u1, u2
type(Vector)  :: v1, v2

u1 = DVector('gaussian', 6, name='z')
u2 = DVector('gaussian', 6)
v1 = mcmplx(u1, u2)
call metadata_copy(v1, v2)
allocate(v2 % v(v2 % is : v2 % ie))
v2 = NaN()
print *, 'v2 has NaN?', isNaN(v2)
\end{lstlisting}


% Section 5: Class DVector
% ===================================================================
% Section 5: Class DVector
% ===================================================================
\section{Class \texttt{DVector}}
\label{sec:class_dvector}

The \typename{DVector} derived type provides a structured interface for one-dimensional double-precision real (\typename{real(8)} / \typename{real(real64)}) vectors with custom index bounds and an optional label.

% -------------------------------------------------------------------
\subsection{Components}
\label{subsec:dvec_components}

% ===================================================================
% type(DVector)
% ===================================================================
\manentry{Derived Type}{type(DVector)}{DVector (type)}{Double-precision real 1D vector class with index bounds and label metadata}{dvec:type}

\manheading{Synopsis}
\lstset{language=[95]Fortran}
\begin{lstlisting}
type(DVector) :: vec
! Components: vec % n, vec % is, vec % ie, vec % name, vec % v(is:ie)
\end{lstlisting}

\manheading{Brief}
Encapsulates a 1D allocatable \typename{real(8)} array \texttt{v(is:ie)} of length \texttt{n} along with bounds and label metadata.

\manheading{Description}
Before initialization, \texttt{vec \% n = vec \% is = vec \% ie = 0} and \texttt{vec \% v} is unallocated. After construction, valid indices run from \texttt{is} to \texttt{ie = is + n - 1}.

\manheading{Data Components}
\begin{description}[leftmargin=2.2cm, style=nextline]
  \argitem{n}{integer, default 0}{Length of the vector (\texttt{n == 0} means uninitialized).}
  \argitem{is, ie}{integer, default 0}{Lower and upper index bounds satisfying \texttt{ie = is + n - 1}.}
  \argitem{name}{character(20), default \texttt{''}}{User label (\texttt{name\_len = 20}) displayed by \procname{\% print}.}
  \argitem{v(:)}{real(8), allocatable}{Allocated 1D real array \texttt{v(is:ie)}.}
\end{description}

\manheading{Example}
\lstset{language=[95]Fortran}
\begin{lstlisting}
use :: matrix_mod_matrix
type(DVector) :: u

u = DVector(6, is=-2, name='coords')
print *, u % n, u % is, u % ie    ! 6, -2, 3
\end{lstlisting}

% -------------------------------------------------------------------
\subsection{Constructors}
\label{subsec:dvec_constructors}

% ===================================================================
% DVector(...) constructor
% ===================================================================
\manentry{Generic Constructor}{DVector(...)}{DVector (constructor)}{Constructs and allocates a \texttt{type(DVector)} (zero, constant, random, or from 1D real array)}{dvec:constructor}

\manheading{Synopsis}
\lstset{language=[95]Fortran}
\begin{lstlisting}
u = DVector(n,      [is], [name])                ! pure: zero vector
u = DVector(n, r,   [is], [name])                ! pure: filled with real r
u = DVector(n, z,   [is], [name])                ! pure: filled with Re(z)
u = DVector(a,      [is], [name])                ! pure: copied from real a(:)
u = DVector(rtype, n, [is], [name], [sigma])     ! uniform or Gaussian random
\end{lstlisting}

\manheading{Brief}
Allocates a \typename{type(DVector)} of length \argname{n} (or \texttt{size(a)}), configures \texttt{is:ie} and \texttt{name}, and initializes \texttt{u \% v}.

\manheading{Description}
\begin{enumerate}
  \item \textbf{\texttt{DVector(n, [is], [name])}} (\texttt{pure}): Allocates \texttt{v(is:is+n-1)} filled with \texttt{0.0\_8} (defaults \texttt{is = 1}, \texttt{name = ''}).
  \item \textbf{\texttt{DVector(n, r, [is], [name])}} and \textbf{\texttt{DVector(n, z, [is], [name])}} (\texttt{pure}): Allocate a length-\argname{n} vector filled with \argname{r} or \texttt{REAL(z, kind=8)}.
  \item \textbf{\texttt{DVector(a, [is], [name])}} (\texttt{pure}): Copies a 1D \typename{real(8)} array \argname{a(:)} (pass \texttt{is = lbound(a,1)} to keep custom lower bounds).
  \item \textbf{\texttt{DVector(rtype, n, [is], [name], [sigma])}}: \texttt{'r'}/\texttt{'R'}/\texttt{'u'}/\texttt{'U'} fills with uniform $[0,1)$ numbers; \texttt{'g'}/\texttt{'G'} fills with Gaussian variates of standard deviation \argname{sigma} (default \texttt{1.0\_8}). Passing an empty string or any other leading character calls \texttt{matrix\_error}.
\end{enumerate}

\begin{warnbox}[Size-First Argument Order]
As with \typename{Vector}, the constant-value constructors take \argname{n} \emph{before} the scalar (\lcodef{DVector(n, r)} and \lcodef{DVector(n, z)}), opposite to \typename{DMatrix}.
\end{warnbox}

\manheading{Dummy Arguments}
\begin{description}[leftmargin=2.2cm, style=nextline]
  \argitem{n}{integer, intent(in)}{Number of vector elements.}
  \argitem{r / z}{real(8) / complex(8), intent(in)}{Scalar fill value (\texttt{REAL(z, kind=8)} if complex), passed \emph{after} \argname{n}.}
  \argitem{a}{real(8), dimension(:), intent(in)}{Source 1D real array for \texttt{DVector(a, ...)}.}
  \argitem{is}{integer, optional, intent(in)}{Lower index bound (defaults to \texttt{1}).}
  \argitem{name}{character(*), optional, intent(in)}{Optional vector label (up to 20 characters).}
  \argitem{rtype}{character(*), intent(in)}{Generator type: \texttt{'uniform'}/\texttt{'random'} or \texttt{'gaussian'}.}
  \argitem{sigma}{real(8), optional, intent(in)}{Gaussian standard deviation (defaults to \texttt{1.0\_8}).}
  \argitem{vec}{type(DVector), return value}{Allocated \typename{DVector}.}
\end{description}

\manheading{Example}
\lstset{language=[95]Fortran}
\begin{lstlisting}
use :: matrix_mod_matrix
real(8)       :: a(0:3) = [1.0_8, 2.0_8, 3.0_8, 4.0_8]
type(DVector) :: u0, u1, u2, u3

u0 = DVector(4)
u1 = DVector(4, 3.14_8, name='pi')
u2 = DVector('gaussian', 8, sigma=0.5_8)
u3 = DVector(a, is=lbound(a,1))
\end{lstlisting}

% -------------------------------------------------------------------
\subsection{Component Procedures}
\label{subsec:dvec_comp_procs}

\typename{type(DVector)} defines five type-bound generic procedures: \procname{save}, \procname{read}, \procname{print}, \procname{random}, and \procname{gaussian}.

% ===================================================================
% DVEC % save, read, print
% ===================================================================
\manentry{Type-Bound Subroutines}{DVEC \% save, DVEC \% read, DVEC \% print}{DVector \% save, DVector \% read, DVector \% print}{Sequential I/O and formatted display for \texttt{DVector}}{dvec:io}
\addcontentsline{toc}{subsubsection}{\texorpdfstring{\texttt{\textbf{DVector \% read}}}{DVector \% read}}
\addcontentsline{toc}{subsubsection}{\texorpdfstring{\texttt{\textbf{DVector \% print}}}{DVector \% print}}

\manheading{Synopsis}
\lstset{language=[95]Fortran}
\begin{lstlisting}
call vec % save ([unit], [fmt])
call vec % read ([unit])
call vec % print([unit], [fmt], [form], [name], [ips], [ipe])
\end{lstlisting}

\manheading{Brief}
Write raw real values (\procname{\% save}), read raw real values (\procname{\% read}), or print formatted values, metadata summaries, or Mathematica lists (\procname{\% print}) for a \typename{DVector}.

\manheading{Description}
\begin{itemize}
  \item \procname{vec \% save}: Writes one \typename{real(8)} number per line using \texttt{'('//fmt//')'} (defaults: \texttt{unit = f\_mout}, \texttt{fmt = 'G28.17'}). Returns silently if unallocated.
  \item \procname{vec \% read}: Reads real scalars from \argname{unit} (default internal \texttt{f\_minput = stdin}) into pre-allocated \texttt{vec \% v(is:ie)}. Returns silently if unallocated.
  \item \procname{vec \% print}: Supports \texttt{form = 'Default'}, \texttt{'Info'}, and \texttt{'Mathematica'} (emitting \texttt{<name> = \{...\};}), optional slice bounds \argname{ips:ipe}, and optional header \argname{name}. Aborts via \texttt{matrix\_error} if uninitialized (\texttt{n == 0}), unallocated, or given an unrecognized \argname{form}.
\end{itemize}

\manheading{Dummy Arguments}
\begin{description}[leftmargin=2.2cm, style=nextline]
  \argitem{vec}{class(DVector), intent(in) / intent(inout)}{Target real vector.}
  \argitem{unit}{integer, optional, intent(in)}{Fortran I/O unit.}
  \argitem{fmt}{character(*), optional, intent(in)}{Numeric format descriptor (defaults to \texttt{'G28.17'}).}
  \argitem{form}{character(*), optional, intent(in)}{\texttt{'Default'}, \texttt{'Info'}, or \texttt{'Mathematica'} (for \procname{\% print}).}
  \argitem{name}{character(*), optional, intent(in)}{Label override for \procname{\% print}.}
  \argitem{ips, ipe}{integer, optional, intent(in)}{Optional index slice for \procname{\% print}.}
\end{description}

\manheading{Example}
\lstset{language=[95]Fortran}
\begin{lstlisting}
use :: matrix_mod_matrix
type(DVector) :: u1

u1 = DVector('uniform', 12, name='x')
call u1 % print(ips=1, ipe=5, fmt='F8.4')
call u1 % print(form='Mathematica')
\end{lstlisting}

% ===================================================================
% DVEC % random and DVEC % gaussian
% ===================================================================
\manentry{Type-Bound Subroutines}{DVEC \% random, DVEC \% gaussian}{DVector \% random, DVector \% gaussian}{Fill a pre-allocated \texttt{DVector} with uniform or Gaussian real random numbers}{dvec:rand_gauss}
\addcontentsline{toc}{subsubsection}{\texorpdfstring{\texttt{\textbf{DVector \% gaussian}}}{DVector \% gaussian}}

\manheading{Synopsis}
\lstset{language=[95]Fortran}
\begin{lstlisting}
call vec % random()
call vec % gaussian([sigma])
\end{lstlisting}

\manheading{Brief}
Populate an allocated \typename{DVector} with uniform $[0,1)$ or zero-mean Gaussian random variates (default \texttt{sigma = 1.0\_8}).

\manheading{Description}
Call \cbox{call random\_number(vec \% v)} and \cbox{call random\_number(vec \% v, sigma=s)}, returning silently if \texttt{vec \% v} is not allocated.

\manheading{Dummy Arguments}
\begin{description}[leftmargin=2.2cm, style=nextline]
  \argitem{vec}{class(DVector), intent(inout)}{Pre-allocated real vector.}
  \argitem{sigma}{real(8), optional}{Gaussian standard deviation (defaults to \texttt{1.0\_8}).}
\end{description}

\manheading{Example}
\lstset{language=[95]Fortran}
\begin{lstlisting}
use :: matrix_mod_matrix
type(DVector) :: u

u = DVector(8)
call u % random()
call u % gaussian(sigma=2.0_8)
\end{lstlisting}

% -------------------------------------------------------------------
\subsection{Operators}
\label{subsec:dvec_operators}

\mansubgroup{Assignment}

% ===================================================================
% assignment(=) DVector
% ===================================================================
\manentry{Overloaded Operator}{assignment(=) (DVector)}{DVector: assignment(=)}{Assigns a \texttt{DVector}, \texttt{Vector}, 1D array, or scalar to a \texttt{DVector}}{dvec:op_assign}

\manheading{Synopsis}
\lstset{language=[95]Fortran}
\begin{lstlisting}
u1 = u2    ! type(DVector)        -> intrinsic copy (metadata + values)
u1 = v1    ! type(Vector)         -> assigns Re(v1); aborts on size mismatch
u1 = b     ! real(8)    :: b(:)   -> copies array b(:) into u1 % v
u1 = a     ! complex(8) :: a(:)   -> copies Re(a(:)) into u1 % v
u1 = r     ! real(8)    :: r      -> fills pre-allocated u1 % v = r
u1 = z     ! complex(8) :: z      -> fills pre-allocated u1 % v = Re(z)
\end{lstlisting}

\manheading{Brief}
Overloads \texttt{=} for assigning real or complex vectors, 1D arrays, or scalars to a \typename{DVector}.

\manheading{Description}
\begin{itemize}
  \item \texttt{u1 = u2}: Intrinsic derived-type copy of all metadata and values.
  \item \texttt{u1 = v1}, \texttt{u1 = b}, \texttt{u1 = a}: If \texttt{u1 \% n == 0}, copies metadata (from \argname{v1} or defaults \texttt{is = 1..size}) and values (taking the real part of complex sources). If \argname{u1} is already constructed, preserves \texttt{is, ie, name}; on a size mismatch the program stops with an \texttt{error stop} message naming the procedure (exit status 1).
  \item \texttt{u1 = r} and \texttt{u1 = z}: Require \texttt{u1 \% v} to be allocated (aborting with Fortran \texttt{error stop} if unallocated), filling \texttt{u1 \% v = r} or \texttt{REAL(z, kind=8)}.
\end{itemize}

\manheading{Dummy Arguments}
\begin{description}[leftmargin=2.2cm, style=nextline]
  \argitem{vecb}{type(DVector), intent(inout)}{Destination real vector.}
  \argitem{rhs}{type(DVector), type(Vector), real(8)(:), complex(8)(:), real(8), or complex(8), intent(in)}{Source operand.}
\end{description}

\manheading{Example}
\lstset{language=[95]Fortran}
\begin{lstlisting}
use :: matrix_mod_matrix
type(DVector) :: u1
real(8)       :: b(4) = [1.0_8, 2.0_8, 3.0_8, 4.0_8]

u1 = DVector(4, is=0, name='Grid')
u1 = b               ! preserves is=0, name='Grid'
u1 = 0.0_8           ! OK because u1 is constructed
\end{lstlisting}

\mansubgroup{Addition, Subtraction, Scalar Multiplication, Division, and Power}

% ===================================================================
% operator(+), operator(-), *, /, ** (DVector)
% ===================================================================
\manentry{Overloaded Operators}{operator(+), operator(-), operator(*), operator(/), operator(**) (DVector)}{DVector: +, -, scalar *, /, **}{Real arithmetic and element-wise power operators for \texttt{DVector}}{dvec:op_arith}
\addcontentsline{toc}{subsubsection}{\texorpdfstring{\texttt{\textbf{DVector: operator(+)}}}{DVector: operator(+)}}
\addcontentsline{toc}{subsubsection}{\texorpdfstring{\texttt{\textbf{DVector: operator(-)}}}{DVector: operator(-)}}
\addcontentsline{toc}{subsubsection}{\texorpdfstring{\texttt{\textbf{DVector: operator(/), operator(**)}}}{DVector: operator(/), operator(**)}}

\manheading{Synopsis}
\lstset{language=[95]Fortran}
\begin{lstlisting}
u3 = u1 + u2;   u3 = u1 - u2;   u3 = -u1      ! DVector +/- DVector, unary -
u3 = u1 + b;    u3 = b  + u1;   u3 = u1 - b   ! real(8)(:) array (also b - u1)
u3 = u1 + r;    u3 = r  + u1;   u3 = u1 - r   ! real(8) scalar   (also r - u1)
u2 = r  * u1;   u2 = u1 * r;    u2 = u1 / r   ! real(8) scalar * and /
u2 = u1 ** k;   u2 = u1 ** x                  ! integer k or real(8) x power
\end{lstlisting}

\manheading{Brief}
Perform \texttt{pure} element-wise real vector arithmetic and element-wise exponentiation, preserving metadata from the \typename{DVector} operand.

\manheading{Description}
All unary and binary \procname{+}, \procname{-}, scalar \procname{*}, \procname{/}, and \procname{**} operations on \typename{DVector} are \texttt{pure} functions returning a \typename{type(DVector)}. A length mismatch in array/vector \procname{+} or \procname{-} stops the program with an \texttt{error stop} message naming the procedure.

\begin{notebox}[No Complex-Scalar Overloads on \texttt{DVector}]
Despite appearing in older draft notes, \filename{matrix\_mod\_matrix.f90} does \emph{not} overload \procname{+}, \procname{-}, \procname{*}, \procname{/}, or \procname{**} between a \typename{DVector} and a \typename{complex(8)} scalar or exponent. To combine a \typename{DVector} with complex scalars, promote it to \typename{Vector} first (e.g.\ \lcodef{v = u1; v2 = z * v}).
\end{notebox}

\manheading{Dummy Arguments}
\begin{description}[leftmargin=2.2cm, style=nextline]
  \argitem{lhs, rhs}{class(DVector), real(8)(:), real(8), or integer, intent(in)}{Real vector, real array, real scalar, or integer/real exponent.}
  \argitem{result}{type(DVector), return value}{Allocated real result vector.}
\end{description}

\manheading{Example}
\lstset{language=[95]Fortran}
\begin{lstlisting}
use :: matrix_mod_matrix
type(DVector) :: u1, u2, u3

u1 = DVector('uniform', 6) + 1.0_8
u2 = 2.0_8 * u1 / 3.0_8
u3 = -u1 + (u2 ** 0.5_8)
\end{lstlisting}

\mansubgroup{Matrix--Vector Multiplication}

% ===================================================================
% DMatrix * DVector, mvmult, vmmult
% ===================================================================
\manentry{Overloaded Operator / BLAS Subroutines}{operator(*), mvmult, vmmult (DMatrix--DVector)}{DVector: * and mvmult, vmmult}{Matrix--vector multiplication via \texttt{*} or zero-allocation \texttt{BLAS} subroutines}{dvec:matvec}
\addcontentsline{toc}{subsubsection}{\texorpdfstring{\texttt{\textbf{DVector: mvmult}}}{DVector: mvmult}}
\addcontentsline{toc}{subsubsection}{\texorpdfstring{\texttt{\textbf{DVector: vmmult}}}{DVector: vmmult}}

\manheading{Synopsis}
\lstset{language=[95]Fortran}
\begin{lstlisting}
u2 = d1 * u1             ! u2 % v(d1%is:d1%ie), u2 % name = u1 % name
u2 = u1 * d1             ! u2 % v(d1%js:d1%je), u2 % name = u1 % name
call mvmult(d1, u1, u2)  ! pure: zero-alloc u2 % v = d1 % v . u1 % v
call vmmult(u1, d1, u2)  ! pure: zero-alloc u2 % v = transpose(d1%v) . u1%v
\end{lstlisting}

\manheading{Brief}
Multiply a \typename{DMatrix} \argname{d1} and a \typename{DVector} \argname{u1} using \texttt{BLAS} \texttt{dgemv} via either operator notation (\procname{*}) or zero-overhead subroutines (\procname{mvmult}, \procname{vmmult}).

\manheading{Description}
Operator \procname{*} allocates \argname{u2} with bounds taken from \argname{d1} (\texttt{is:ie} for \texttt{d1 * u1}, \texttt{js:je} for \texttt{u1 * d1}), copies \texttt{u1 \% name}, and aborts via \texttt{matrix\_error} on a dimension mismatch. The \texttt{pure} subroutines \procname{mvmult} and \procname{vmmult} automatically construct or resize \argname{u2} whenever \texttt{u2 \% v} is unallocated or its length mismatches (\texttt{u2 \% n /= d1 \% m} for \procname{mvmult}, \texttt{u2 \% n /= d1 \% n} for \procname{vmmult}); when \argname{u2} is already allocated with matching length, they write in-place preserving its existing bounds and \texttt{name}. Note that \procname{mvmult} forwards \texttt{type = 'S'} (\texttt{dsymv}) when \texttt{d1 \% mtype = 'DS'}, whereas \procname{vmmult} calls \modname{matrix\_mod\_array}'s \texttt{vmmult(u1 \% v, d1 \% v, u2 \% v)} without a \texttt{type} argument; inner-dimension mismatches stop the program with an \texttt{error stop} message.

\manheading{Dummy Arguments}
\begin{description}[leftmargin=2.2cm, style=nextline]
  \argitem{MATA}{class(DMatrix), intent(in)}{Real matrix operand.}
  \argitem{vecb}{class(DVector), intent(in)}{Real input vector.}
  \argitem{vecc}{type(DVector), return value / intent(inout)}{Output real vector (auto-allocated/resized if needed for \procname{mvmult}/\procname{vmmult}).}
\end{description}

\manheading{Example}
\lstset{language=[95]Fortran}
\begin{lstlisting}
use :: matrix_mod_matrix
type(DMatrix) :: d1
type(DVector) :: u1, u2

d1 = DMatrix('gaussian', 8)
u1 = DVector('gaussian', 8)
u2 = d1 * u1
call vmmult(u1, d1, u2)   ! overwrite u2 in-place without temporary allocation
\end{lstlisting}

% -------------------------------------------------------------------
\subsection{Procedures}
\label{subsec:dvec_procs}

\mansubgroup{Random Numbers}

% ===================================================================
% random_number (DVector)
% ===================================================================
\manentry{Generic Subroutine}{random\_number (DVector)}{DVector: random\textunderscore number}{Overloaded uniform and Gaussian random generator for \texttt{DVector}}{dvec:random_number}

\manheading{Synopsis}
\lstset{language=[95]Fortran}
\begin{lstlisting}
call random_number(vec)          ! uniform random numbers in [0,1)
call random_number(vec, sigma)   ! Gaussian random numbers with std dev sigma
\end{lstlisting}

\manheading{Brief}
Fills a pre-allocated \typename{DVector} with uniform $[0,1)$ or Gaussian (\argname{sigma} required) real random numbers.

\manheading{Description}
Calls \texttt{dvector\_random\_set} or \texttt{random\_number(vec \% v, sigma=sigma)}, returning silently if \texttt{vec \% v} is unallocated.

\manheading{Dummy Arguments}
\begin{description}[leftmargin=2.2cm, style=nextline]
  \argitem{vec}{class(DVector), intent(inout)}{Pre-allocated real vector.}
  \argitem{sigma}{real(8)}{Target Gaussian standard deviation (non-optional in the Gaussian overload).}
\end{description}

\manheading{Example}
\lstset{language=[95]Fortran}
\begin{lstlisting}
use :: matrix_mod_matrix
type(DVector) :: u

call matrix_random_init
u = DVector(8)
call random_number(u)
call random_number(u, sigma=0.5_8)
\end{lstlisting}

\mansubgroup{Math Functions}

% ===================================================================
% abs, sin, cos, exp, log, sqrt (DVector)
% ===================================================================
\manentry{Pure Generic Functions}{abs, sin, cos, exp, log, sqrt (DVector)}{DVector: abs, sin, cos, exp, log, sqrt}{Element-wise real mathematical functions on \texttt{DVector}}{dvec:math_funcs}
\addcontentsline{toc}{subsubsection}{\texorpdfstring{\texttt{\textbf{DVector: sin}}}{DVector: sin}}
\addcontentsline{toc}{subsubsection}{\texorpdfstring{\texttt{\textbf{DVector: cos}}}{DVector: cos}}
\addcontentsline{toc}{subsubsection}{\texorpdfstring{\texttt{\textbf{DVector: exp}}}{DVector: exp}}
\addcontentsline{toc}{subsubsection}{\texorpdfstring{\texttt{\textbf{DVector: log}}}{DVector: log}}
\addcontentsline{toc}{subsubsection}{\texorpdfstring{\texttt{\textbf{DVector: sqrt}}}{DVector: sqrt}}

\manheading{Synopsis}
\lstset{language=[95]Fortran}
\begin{lstlisting}
u2 = abs(u1);  u2 = sin(u1);  u2 = cos(u1);  u2 = exp(u1)
u2 = log(u1);  u2 = sqrt(u1)   ! log: 0 -> -Inf, <0 -> NaN;  sqrt: <0 -> NaN
\end{lstlisting}

\manheading{Brief}
Evaluate $|u_i|, \sin u_i, \cos u_i, \e{u_i}, \ln u_i, \sqrt{u_i}$ element-wise on \argname{u1}, returning a new \typename{DVector} with copied metadata.

\manheading{Description}
All six functions are \texttt{pure}. Negative elements give IEEE \texttt{NaN} in \procname{log} and \procname{sqrt}; a zero element gives $-\infty$ in \procname{log}, which \procname{isNaN} does \emph{not} detect.

\manheading{Dummy Arguments}
\begin{description}[leftmargin=2.2cm, style=nextline]
  \argitem{vec}{class(DVector), intent(in)}{Input real vector.}
  \argitem{vecb}{type(DVector), return value}{Allocated \typename{DVector} with preserved metadata.}
\end{description}

\manheading{Example}
\lstset{language=[95]Fortran}
\begin{lstlisting}
use :: matrix_mod_matrix
type(DVector) :: u1, u2

u1 = DVector('uniform', 6) + 0.1_8
u2 = log(u1) + sqrt(u1)
\end{lstlisting}

\mansubgroup{Vector and Misc Procedures}

% ===================================================================
% dot_product, sort, maxval, minval, norm (DVector)
% ===================================================================
\manentry{Generic Functions}{dot\_product, sort, maxval, minval, norm (DVector)}{DVector: dot\textunderscore product, sort, maxval, minval, norm}{Inner product, Quicksort, extrema, and normalized $L_1$ norm for \texttt{DVector}}{dvec:vec_ops}
\addcontentsline{toc}{subsubsection}{\texorpdfstring{\texttt{\textbf{DVector: sort}}}{DVector: sort}}
\addcontentsline{toc}{subsubsection}{\texorpdfstring{\texttt{\textbf{DVector: maxval}}}{DVector: maxval}}
\addcontentsline{toc}{subsubsection}{\texorpdfstring{\texttt{\textbf{DVector: minval}}}{DVector: minval}}
\addcontentsline{toc}{subsubsection}{\texorpdfstring{\texttt{\textbf{DVector: norm}}}{DVector: norm}}

\manheading{Synopsis}
\lstset{language=[95]Fortran}
\begin{lstlisting}
r  = dot_product(u1, u2)      ! pure: sum( u1 % v * u2 % v )
u2 = sort(u1, [by])           ! by='RF','RR','AF','AR' (default 'RF')
r  = maxval(u1, [mask])       ! pure: max value (optionally masked)
r  = minval(u1, [mask])       ! pure: min value (optionally masked)
r  = norm(u1)                 ! pure: sum( abs(u1 % v) ) / u1 % n
\end{lstlisting}

\manheading{Brief}
Provide the real Euclidean inner product, real Quicksort, masked extrema (\procname{maxval}/\procname{minval}), and mean absolute norm for \typename{DVector}.

\manheading{Description}
\begin{itemize}
  \item \procname{dot\_product(u1, u2)} (\texttt{pure}): Evaluates $\sum_i (u_1)_i (u_2)_i$ as a \typename{real(8)} scalar. Vectors of different length abort the program.
  \item \procname{sort(u1, [by])}: Sorts \texttt{u1 \% v} by signed value (\texttt{'R'}) or absolute value (\texttt{'A'}), forward ascending (\texttt{'F'}) or reverse descending (\texttt{'R'}), defaulting to \texttt{'RF'}.
  \item \procname{maxval(u1, [mask])} and \procname{minval(u1, [mask])} (\texttt{pure}): Wrap Fortran intrinsics \texttt{maxval(u1 \% v, mask)} and \texttt{minval(u1 \% v, mask)}. Note that \procname{maxval} and \procname{minval} are defined only for \typename{DVector}, not \typename{Vector}.
  \item \procname{norm(u1)} (\texttt{pure}): Returns $\frac{1}{N}\sum_{i=1}^N |u_i|$.
\end{itemize}

\manheading{Dummy Arguments}
\begin{description}[leftmargin=2.2cm, style=nextline]
  \argitem{veca, vecb}{class(DVector), intent(in)}{Input real vector(s).}
  \argitem{by}{character(*), optional, intent(in)}{Sort specification: \texttt{'RF'} (default), \texttt{'RR'}, \texttt{'AF'}, or \texttt{'AR'}.}
  \argitem{mask}{logical, dimension(:), optional, intent(in)}{Boolean mask conformable with \texttt{veca \% v} for \procname{maxval}/\procname{minval}.}
  \argitem{result}{real(8) or type(DVector)}{Scalar statistic or sorted \typename{DVector}.}
\end{description}

\manheading{Example}
\lstset{language=[95]Fortran}
\begin{lstlisting}
use :: matrix_mod_matrix
type(DVector) :: u1, u2
real(8)       :: rmax, rmin_neg

u1       = DVector('gaussian', 10)
u2       = sort(u1, by='AR')             ! descending by |u1(i)|
rmax     = maxval(u1)
rmin_neg = minval(u1, mask=(u1%v < 0.0_8))
\end{lstlisting}

% ===================================================================
% metadata_copy, isNaN, mcmplx ref (DVector)
% ===================================================================
\manentry{Generic Procedures}{metadata\_copy, isNaN (DVector)}{DVector: metadata\textunderscore copy, isNaN}{Metadata copying and \texttt{NaN} testing for \texttt{DVector}}{dvec:misc}
\addcontentsline{toc}{subsubsection}{\texorpdfstring{\texttt{\textbf{DVector: isNaN}}}{DVector: isNaN}}

\manheading{Synopsis}
\lstset{language=[95]Fortran}
\begin{lstlisting}
call metadata_copy(u1, dst)   ! pure: dst may be type(DVector) or type(Vector)
has_nan = isNaN(u1)           ! .true. if any entry of u1 % v is NaN
v1 = mcmplx(u1, u2)           ! combine two DVectors into a Vector (Sec. 4.5)
\end{lstlisting}

\manheading{Brief}
Copy bounds/length/name metadata between \typename{DVector} and \typename{Vector} objects (\procname{metadata\_copy}) or test for IEEE \texttt{NaN} elements (\procname{isNaN}).

\manheading{Description}
\procname{metadata\_copy} copies \texttt{n, is, ie, name} without allocating \texttt{dst \% v}. \procname{isNaN(u1)} returns \texttt{.true.} iff any element of \texttt{u1 \% v} satisfies \texttt{IEEE\_IS\_NAN}. Use \procname{mcmplx(u1, u2)} (Section~\ref{vec:misc}) to assemble two real \typename{DVector}s into a complex \typename{Vector}, and \procname{NaN(r)} (Section~\ref{mod:NaN}) to produce real \texttt{NaN} scalars.

\manheading{Dummy Arguments}
\begin{description}[leftmargin=2.2cm, style=nextline]
  \argitem{veca, vecb}{class(DVector) or class(Vector)}{Source (\texttt{intent(in)}) and destination (\texttt{intent(inout)}) for \procname{metadata\_copy}.}
  \argitem{vec}{class(DVector), intent(in)}{Real vector inspected by \procname{isNaN}.}
\end{description}

\manheading{Example}
\lstset{language=[95]Fortran}
\begin{lstlisting}
use :: matrix_mod_matrix
type(DVector) :: u1, u2
real(8)       :: r = 0.0_8

u1 = DVector(4, is=1, name='realvec')
call metadata_copy(u1, u2)
allocate(u2 % v(u2 % is : u2 % ie))
u2 % v = NaN(r)
print *, 'isNaN(u2)?', isNaN(u2)
\end{lstlisting}


% Sections 6-8: Module array_mod, File Structure, Performance, and References
% ===================================================================
% Sections 6-8: Module array_mod, File Structure, Performance, and References
% ===================================================================

% ===================================================================
% Section 6: Module array_mod
% ===================================================================
\section{Module \texttt{array\_mod}}
\label{sec:module_array}

Under the hood, \modname{matrix\_mod\_matrix} delegates numerical linear algebra kernels---matrix/vector products, LU inversions, determinants, Pfaffians, eigensolvers, sorting, norms, and Gaussian random number generation---to \modname{matrix\_mod\_array}. The standalone public interface to \modname{matrix\_mod\_array} (along with \modname{tensorprod\_mod} and \modname{timer\_mod}) is exposed by the \modname{array\_mod} module (\filename{array\_mod.f90}), which operates directly on standard Fortran allocatable and assumed-shape arrays without derived-type wrapper overhead.

For complete documentation of \modname{array\_mod}, consult its companion manual:
\begin{itemize}
  \item Local manual (in \filename{Documentation/}): \filename{array\_mod\_doc.pdf} (source: \filename{array\_mod\_doc.tex})
  \item Online Google Drive draft: \href{https://docs.google.com/document/d/19KoFFvpxTcm9FN1zGqdGaBopyJhl2vYA0D1mEHbTUZ8/edit?usp=sharing}{array module documentation (Google Docs)}
\end{itemize}
You can freely use \modname{array\_mod} alongside \modname{matrix\_mod\_matrix} (or pass \texttt{mat \% v} / \texttt{vec \% v} directly to \modname{array\_mod} routines such as \procname{lmatmul}, \procname{tensorprod}, or \procname{eigenvectors(..., sortby=...)}).

% ===================================================================
% Section 7: File Structure
% ===================================================================
\section{File Structure}
\label{sec:file_structure}

The source files of the \texttt{Matrix} library in \filename{Software/Tools/Matrix/} follow a modular naming convention. The third component of the filename identifies the module:
\lstset{language=[95]Fortran}
\begin{lstlisting}
matrix_mod_common.f90   ! Common constants (dp, PI, ONE, ZERO, IMU), f_mout
matrix_mod_matrix.f90   ! Public types, generic interfaces, operator bindings
matrix_mod_array.f90    ! Core array linear algebra procedures
array_mod.f90           ! Clean public facade (array, timer, tensorprod)
tensorprod_mod.f90      ! Kronecker tensor products and vectorization helpers
timer_mod.f90           ! CPU and wall-clock execution timers
pfapack/fortran/        ! Bundled PFAPACK library (f95_pfapack, skpfa, skpf10)
\end{lstlisting}

Files with a fourth filename component are included directly into their parent module source file via Fortran \texttt{include} statements:
\lstset{language=[95]Fortran}
\begin{lstlisting}
matrix_mod_matrix_matrixClass.f90  ! Private base MatrixClass methods
matrix_mod_matrix_matrix.f90       ! type(Matrix)  procedures and operators
matrix_mod_matrix_dmatrix.f90      ! type(DMatrix) procedures and operators
matrix_mod_matrix_vector.f90       ! type(Vector)  procedures and operators
matrix_mod_matrix_dvector.f90      ! type(DVector) procedures and operators
matrix_mod_array_lapack.f90        ! BLAS/LAPACK/PFAPACK array wrappers
\end{lstlisting}

Compilation is managed by the provided \filename{Makefile}:
\lstset{language=bash}
\begin{lstlisting}
make gnu   # Compiles PFAPACK and gfortran *.mod/*.o -> libMatrix.a
make intel # Compiles PFAPACK and ifx/MKL *.mod/*.o -> libMatrix.a
make mods  # Compiles the Fortran modules and produces object/.mod files
make lib   # Compiles the modules (if needed) and archives them into libMatrix.a
make FC=ifort FFLAGS='-O3 -heap-arrays -xAVX -parallel' \
              LIBS='-mkl=parallel' lib
\end{lstlisting}

% ===================================================================
% Section 8: Performance
% ===================================================================
\section{Performance}
\label{sec:performance}

A comprehensive suite of matrix--matrix and matrix--vector multiplication benchmarks comparing \modname{matrix\_mod\_matrix} derived types against raw Fortran arrays was carried out on a desktop workstation\footnote{Desktop configuration: Intel(R) Core(TM) i7-4790 CPU @ 3.60\,GHz with Intel \texttt{ifort} Version 19.1.1.217 Build 20200306, compiled with \texttt{-O3 -heap-arrays -xHost -parallel -mkl=parallel}, \texttt{nthreads=8} (most likely up to 4 physical cores were utilized).} and on the ARIS national supercomputer.\footnote{Supercomputer configuration (ARIS): Intel Xeon E5-2680v2, Intel \texttt{ifort} Version 19.0.1.144 Build 20181018, compiled with \texttt{-O3 -heap-arrays -xCORE-AVX-I -parallel -mkl=parallel}, \texttt{nthreads=20}.} The benchmark programs and raw outputs (\filename{matrix\_benchmark*}, cited in the table captions) are not part of the current \filename{Software/Tools/Matrix/} distribution; the tables below reproduce the published results.

\begin{notebox}[Summary of Performance Findings]
\begin{itemize}
  \item \textbf{General Design Trade-Off}: The derived types \typename{Matrix}, \typename{DMatrix}, \typename{Vector}, and \typename{DVector} are intended primarily for expressive scientific programming and educational use rather than raw high-performance computing. Their object abstraction carries a cost, though for sufficiently large matrices ($N \ge 200$) the overhead becomes small in several critical operations.
  \item \textbf{Matrix--Matrix Multiplication (Table~\ref{tab:mat_b1}: B.1 \& B.3)}: Overloaded operator syntax (\lcodef{m3 = m1 * m2}) incurs significant allocation/copy overhead for $N \le 500$. Using the subroutine interface \lcodef{call mmmult(m1, m2, m3)} matches direct raw-array \lcodef{call mmmult(a1, a2, a3)} performance across almost all tested sizes.
  \item \textbf{Matrix--Vector Multiplication (Table~\ref{tab:mat_b4}: B.4 \& B.6)}: On the desktop workstation, \lcodef{call mvmult(m1, v2, v3)} matches raw-array \lcodef{call mvmult(a1, u2, u3)} closely. On the 20-thread ARIS supercomputer node, however, the derived-type \procname{mvmult} pays a $27\%$--$36\%$ penalty for $N \ge 2000$, while operator overloading \lcodef{v3 = m1 * v2} pays a $27\%$--$52\%$ penalty (Table~B.6).
\end{itemize}
\end{notebox}

% -------------------------------------------------------------------
% Table B.2: Legend for Matrix-Matrix Benchmarks
% -------------------------------------------------------------------
\begin{table}[htbp]
\centering
\begin{tabular}{c l}
\toprule
\textbf{Column} & \textbf{Matrix--Matrix Multiplication Method} \\
\midrule
\texttt{(3)} & Operator overloading: \lcodef{m3 = m1 * m2} \\
\texttt{(7)} & Derived-type subroutine: \lcodef{call mmmult(m1, m2, m3)} \\
\texttt{(8)} & Raw-array subroutine: \lcodef{call mmmult(a1, a2, a3)} \\
\bottomrule
\end{tabular}
\caption{\label{tab:mat_b2}Legend of the methods benchmarked in Table~\ref{tab:mat_b1} (sub-tables B.1 and B.3) with \texttt{type(Matrix) :: m1, m2, m3} and \texttt{complex(8), allocatable :: a1(:,:), a2(:,:), a3(:,:)}.}
\end{table}

% -------------------------------------------------------------------
% Table B.1 & B.3: Desktop and Supercomputer Matrix-Matrix Benchmarks
% -------------------------------------------------------------------
\begin{table}[htbp]
\centering
\scriptsize
\setlength{\tabcolsep}{4pt}
\begin{tabular}{r c c c @{\hspace{1.2em}} c c c}
\toprule
 & \multicolumn{3}{c}{\textbf{Table B.1: Desktop (Core i7-4790)}} & \multicolumn{3}{c}{\textbf{Table B.3: ARIS (Xeon E5-2680v2)}} \\
\cmidrule(lr){2-4} \cmidrule(lr){5-7}
$N$ & \texttt{(3) m1*m2} & \texttt{(7) mmmult(m)} & \texttt{(8) mmmult(a)} & \texttt{(3) m1*m2} & \texttt{(7) mmmult(m)} & \texttt{(8) mmmult(a)} \\
\midrule
30   & \texttt{5.38e-06} & \texttt{3.43e-06} & \texttt{3.54e-06} & \texttt{2.12e-05} & \texttt{1.36e-05} & \texttt{1.38e-05} \\
60   & \texttt{2.09e-05} & \texttt{1.37e-05} & \texttt{1.37e-05} & \texttt{4.01e-05} & \texttt{2.28e-05} & \texttt{2.28e-05} \\
120  & \texttt{1.08e-04} & \texttt{8.03e-05} & \texttt{7.98e-05} & \texttt{1.26e-04} & \texttt{6.78e-05} & \texttt{6.76e-05} \\
190  & \texttt{3.69e-04} & \texttt{2.91e-04} & \texttt{2.92e-04} & \texttt{3.71e-04} & \texttt{2.40e-04} & \texttt{2.37e-04} \\
250  & \texttt{7.89e-04} & \texttt{6.52e-04} & \texttt{6.51e-04} & \texttt{6.83e-04} & \texttt{4.47e-04} & \texttt{4.56e-04} \\
500  & \texttt{6.31e-03} & \texttt{5.33e-03} & \texttt{5.34e-03} & \texttt{3.93e-03} & \texttt{2.72e-03} & \texttt{2.74e-03} \\
750  & \texttt{2.00e-02} & \texttt{1.72e-02} & \texttt{1.72e-02} & \texttt{1.31e-02} & \texttt{9.75e-03} & \texttt{9.69e-03} \\
1000 & \texttt{4.49e-02} & \texttt{4.06e-02} & \texttt{4.06e-02} & \texttt{2.83e-02} & \texttt{2.24e-02} & \texttt{2.22e-02} \\
2000 & \texttt{3.40e-01} & \texttt{3.43e-01} & \texttt{3.43e-01} & \texttt{2.00e-01} & \texttt{1.71e-01} & \texttt{1.70e-01} \\
4000 & \texttt{2.56e+00} & \texttt{2.44e+00} & \texttt{2.47e+00} & \texttt{1.47e+00} & \texttt{1.28e+00} & \texttt{1.28e+00} \\
\bottomrule
\end{tabular}
\caption{\label{tab:mat_b1}\label{tab:mat_b3}Time per $N\times N$ complex matrix--matrix multiplication in seconds on the desktop workstation (\texttt{matrix\_benchmark\_v01.dat}, Table~B.1) and on the ARIS supercomputer (\texttt{matrix\_benchmark\_v00\_aris.dat}, Table~B.3). Columns correspond to lines starting with \texttt{mm}; definitions are given in Table~\ref{tab:mat_b2}.}
\end{table}

% -------------------------------------------------------------------
% Table B.5: Legend for Matrix-Vector Benchmarks
% -------------------------------------------------------------------
\begin{table}[htbp]
\centering
\begin{tabular}{c l}
\toprule
\textbf{Column} & \textbf{Matrix--Vector Multiplication Method} \\
\midrule
\texttt{(3)} & Operator overloading: \lcodef{v3 = m1 * v2} \\
\texttt{(6)} & Derived-type subroutine: \lcodef{call mvmult(m1, v2, v3)} \\
\texttt{(7)} & Raw-array subroutine: \lcodef{call mvmult(a1, u2, u3)} \\
\bottomrule
\end{tabular}
\caption{\label{tab:mat_b5}Legend of the methods benchmarked in Table~\ref{tab:mat_b4} (sub-tables B.4 and B.6) with \texttt{type(Matrix) :: m1}, \texttt{type(Vector) :: v2, v3}, and \texttt{complex(8), allocatable :: a1(:,:), u2(:), u3(:)}.}
\end{table}

% -------------------------------------------------------------------
% Table B.4 & B.6: Desktop and Supercomputer Matrix-Vector Benchmarks
% -------------------------------------------------------------------
\begin{table}[htbp]
\centering
\scriptsize
\setlength{\tabcolsep}{4pt}
\begin{tabular}{r c c c @{\hspace{1.2em}} c c c}
\toprule
 & \multicolumn{3}{c}{\textbf{Table B.4: Desktop (Core i7-4790)}} & \multicolumn{3}{c}{\textbf{Table B.6: ARIS (Xeon E5-2680v2)}} \\
\cmidrule(lr){2-4} \cmidrule(lr){5-7}
$N$ & \texttt{(3) m1*v2} & \texttt{(6) mvmult(m)} & \texttt{(7) mvmult(a)} & \texttt{(3) m1*v2} & \texttt{(6) mvmult(m)} & \texttt{(7) mvmult(a)} \\
\midrule
30    & \texttt{9.88e-07} & \texttt{5.63e-07} & \texttt{5.63e-07} & \texttt{1.75e-06} & \texttt{9.15e-07} & \texttt{9.12e-07} \\
60    & \texttt{2.16e-06} & \texttt{1.71e-06} & \texttt{1.71e-06} & \texttt{3.99e-06} & \texttt{3.17e-06} & \texttt{3.17e-06} \\
120   & \texttt{4.81e-06} & \texttt{4.14e-06} & \texttt{4.44e-06} & \texttt{9.59e-06} & \texttt{6.73e-06} & \texttt{1.01e-05} \\
190   & \texttt{1.06e-05} & \texttt{9.82e-06} & \texttt{9.89e-06} & \texttt{1.43e-05} & \texttt{1.05e-05} & \texttt{1.07e-05} \\
250   & \texttt{1.65e-05} & \texttt{1.55e-05} & \texttt{1.58e-05} & \texttt{3.03e-05} & \texttt{2.58e-05} & \texttt{2.84e-05} \\
500   & \texttt{4.11e-05} & \texttt{4.06e-05} & \texttt{4.09e-05} & \texttt{6.54e-05} & \texttt{5.84e-05} & \texttt{5.98e-05} \\
750   & \texttt{1.73e-04} & \texttt{1.70e-04} & \texttt{1.61e-04} & \texttt{1.44e-04} & \texttt{1.14e-04} & \texttt{1.20e-04} \\
1000  & \texttt{6.04e-04} & \texttt{6.38e-04} & \texttt{6.91e-04} & \texttt{1.91e-04} & \texttt{1.58e-04} & \texttt{1.48e-04} \\
2000  & \texttt{2.92e-03} & \texttt{2.87e-03} & \texttt{2.87e-03} & \texttt{2.25e-03} & \texttt{1.97e-03} & \texttt{1.48e-03} \\
4000  & \texttt{1.17e-02} & \texttt{1.15e-02} & \texttt{1.15e-02} & \texttt{7.56e-03} & \texttt{7.45e-03} & \texttt{5.49e-03} \\
8000  & \texttt{4.64e-02} & \texttt{4.62e-02} & \texttt{4.62e-02} & \texttt{2.79e-02} & \texttt{2.77e-02} & \texttt{2.15e-02} \\
16000 & \texttt{1.89e-01} & \texttt{1.87e-01} & \texttt{1.90e-01} & \texttt{1.01e-01} & \texttt{1.01e-01} & \texttt{7.97e-02} \\
\bottomrule
\end{tabular}
\caption{\label{tab:mat_b4}\label{tab:mat_b6}Time per complex matrix--vector multiplication in seconds on the desktop workstation (\texttt{matrix\_benchmark\_v01.dat}, Table~B.4) and on the ARIS supercomputer (\texttt{matrix\_benchmark\_v01\_aris.dat}, Table~B.6). Columns correspond to lines starting with \texttt{mv}; definitions are given in Table~\ref{tab:mat_b5}.}
\end{table}

% ===================================================================
% References
% ===================================================================
\begin{thebibliography}{9}

\bibitem{Wimmer_2012}
M.~Wimmer,
\emph{Algorithm 923: Efficient Numerical Computation of the Pfaffian for Dense and Banded Skew-Symmetric Matrices},
ACM Transactions on Mathematical Software (TOMS) \textbf{38} (2012) 30,
\href{https://arxiv.org/abs/1102.3440}{arXiv:1102.3440 [cond-mat.mes-hall]}.

\bibitem{pfapack}
M.~Wimmer,
\emph{PFAPACK: A library for computing the Pfaffian of skew-symmetric matrices},
\url{https://www.michaelwimmer.org/downloads.html}.

\end{thebibliography}


\end{document}
